% ===========================================================================
%  Chapter 2 -- Literature Review
%  Automated Classification of Lumbar Degenerative Disease from MRI
%
%  Compiles standalone:   pdflatex chapter2 ; biber chapter2 ; pdflatex chapter2
%
%  To fold into a larger thesis instead, delete everything from \documentclass
%  down to \begin{document}, and the \printbibliography/\end{document} at the
%  foot, then % ===========================================================================
%  Chapter 2 -- Literature Review
%  Automated Classification of Lumbar Degenerative Disease from MRI
%
%  Compiles standalone:   pdflatex chapter2 ; biber chapter2 ; pdflatex chapter2
%
%  To fold into a larger thesis instead, delete everything from \documentclass
%  down to \begin{document}, and the \printbibliography/\end{document} at the
%  foot, then % ===========================================================================
%  Chapter 2 -- Literature Review
%  Automated Classification of Lumbar Degenerative Disease from MRI
%
%  Compiles standalone:   pdflatex chapter2 ; biber chapter2 ; pdflatex chapter2
%
%  To fold into a larger thesis instead, delete everything from \documentclass
%  down to \begin{document}, and the \printbibliography/\end{document} at the
%  foot, then % ===========================================================================
%  Chapter 2 -- Literature Review
%  Automated Classification of Lumbar Degenerative Disease from MRI
%
%  Compiles standalone:   pdflatex chapter2 ; biber chapter2 ; pdflatex chapter2
%
%  To fold into a larger thesis instead, delete everything from \documentclass
%  down to \begin{document}, and the \printbibliography/\end{document} at the
%  foot, then \input{chapter2} from your main file.
% ===========================================================================

\documentclass[12pt,a4paper]{report}

\usepackage{lmodern}          % scalable fonts; required for microtype expansion
\usepackage[T1]{fontenc}
\usepackage[utf8]{inputenc}
\usepackage[english]{babel}
\usepackage{microtype}
\usepackage[margin=2.5cm]{geometry}
\usepackage{setspace}
\onehalfspacing

\usepackage{amsmath,amssymb}
\usepackage{booktabs}
\usepackage{longtable}
\usepackage{array}
\usepackage{graphicx}
\usepackage[hidelinks]{hyperref}

\usepackage[
  backend=biber,
  style=numeric-comp,
  sorting=none,
  maxbibnames=6,
  minbibnames=3,
  giveninits=true,
  doi=true,
  url=false,
  isbn=false
]{biblatex}

\addbibresource{../lumbar_spine_mri_ai_literature_inventory.bib}
\addbibresource{references-extra.bib}

% The inventory .bib carries internal bookkeeping fields (provenance note, the
% per-paper summary, the local PDF path). These are working metadata, not
% bibliographic data, so they are suppressed from the printed reference list.
% Leaving them in also breaks compilation, since they contain raw underscores.
\AtEveryBibitem{%
  \clearfield{note}%
  \clearfield{comment}%
  \clearfield{annotation}%
  \clearfield{abstract}%
  \clearfield{file}%
}

\begin{document}

\chapter{Literature Review}
\label{ch:litreview}

% ===========================================================================
\section{Introduction}
\label{sec:lr-intro}

\subsection{Scope and Objectives of the Review}
\label{sec:lr-scope}

Lumbar degenerative disease is among the most common indications for magnetic
resonance imaging (MRI) worldwide, and the interpretation of those images is a
repetitive, subjective and time-consuming task that has attracted sustained
attention from the medical image analysis community \cite{lee2024artificial,zhang2024a}.
This chapter reviews that body of work. Its purpose is not to catalogue every
published model, but to establish what is genuinely settled, what remains
contested, and where a controlled experiment can still add knowledge.

The review is organised around a single guiding question: \emph{what determines
whether an automated system can grade lumbar degenerative severity reliably
enough to be clinically useful?} Four candidate answers recur in the
literature---the quality and definition of the reference labels, the degree to
which the anatomy is localised before it is classified, the choice of
representational backbone, and the configuration of imaging sequences supplied
to the model. Each is examined in turn, and the chapter closes by identifying
the two of these that remain inadequately resolved and that this thesis
therefore addresses.

Three specific objectives follow from this. First, to characterise the clinical
grading systems that supply the ground truth for every supervised model in this
field, and to quantify the reliability ceiling those systems impose. Second, to
trace the methodological evolution from handcrafted features to contemporary
convolutional and attention-based architectures, with particular attention to
comparative claims between the two. Third, to assess the evidence on how much
diagnostic information is carried by each MRI sequence and plane, and thus
whether the computational expense of multi-sequence processing is warranted.

\subsection{Search Strategy, Sources, and Inclusion Criteria}
\label{sec:lr-search}

The literature underpinning this chapter was assembled into a structured
inventory of 92 records spanning 2001 to 2026. Sources were identified through
PubMed, Scopus, IEEE~Xplore and arXiv, using combinations of the terms
\emph{lumbar spine}, \emph{MRI}, \emph{spinal stenosis}, \emph{disc
degeneration}, \emph{deep learning}, \emph{convolutional neural network},
\emph{transformer}, \emph{segmentation} and \emph{grading}. Reference lists of
the retrieved systematic reviews \cite{compte2023are,wang2024machine,liawrungrueang2025artificial,mendes2026deep}
were searched by hand for records the database queries had missed, and the
citation graphs of the most influential frameworks were traced forward.

Records were retained if they satisfied at least one of four criteria: they
define or validate a clinical grading standard for lumbar degenerative disease;
they present an automated method for detecting, segmenting or grading lumbar
pathology on MRI; they externally validate such a method; or they describe a
dataset or benchmark on which such methods are trained. Studies confined to
computed tomography or plain radiography were excluded unless they establish a
methodological principle transferable to MRI, as were studies of the cervical
or thoracic spine except where explicitly used for contrast.

Two caveats on the evidence base should be recorded at the outset. First, the
field's publication rate accelerated sharply after 2024, and a substantial
fraction of the most directly relevant work
\cite{batra2025mscan,acharya2026disccentric,nguyen2026crossspine,islamsk2026an}
exists as preprints that have not completed peer review; these are cited as
indicative of methodological direction rather than as settled results. Second,
performance figures reported across studies are frequently not comparable,
because they derive from different datasets, different label definitions and
different evaluation protocols. Wherever numerical results are quoted in this
chapter, the cohort and validation regime that produced them are stated
alongside, and cross-study comparisons are avoided unless the underlying
protocols genuinely align.

\subsection{Structure of the Chapter}
\label{sec:lr-structure}

Sections~\ref{sec:lr-clinical} to~\ref{sec:lr-grading} establish the clinical
substrate: the disease, the imaging protocol used to observe it, and the
grading systems that convert observation into ordinal labels. This material is
not background decoration. The severity classes a model is trained to predict
are defined entirely by these systems, and their documented unreliability
propagates directly into every reported accuracy figure in the chapters that
follow.

Section~\ref{sec:lr-evolution} traces the methodological history of the field,
and Section~\ref{sec:lr-localization} isolates anatomical localisation as the
design principle that most consistently separates effective systems from
ineffective ones. Sections~\ref{sec:lr-backbones} and~\ref{sec:lr-context} then
address the two axes on which this thesis experiments: the representational
backbone, and the spatial and sequential context supplied to it.
Section~\ref{sec:lr-imbalance} treats class imbalance and the ordinal structure
of the label space, and Section~\ref{sec:lr-datasets} the datasets and
benchmarks---principally the RSNA LumbarDISC collection---on which the field now
converges. Sections~\ref{sec:lr-validation} and~\ref{sec:lr-interpretability}
consider validation practice and interpretability respectively.
Section~\ref{sec:lr-gap} synthesises the whole and states the research gap.

% ===========================================================================
\section{Clinical Background: Lumbar Degenerative Disease}
\label{sec:lr-clinical}

\subsection{Epidemiology and Socioeconomic Burden}
\label{sec:lr-epidemiology}

Low back pain is the single largest contributor to years lived with disability
worldwide, and lumbar degenerative change is among its most frequently
implicated structural correlates \cite{compte2023are,zheng2022deep}. Lifetime
prevalence is reported as high as 84\%, and in the United States alone the
combined direct medical and indirect productivity costs are estimated between
\$50 and \$100 billion annually \cite{richards2026the}. Lumbar spinal stenosis
in particular is the most common indication for spinal surgery in patients over
65 \cite{schizas2010qualitative,bharadwaj2023deep}, and demographic ageing is
expected to increase both its incidence and the imaging workload it generates
\cite{lee2024artificial}.

The operational consequence for this thesis is one of volume. MRI referrals for
low back pain have grown faster than the radiology workforce able to report
them \cite{ghosh2014supervised,compte2023are}, and each lumbar study requires
the reader to assign severity at multiple anatomical levels across several
sequences---an average of 14 to 19 minutes of expert time per study
\cite{compte2023are}. It is this combination of high volume, repetitive
structure and documented inter-reader variability that makes the task a
plausible target for automation, rather than any claim that machines perceive
the underlying pathology more acutely than radiologists do.

\subsection{The Degenerative Cascade}
\label{sec:lr-cascade}

Degeneration of the lumbar spine is a progressive, multifactorial process
driven by ageing, mechanical loading, genetic predisposition and lifestyle
factors \cite{esposito2025mribased,zhao2025advances}. It begins in the
intervertebral disc, whose nucleus pulposus loses proteoglycan content and
therefore water from approximately the third decade of life onward. The
resulting dehydration reduces the disc's height, its capacity to distribute
axial load, and its signal intensity on T2-weighted imaging---the last of these
being precisely what the Pfirrmann scale measures \cite{pfirrmann2001magnetic}.

Disc height loss initiates a compensatory cascade through the rest of the motion
segment. The annulus fibrosus weakens and may permit bulging or frank herniation
of nuclear material; the ligamentum flavum buckles and hypertrophies; the
posterior facet joints, now bearing altered load, develop osteoarthritic change
and osteophytosis \cite{bharadwaj2023deep,nikpasand2024automated}. Changes in
the adjacent vertebral endplates and marrow---the Modic
phenotypes~\cite{modic1988degenerative}---frequently accompany this process and
have themselves been the target of automated detection
\cite{gao2022automatic,athertya2021classification}. The clinical significance of
the cascade is that these individually modest changes act cumulatively to
encroach on the spaces occupied by neural tissue.

\subsection{The Three Anatomical Forms of Stenosis}
\label{sec:lr-stenosis-types}

Encroachment is conventionally partitioned by anatomical compartment, and this
partition is directly mirrored in the label structure of the dataset used in
this thesis \cite{richards2026the}:

\begin{enumerate}
  \item \textbf{Spinal canal stenosis} is narrowing of the central canal, which
  contains the thecal sac and the cauda equina. It is best appreciated on axial
  T2-weighted images, where the relationship between cerebrospinal fluid and
  nerve rootlets is visible \cite{schizas2010qualitative,lee2011a}. Severe
  central stenosis characteristically produces neurogenic claudication.

  \item \textbf{Neural foraminal narrowing} is constriction of the lateral
  foramina through which the exiting nerve roots pass. It is assessed on
  parasagittal images, where the epidural fat surrounding the root provides the
  contrast that makes compression visible---T1 weighting being particularly
  suited to this \cite{lee2010a}. It produces dermatomal radicular symptoms.

  \item \textbf{Subarticular stenosis}, also termed lateral recess stenosis,
  affects the zone between the medial edge of the facet and the foramen, where
  it compresses the traversing root before it exits at the level below
  \cite{lurie2008reliability,andreisek2014consensus}.
\end{enumerate}

The distinction matters computationally as well as clinically, because the three
compartments are not equally visible in any single plane. Foraminal narrowing is
poorly assessed from axial images and subarticular stenosis poorly from
sagittal ones, a constraint that Section~\ref{sec:lr-context} shows to have
concrete consequences for model design. It is also the compartment that
correlates with which sequence a model must be given, and thus sits at the
centre of the multi-sequence question this thesis examines.

\subsection{Clinical Presentation and the Treatment Decision}
\label{sec:lr-treatment}

Severity grading is not an academic exercise; it partitions patients into
distinct management pathways. Conservative management---physiotherapy, analgesia,
epidural steroid injection---is appropriate for mild and most moderate disease,
whereas severe compression with corresponding symptoms may warrant decompressive
surgery or fusion \cite{schizas2010qualitative,mysliwiec2010msu}. The evidence
that morphological grade carries genuine prognostic weight is strongest for the
Schizas classification, where grade~C and~D patients were substantially more
likely to fail conservative treatment \cite{schizas2010qualitative}, and for the
MSU herniation grades, where size~2 and~3 lesions identify microdiscectomy
candidates \cite{mysliwiec2010msu}.

Two qualifications temper this. First, the correlation between imaging findings
and symptoms is imperfect: anatomical stenosis is demonstrable in a substantial
proportion of asymptomatic individuals \cite{lurie2008reliability,esposito2025mribased},
so imaging severity constrains rather than determines management. Second, and
more consequentially for automated systems, the clinically decisive boundary is
not evenly distributed across the ordinal scale. The distinction between
\emph{normal} and \emph{mild} rarely changes management, whereas the distinction
between \emph{moderate} and \emph{severe} frequently does. A model that
distributes its errors uniformly across grade boundaries is therefore not
clinically equivalent to one of the same overall accuracy that concentrates its
errors at the harmless end---a point developed in
Section~\ref{sec:lr-error-asymmetry}.

% ===========================================================================
\section{MRI Acquisition: Sequences, Planes, and Diagnostic Yield}
\label{sec:lr-mri}

\subsection{Why MRI is the Reference Modality}
\label{sec:lr-why-mri}

MRI is the reference standard for non-invasive assessment of the lumbar spine.
Computed tomography and plain radiography depict osseous anatomy well but
resolve soft tissue poorly, and both expose the patient to ionising radiation.
MRI resolves the intervertebral disc, thecal sac, individual nerve roots and
ligamentous structures with a contrast unavailable to the alternatives
\cite{zhao2025advances,liawrungrueang2025artificial}, and it is the disposition
of exactly these soft-tissue structures that defines every grading system
reviewed in Section~\ref{sec:lr-grading}. MRI has correspondingly displaced CT
myelography for stenosis assessment \cite{tumko2024a}.

\subsection{Complementary Information Across Sequences and Planes}
\label{sec:lr-sequences}

Lumbar protocols are inherently multiplanar and multisequence, because the spine
is a three-dimensional structure whose pathologies are not all visible in one
projection. Three acquisitions dominate clinical practice and constitute the
input available in the dataset used here \cite{richards2026the}.

\emph{Sagittal T2-weighted} imaging renders fluid hyperintense, so cerebrospinal
fluid in the thecal sac appears bright against darker ligament and bone. It is
the sequence on which disc hydration---and therefore Pfirrmann
grade---is judged \cite{pfirrmann2001magnetic}, and on which the
craniocaudal extent of canal compromise is surveyed. Short tau inversion
recovery (STIR) suppresses the signal from fat and thereby exposes marrow
oedema and inflammatory change, making it the sequence of choice for active
Modic type~1 change \cite{gao2022automatic}.

\emph{Sagittal T1-weighted} imaging renders fat hyperintense. Because the exiting
nerve root within the neural foramen is normally surrounded by epidural fat, T1
provides the contrast against which obliteration of that fat---the basis of the
Lee foraminal grading system---becomes visible \cite{lee2010a}.

\emph{Axial T2-weighted} imaging provides the transverse cross-section through
the disc level. It is the only view in which the morphology of the dural sac,
the arrangement of the cauda equina rootlets and the lateral recesses can be
directly assessed, and it therefore underpins both the Schizas and Lee central
canal systems \cite{schizas2010qualitative,lee2011a} and the assessment of
subarticular stenosis \cite{andreisek2014consensus}.

\subsection{Matching Pathology to Sequence}
\label{sec:lr-sequence-matching}

The correspondence between pathology and sequence is not a matter of
convention but of physics, and the empirical literature bears this out with
unusual clarity. The most direct evidence comes from external validation of
sagittal-only systems. Nigru et al.\ evaluated SpineNetV2 across eleven
radiological features on 1{,}747 lumbosacral discs and found agreement above
$\kappa = 0.7$ for most targets but distinctly lower agreement for foraminal
stenosis and disc herniation, attributing the shortfall specifically to the
absence of transverse anatomical context in sagittal input
\cite{nigru2024external}. Hallinan et al.\ approached the same constraint from
the opposite direction, deliberately assigning axial T2 to central canal and
lateral recess grading while reserving sagittal T1 for the neural foramina
\cite{hallinan2021deep}.

This has a direct bearing on the experiments in this thesis. If a
single-sequence model is to be evaluated, the choice of which sequence to
retain is not arbitrary: sagittal T2 is the most defensible candidate because it
is acquired in essentially every back-pain protocol and because it carries the
signal for canal stenosis and disc degeneration, but the literature predicts
that it should be weakest precisely on the foraminal and subarticular targets.
Whether a sufficiently expressive architecture can compensate for that
predicted weakness is an open question, and one this thesis is designed to
answer.

\subsection{Sources of Acquisition Variability}
\label{sec:lr-variability}

Models trained on one institution's images frequently degrade on another's, and
the proximate causes lie in acquisition. Scanner manufacturer and field strength
(typically 1.5 versus 3.0~T), slice thickness and spacing, in-plane resolution,
echo and repetition times, the presence or absence of fat suppression, and
patient positioning all vary between sites \cite{senzgamboa2023automatic,mendes2026deep}.
Axial acquisitions vary further in whether slices form a continuous stack or are
angled to lie parallel to each disc level---a difference the LumbarDISC
inclusion criteria explicitly accommodate \cite{richards2026the}, and one that
any model consuming axial data must tolerate.

S\'aenz~Gamboa et al.\ treat this variability as the research problem in its own
right rather than as a nuisance, evaluating segmentation across multiple centres
and acquisition parameters and finding that ensembles of U-Net variants were
required to maintain performance \cite{senzgamboa2023automatic}. Batra et al.\
make the complementary design argument: rather than assuming uniform histograms
and a fixed slice count per study, their pipeline is built to accept a variable
number of axial and sagittal slices, on the grounds that this reflects the
conditions under which clinical data actually arrives \cite{batra2025mscan}.
Variability is thus not merely a threat to validity but a specification
constraint, and Section~\ref{sec:lr-generalization} returns to how effectively
the field has met it.

% ===========================================================================
\section{Clinical Grading Systems as the Origin of Ground Truth}
\label{sec:lr-grading}

Every supervised model reviewed in this chapter learns to reproduce labels
generated by a human applying one of the systems described below. The
properties of those systems---their definitions, their granularity and above all
their reliability---therefore bound what any model trained on them can achieve.
This section treats them as the measurement instruments they are.

\subsection{Disc Degeneration: Pfirrmann and its Modified Extension}
\label{sec:lr-pfirrmann}

The Pfirrmann classification, introduced in 2001, remains the most widely
adopted scheme for intervertebral disc degeneration and is applied to sagittal
T2-weighted images \cite{pfirrmann2001magnetic}. It assigns one of five ordinal
grades on the basis of four morphological criteria: homogeneity of the nucleus
pulposus, distinction between nucleus and annulus, T2 signal intensity as a
proxy for hydration, and disc height. Grade~I denotes a homogeneous, bright,
normally proportioned disc; grade~V a collapsed disc of uniformly low signal.

Its virtues are simplicity and universality; its limitations are well
documented. It is a subjective visual judgement, it is insensitive to change
over short intervals, and it captures only intrinsic disc state, saying nothing
about Modic change, high-intensity zones or the downstream compression of neural
structures \cite{esposito2025mribased,zhao2025advances}. It also suffers a
ceiling effect in elderly cohorts, where most discs present as grade~IV or~V and
clinically meaningful variation is compressed into the top of the scale.
Griffith et al.\ addressed this with an eight-level modified scale supported by a
24-image reference panel, validated on 260 discs in 52 subjects of mean age 73;
intra-observer agreement was excellent (weighted $\kappa$ 0.79--0.91) and
inter-observer agreement substantial (0.65--0.67), with most disagreement
confined to a single grade \cite{griffith2007modified}.

Esposito et al.\ place these two systems in a wider context that is worth
recording: a scoping review of 569 studies identified \emph{93} distinct grading
systems for lumbar disc degeneration---63 subjective, 25 quantitative---of which
the Pfirrmann method accounted for just over half of all reported use
\cite{esposito2025mribased}. The apparent standardisation of the field is
therefore partly illusory, and comparisons of automated grading accuracy across
studies must confirm that the same instrument was being reproduced.

\subsection{Central Canal Stenosis: Schizas and Lee}
\label{sec:lr-central-grading}

Two morphological systems have largely displaced simple dural sac
cross-sectional area (DSCA) thresholds for central canal assessment, on the
grounds that area measurements correlate poorly with symptoms and overlap
substantially between symptomatic and asymptomatic individuals
\cite{schizas2010qualitative}.

The \textbf{Schizas classification} grades the morphology of the dural sac on
axial T2 images by the ratio of cerebrospinal fluid to nerve rootlets. Grades~A
and~B retain visible CSF---in~A the rootlets occupy part of the sac, in~B they
fill it but remain individually discernible, giving a characteristic grainy
appearance. Grades~C and~D show no discernible rootlets and no CSF signal,
distinguished from one another by whether posterior epidural fat is preserved
(C) or obliterated (D) \cite{schizas2010qualitative}. Applied to 95 subjects
across surgical, conservatively managed and low-back-pain groups, intra- and
inter-observer agreement were substantial and moderate respectively
($\kappa = 0.65$ and $0.44$). Critically, area measurement over-diagnosed
stenosis in 35 patients and under-diagnosed it in 12 relative to the
morphological grade, and grades~C and~D predicted failure of conservative
management \cite{schizas2010qualitative}.

The \textbf{Lee central canal system} offers a simpler four-grade alternative
based on the degree of separation of the cauda equina rootlets on axial T2:
grade~0 no stenosis; grade~1 obliteration of the anterior CSF space with all
rootlets still separated; grade~2 some rootlets aggregated; grade~3 none
separated \cite{lee2011a}. Assessed by four musculoskeletal radiologists across
81 patients, it achieved inter-reader ICC of 0.730--0.953 and intra-reader
$\kappa$ of 0.863--0.900, and grades tracked cross-sectional area and
anteroposterior diameter at nearly all levels---notably better agreement than
the Schizas system, at the cost of finer morphological detail.

\subsection{Neural Foraminal Stenosis: Lee and Sartoretti}
\label{sec:lr-foraminal-grading}

The \textbf{Lee foraminal system} grades parasagittal images by the pattern of
perineural fat obliteration and the morphology of the nerve root: grade~0
normal; grade~1 fat obliteration in two opposing directions; grade~2
obliteration in all four directions without morphological change to the root;
grade~3 nerve root collapse or distortion \cite{lee2010a}. Evaluated on 576
foramina in 96 patients from L3--L4 to L5--S1 by two radiologists, inter- and
intra-observer agreement were near-perfect ($\kappa$ approximately 0.80--1.0).
This is the direct clinical precursor of the left and right foraminal narrowing
targets in the benchmark used in this thesis. The Sartoretti classification
offers a six-point sagittal alternative (grades~A--F) indexed to progressive
contact between the exiting root and each foraminal boundary.

The near-perfect agreement Lee et al.\ report deserves scepticism when
generalised. It was obtained by two experienced readers applying explicit
criteria to a curated series; the multi-reader, multi-institution agreement
figures in Section~\ref{sec:lr-reliability} are considerably lower, and it is
the latter that describe the conditions under which large training corpora are
actually annotated.

\subsection{Subarticular and Lateral Recess Stenosis}
\label{sec:lr-subarticular-grading}

Subarticular stenosis is graded on axial T2 by the proportion of the lateral
recess compromised: mild at up to one third, moderate between one and two
thirds, severe beyond two thirds, with nerve root involvement described as
touching, displacing or compressing \cite{lurie2008reliability,andreisek2014consensus}.
Compared with the disc and central canal systems this is a notably less
crisply specified instrument, resting on a visually estimated ratio rather than
a discrete morphological pattern. That imprecision shows in the reliability
data: subarticular stenosis is consistently the least reproducible of the
targets \cite{lurie2008reliability}, and it is also the target on which the
severity distribution in LumbarDISC is least skewed
(Section~\ref{sec:lr-lumbardisc}). Its difficulty is thus intrinsic, not merely
a consequence of scarce severe examples.

\subsection{Disc Herniation: Nomenclature and the MSU Classification}
\label{sec:lr-herniation-grading}

Herniation is described qualitatively by the standardised nomenclature of
bulge, protrusion, extrusion and sequestration \cite{fardon2014lumbar}, and
quantitatively by the Michigan State University (MSU) classification, which
grades size against a transverse line joining the medial edges of the facet
joints on the axial slice of maximal herniation \cite{mysliwiec2010msu}.
Grade~1 extends less than half the distance to that line, grade~2 more than
half without crossing it, grade~3 completely beyond it. The MSU system supplies
the reference standard for several automated herniation studies
\cite{zhang2023deep}, and its size~1 versus size~2/3 boundary maps onto the
conservative versus surgical decision.

\subsection{Inter- and Intra-Observer Reliability}
\label{sec:lr-reliability}

The most consequential study for present purposes is that of Lurie et al., who
assessed standardised MRI interpretation of lumbar stenosis using 58 randomly
selected studies from the SPORT trial read by four independent readers, with 20
re-read after at least a month \cite{lurie2008reliability}. Inter-reader
agreement was substantial for central stenosis ($\kappa = 0.73$) but only
moderate for foraminal stenosis ($0.58$) and nerve root impingement ($0.51$),
and poorest for subarticular stenosis ($0.49$), with marked variation between
reader pairs. Quantitative measurement fared no better in absolute terms: mean
absolute inter-reader difference in thecal sac area was 128~mm$^2$, or 13\%.
Intra-reader agreement exceeded inter-reader agreement for every feature.

Two structural findings recur across the wider reliability literature. First,
agreement is systematically worse for the lateral compartments than for the
central canal---the same ordering that later appears in automated model
performance. Second, disagreement is overwhelmingly confined to adjacent
grades: Griffith et al.\ report 84\% of intra-observer and 91\% of
inter-observer disagreement spanning a single grade
\cite{griffith2007modified}, and the same pattern recurs in automated
evaluation \cite{mcsweeney2023external,udomluck2026multitask}. Earlier consensus
efforts confirm that the underlying definitions were themselves contested: the
Delphi survey of Mamisch et al.\ \cite{mamisch2012radiologic} and the consensus
conference reported by Andreisek et al.\ \cite{andreisek2014consensus} were both
convened precisely because radiologic criteria for stenosis varied between
institutions.

\subsection{Label Noise as a Performance Ceiling}
\label{sec:lr-label-ceiling}

The reliability figures above are not a peripheral caveat; they are a hard
constraint on the interpretation of every result in this field. If four expert
readers agree on subarticular severity at $\kappa = 0.49$, then a model trained
on one reader's labels and evaluated against another's cannot exceed that
agreement, and a reported accuracy approaching unity on such a target should
prompt scrutiny of the evaluation protocol rather than admiration. The
LumbarDISC dataset description states the point explicitly, noting that
published inter-rater agreement on these grades ranges from poor to substantial
depending on anatomical location \cite{richards2026the}.

The field has responded to this in two ways. The first is to make the reference
standard more robust by construction---majority voting across a panel with
adjudication of disagreements, as in Tumko et al.\ (seven radiologists)
\cite{tumko2024a} and Su et al.\ (two clinicians with expert adjudication)
\cite{su2022automatic}. The second is to reframe the evaluation question: rather
than asking whether a model matches a single reader, asking whether its
disagreement with readers falls inside the band of disagreement between readers.
Won et al.\ furnish the cleanest example, reporting that the difference between
their two experts was statistically significant while the difference between
each expert and the model trained on that expert's labels was not
\cite{won2020spinal}. Comparatively few studies adopt either safeguard, and
Section~\ref{sec:lr-reporting} returns to this as a systemic weakness.

\subsection{The Ordinal Structure of the Label Space}
\label{sec:lr-ordinality}

Severity labels are ordinal, not categorical: \emph{normal}, \emph{mild},
\emph{moderate} and \emph{severe} stand in a fixed order, and the distance
between them is neither uniform nor clinically symmetric. Treating them as
unordered classes under categorical cross-entropy---as the majority of the
reviewed work does---discards that structure, penalising a
severe-graded-as-normal error exactly as heavily as a moderate-graded-as-severe
one.

Two empirical observations follow. First, errors concentrate overwhelmingly at
adjacent grade boundaries. McSweeney et al.\ report that SpineNet disagreed with
human raters on 20.83\% of discs but that only 0.85\% of discs differed by more
than one grade \cite{mcsweeney2023external}; Udomluck et al.\ observe the same
concentration \cite{udomluck2026multitask}; Ishimoto et al.\ find exact-grade
agreement of 65.7\% rising to 94.1\% when collapsed to a severe-versus-rest
dichotomy \cite{ishimoto2020could}. Much of the apparent error in ordinal
grading is therefore boundary error rather than gross misclassification.

Second, and less expectedly, explicitly ordinal objectives have not reliably
improved on this. Niemeyer et al.\ tested soft-kappa loss, ordinal
cross-entropy and regression losses against plain cross-entropy for Pfirrmann
grading on 1{,}599 patients and 7{,}948 discs, together with class pooling and
class-weighted losses, and found that none improved overall classification
performance; the default cross-entropy classifier reached $\kappa = 0.92$ with
predictions within one grade of ground truth in 99.2\% of cases
\cite{niemeyer2021a}. Patil et al.\ likewise found performance strongest at the
extremes of the Pfirrmann scale and on the binary low-versus-high distinction,
with intermediate grades remaining the difficulty \cite{patil2026mribased}. The
ordinal structure of the problem is thus real and measurable, but the question
of how to exploit it is genuinely open---a point that informs the loss functions
adopted in Chapter~3.

% ===========================================================================
\section{Evolution of Automated Lumbar MRI Analysis}
\label{sec:lr-evolution}

\subsection{Handcrafted Features and Classical Machine Learning}
\label{sec:lr-classical}

Before deep learning, automated analysis of lumbar MRI depended on features
designed by hand. Investigators specified mathematical descriptors of local
appearance and passed them to a conventional classifier. Ghosh et al.\ combined
Histogram of Oriented Gradients (HOG) descriptors with a support vector machine
and a set of geometric heuristics, reporting 99\% disc localisation accuracy
across 53 clinical cases and 318 discs, and---unusually for the period---emitting
a tight bounding box for each disc rather than a single interior point, so that
a downstream system could operate on the region directly
\cite{ghosh2012a}. Athertya and Kumar took a comparable approach to vertebral
degeneration, extracting local binary patterns, Hu moments and gray-level
co-occurrence matrix (GLCM) statistics to classify Modic changes, endplate
defects and focal changes, reaching 85.91\% accuracy for the extent of Modic
change on 100 patients \cite{athertya2021classification}. Related efforts applied
hybrid classical models to disc abnormality detection
\cite{unal2015detection,castromateos2016intervertebral} and semi-automatic
computer-aided classification of degenerative disease \cite{ruizespaa2015semiautomatic},
while Huber et al.\ pursued texture analysis with decision trees as a
quantitative alternative to visual stenosis grading \cite{huber2019qualitative}.
He et al.\ applied a synchronised superpixel representation to define consistent
regions of interest across levels before classifying neural foraminal stenosis,
an early attack on one of the exact conditions graded in this thesis
\cite{he2016automated}.

These systems established feasibility but generalised poorly. Handcrafted
descriptors are sensitive to scanner hardware, field strength, slice thickness
and patient positioning, and the anatomical variability of the lumbar spine
defeats fixed geometric assumptions. One finding from this era nevertheless
retains value: Athertya and Kumar's feature-sensitivity analysis showed that
GLCM entropy alone sufficed to detect focal changes, while endplate-defect
classification depended almost entirely on the first two Hu moments
\cite{athertya2021classification}. Such explicit attribution of predictive power
to individual features is something the deep models that displaced these methods
recovered only later, and imperfectly, through the interpretability techniques of
Section~\ref{sec:lr-interpretability}.

\subsection{The First End-to-End Deep Learning Systems}
\label{sec:lr-first-dl}

Convolutional neural networks removed the need for explicit feature design by
learning hierarchical representations directly from pixels \cite{he2016deep}. In
lumbar imaging the decisive contribution was SpineNet
\cite{jamaludin2017spinenet}. Trained on 2{,}009 patients and 12{,}018 disc
volumes from sagittal T2 images, it predicted six radiological scores
simultaneously---Pfirrmann grade, disc narrowing, upper and lower endplate
defects, and upper and lower marrow changes---from a shared convolutional trunk
that branched into task-specific heads. Two design decisions proved durable.
First, multi-task learning let one architecture serve several targets, reducing
the free parameters relative to training separate networks and exploiting the
correlation between co-occurring pathologies. Second, the system required only
disc-level class labels: no voxel-wise segmentation and no slice-level
localisation, with the evidence hotspots emerging implicitly from classification
training. Reported performance was state-of-the-art and near-human across all
six scores.

The same group demonstrated shortly afterwards that longitudinal clinical scans
could serve as free supervision. A Siamese network was pre-trained with a
contrastive loss on whether two scans belonged to the same patient---information
available from DICOM metadata without knowing the patient's identity---together
with an auxiliary vertebral-level prediction task, then fine-tuned for
degeneration grading. On 1{,}016 subjects, 423 with follow-up imaging, the
pre-trained network outperformed training from scratch and reached equivalent
performance from far fewer annotated examples \cite{jamaludin2017selfsupervised}.
This anticipated by several years the self-supervised pretraining strategies now
routine in medical imaging, and it is the direct ancestor of the contrastive
approach discussed in Section~\ref{sec:lr-contrastive}.

SpineNet's principal limitation was structural rather than architectural: it
consumed preprocessed sagittal disc volumes, and could therefore not assess
pathologies that require the transverse plane. SpineNetV2 extended coverage to
eleven radiological features and moved to a multi-stage design, using a U-Net to
detect and label vertebral bodies and discs before a ResNet classifier graded
each \cite{windsor2022spinenetv2}. The reliance on sagittal input nonetheless
persisted, and Section~\ref{sec:lr-generalization} shows this to be the
consistent locus of its external-validation failures.

\subsection{From Binary Detection to Multi-Level, Multi-Pathology Grading}
\label{sec:lr-task-evolution}

The task definition itself has evolved, and the trajectory matters more than any
individual model. Early deep learning work asked binary questions: is stenosis
present or absent \cite{altun2023lssvgg16}; is a herniation visible in this image
\cite{tsai2021lumbar}. Such formulations are tractable but clinically thin,
because management depends on severity and on which level is affected.

Three shifts followed. The first was to \emph{multi-class severity}: Won et al.\
graded stenosis status on 542 L4--5 axial images against two independent experts
\cite{won2020spinal}; Niemeyer et al.\ predicted the full five-point Pfirrmann
scale across 7{,}948 discs \cite{niemeyer2021a}. The second was to
\emph{multi-pathology} output, in which one model reports several findings at
once: Lehnen et al.\ evaluated disc herniation, bulging, canal stenosis,
nerve-root compression and spondylolisthesis across 888 lumbar segments from a
single CNN \cite{lehnen2021detection}; Su et al.\ graded herniation, central
canal stenosis and nerve-root compression from a shared ResNet-50 trunk with
three classification heads \cite{su2022automatic}. The third was to
\emph{per-level, per-side} prediction, in which the unit of analysis is not the
study or the image but the individual disc level and, for the lateral
compartments, the individual side. Hallinan et al.\ exemplify the mature form,
grading central canal, lateral recess and neural foraminal stenosis at each
level with the sequence appropriate to each \cite{hallinan2021deep}. Single-stage
detectors have since been adapted to deliver the same output more cheaply: Wang
et al.\ modified a YOLOv5 network to detect regions of interest and grade
central, lateral recess and foraminal stenosis simultaneously
\cite{wang2025a}, and a related design grades Pfirrmann degeneration, herniation
and high-intensity zones from sagittal and axial input in one pass
\cite{wang2025automated}. Minin et al.\ pursued the same objective for Pfirrmann
grading with an explicitly open model, motivated by the inaccessibility of
proprietary graders \cite{minin2025spinescan}.

This third formulation is precisely that of the benchmark used in this thesis,
which requires twenty-five graded outputs per study: five conditions across five
disc levels \cite{rsna2024rsna,richards2026the}. The progression is therefore not
merely historical. It explains why performance figures from the binary era
cannot be compared with contemporary results, and why a model that classifies a
whole image is solving a different and substantially easier problem than one
that must attribute severity to a specific anatomical location.

\subsection{Chronological Synthesis}
\label{sec:lr-chronology}

Four phases are discernible. Until roughly 2016 the field was dominated by
handcrafted features and classical classifiers, with localisation as the central
problem \cite{ghosh2012a,ghosh2014supervised}. From 2017 to 2020, CNNs displaced
these methods and multi-task grading became feasible
\cite{jamaludin2017spinenet,lu2018deepspine,han2018spinegan}. From 2021 to 2023
the emphasis moved to clinical credibility: external validation
\cite{mcsweeney2023external}, interpretability \cite{bharadwaj2023deep,gao2022automatic},
reader-assistance evaluation \cite{lim2022improved} and multi-centre testing
\cite{yi2023deep,su2022automatic}. From 2024 onward, the release of a large
public benchmark \cite{richards2026the} redirected effort toward anatomically
localised, multi-view severity grading, alongside the arrival of transformer and
vision-language architectures
\cite{batra2025mscan,nguyen2026crossspine,islamsk2026an,zeng2025gpt4lfs}.

The through-line is a steady increase in the specificity of what is being
predicted, accompanied---until recently---by a decrease in the rigour with which
it is validated. Both systematic reviews of the period make the same complaint:
Compte et al.\ found that algorithms performed markedly worse in replication or
external validation than in their development studies \cite{compte2023are}, and
Wang et al.\ concluded that despite pooled sensitivity of 0.84 and specificity of
0.87 across twelve studies and 15{,}044 patients, no model was reliable enough
for clinical practice \cite{wang2024machine}. Section~\ref{sec:lr-validation}
takes up this tension directly.

% ===========================================================================
\section{Anatomical Localization as a Design Principle}
\label{sec:lr-localization}

\subsection{Disc and Vertebra Detection, Labelling, and Level Assignment}
\label{sec:lr-detection}

Before severity can be attributed to L4--L5, the system must know which structure
is L4--L5. Level assignment is thus a precondition of per-level grading, and it
has been treated as a problem in its own right since the earliest work
\cite{ghosh2012a,lootus2015automated}. Pan et al.\ illustrate the mature
solution: a Faster R-CNN locates vertebral bodies L1--S1 on sagittal images, the
geometric relationship between sagittal and axial acquisitions identifies the
corresponding disc in each axial slice, and the disc region of interest is then
extracted---each of the three stages reported at 100\% accuracy, with the residual
difficulty falling entirely on the subsequent classification
\cite{pan2021automatically}. Detection has, in short, largely been solved;
grading has not.

\subsection{Semantic and Instance Segmentation of Lumbar Structures}
\label{sec:lr-segmentation}

Dense segmentation provides a stronger anatomical substrate than bounding boxes.
The U-Net architecture \cite{ronneberger2015unet} and its descendants dominate,
as a dedicated systematic review of deep-learning-assisted disc segmentation
confirms \cite{wang2024deep}.
Ghosh and Chaudhary segmented vertebrae, disc, dural sac and background jointly
using cascaded random forests over image and sparsely sampled context features,
achieving Dice 0.87 and 0.84 for dural sac and disc across 212 clinical cases
\cite{ghosh2014supervised}. Han et al.\ introduced adversarial training for
multi-structure segmentation \cite{han2018spinegan}, S\'aenz-Gamboa et al.\
established CNN semantic segmentation of the structural elements surrounding the
spinal cord \cite{senzgamboa2021automatic}, and Al-Kafri et al.\ applied
SegNet to boundary delineation for stenosis assessment across 515 symptomatic
patients, contributing two purpose-built metrics---confidence and
consistency---for assessing the quality of the ground truth itself rather than
only the model \cite{alkafri2019boundary}.

Subsequent work has largely consisted of architectural refinement. Wang et al.\
added multilevel feature-map attention and residual blocks to an Attention U-Net
\cite{wang2022automatic}; Ahmed et al.\ targeted class imbalance between
anatomical structures with a custom combined loss and modified upsampling
\cite{ahmed2024pioneering}; Silveira et al.\ added an Inception module for
multi-scale extraction and a dual-output mechanism, reporting mIoU 0.8974 and
F1 0.9444 on SPIDER and---notably---outperforming TransUNet
\cite{silveira2025automated}. M\"oller et al.\ extended coverage to the whole
spine with SPINEPS, segmenting fourteen structures including the posterior
elements, and established a useful general result: segmenting semantically first
and instance-wise second outperformed training directly on instance segmentation,
beating an nnU-Net baseline \cite{isensee2021nnunet} on every structure and metric
\cite{mller2025spinepsautomatic}. Chen et al.\ merged convolutional and attention
pathways in parallel with a novel position embedding, outperforming fifteen
competing models \cite{chen2024symtc}.

\subsection{Public Segmentation Resources and Benchmarks}
\label{sec:lr-segmentation-data}

Progress in this subfield has been unusually well served by shared data. The
SPIDER dataset provides 447 sagittal T1 and T2 series from 218 patients across
four hospitals, with reference segmentations of vertebrae, discs and spinal
canal plus per-level radiological gradings, together with reference performance
for a baseline algorithm and nnU-Net and a continuous public challenge on a
sequestered test split \cite{vandergraaf2024lumbar}. Several of the segmentation
models above are trained and benchmarked on it
\cite{ahmed2024pioneering,silveira2025automated,patil2026mribased}, which makes
their results mutually comparable in a way that is rare elsewhere in this
literature.

\subsection{Quantitative Morphometry}
\label{sec:lr-morphometry}

Segmentation enables measurement, and measurement offers an escape from ordinal
subjectivity. The dural sac cross-sectional area is the canonical example.
Ghobrial and Roth compared U-Net, Attention U-Net and MultiResUNet for automated
DSCA quantification on 683 scans with 50 held back for external validation;
MultiResUNet performed best, with Pearson correlation 0.9917 and mean absolute
error 23.70~mm$^2$, holding up externally at 99.95\% accuracy and F1 0.9393
\cite{ghobrial2025deep}. Zheng et al.\ pursued the same objective for disc
degeneration, segmenting degeneration-related regions with a Swin-transformer
network and computing signal-intensity and geometric features, then proposing a
quantitative criterion to replace subjective Pfirrmann grading on the basis of
strong correlation with grade across a large population \cite{zheng2022deep}.
Related work measures the lordosis angle from segmentation masks
\cite{ji2026mdvmunet}.

A caution is warranted. Ghobrial and Roth's model relies on T1 weighting alone
and a modest dataset, as the authors acknowledge \cite{ghobrial2025deep}, and
Schizas' original finding---that area measurement over- and under-diagnosed
stenosis relative to morphological grade in a substantial minority of patients
\cite{schizas2010qualitative}---applies equally whether the area is traced by
hand or by network. Automating a measurement does not repair a weak correlation
between that measurement and the clinical state.

\subsection{Evidence that Localisation Precedes Classification}
\label{sec:lr-localize-first}

The strongest recurring methodological finding in this literature is that
explicit anatomical localisation before classification improves grading. The
argument was made structurally by DeepSPINE, which chained U-Net vertebral
segmentation and disc-level designation to a multi-task grading network across
4{,}075 patients and more than 22{,}000 disc-level annotations
\cite{lu2018deepspine}, and it has been corroborated repeatedly since.

Four lines of evidence converge. Bharadwaj et al.\ segmented dural sac and disc
with V-Net, localised facet and foramen by geometric rule, and only then
classified, reaching Dice 0.93 and 0.94 with foramen and facet localisation IoU
0.72 and 0.83 \cite{bharadwaj2023deep}. \v{S}u\v{s}ter\v{s}i\v{c} et al.\ showed
that segmenting, cropping and enhancing the region of interest before
classification produced 0.87 and 0.91 accuracy on axial and sagittal views for a
multi-class herniation problem \cite{sustersic2022a}. Kowlagi et al.\ made the
comparative case most directly, arguing that a well-tuned three-stage pipeline of
segmentation, localisation and classification allows convolutional networks to
outperform transformer approaches on Pfirrmann grading in a population cohort
\cite{kowlagi2022a}. Acharya and Kansakar tested the principle in its strongest
form, hypothesising that anatomical focus rather than architecture is the binding
constraint, and grading severity from a single localised region of interest per
disc \cite{acharya2026disccentric}.

The implication for the present work is direct. If localisation dominates
architecture in determining performance, then a comparison between backbones is
only meaningful when both receive identically localised input---otherwise the
comparison measures preprocessing rather than representation. This is a
methodological requirement that Section~\ref{sec:lr-cnn-vs-transformer} shows
much of the existing comparative literature fails to satisfy.

\subsection{Reduced-Annotation and Unsupervised Alternatives}
\label{sec:lr-weak-supervision}

Dense annotation is expensive, and several strands of work reduce the
requirement. Kuang et al.\ eliminated manual labels entirely for vertebral
segmentation: sub-optimal masks from a rule-based method were combined through a
voting mechanism to supervise CNN training, with a regional strategy restricting
attention to a specific vertebral region, validated across 243 volunteers
\cite{kuang2020mrisegflow}. SpineNet's evidence hotspots represent the
complementary weakly supervised route, obtaining localisation from
classification labels alone \cite{jamaludin2017spinenet}, while the SPIDER
annotation protocol adopted a human-in-the-loop compromise, training on a small
subset to produce initial masks that were then corrected and fed back
\cite{vandergraaf2024lumbar}. Self-supervised pretraining on unlabelled or
metadata-labelled data \cite{jamaludin2017selfsupervised} completes the picture.
Together these establish that the annotation burden, long cited as the principal
obstacle to clinical deployment, is substantially negotiable.

% ===========================================================================
\section{Backbone Architectures for Severity Classification}
\label{sec:lr-backbones}

\subsection{Convolutional Networks}
\label{sec:lr-cnns}

Convolutional networks exploit two properties well matched to imaging: local
connectivity, since a filter responds to a neighbourhood rather than the whole
field; and translation equivariance, since the same filter is applied at every
position. Stacking convolutions builds a hierarchy in which early layers respond
to edges and texture and later layers to composite anatomical structure.

\textbf{ResNet} addressed the degradation encountered when networks are made
deeper, introducing residual blocks whose skip connections allow a layer's input
to bypass subsequent convolutions and be added to their output, providing an
unimpeded path for gradients and permitting stable training at fifty layers and
beyond \cite{he2016deep}. In lumbar imaging ResNet is ubiquitous: as the
classification head in multi-task grading \cite{su2022automatic}, as the
detection backbone in staged pipelines \cite{zhang2023deep}, as the encoder in
segmentation networks \cite{batra2025mscan}, and as the classifier in
SpineNetV2 \cite{windsor2022spinenetv2}.

\textbf{EfficientNet} attacked a different problem---how to allocate a fixed
computational budget---by scaling depth, width and input resolution together
under fixed coefficients derived from a neural architecture search, rather than
scaling one dimension in isolation \cite{tan2019efficientnet}. The resulting
models reach comparable accuracy at substantially lower parameter and FLOP cost,
which matters for the deployment scenarios that motivate this thesis.
EfficientNet appears as the classifier in Kowlagi et al.'s Pfirrmann pipeline
\cite{kowlagi2022a} and in the final stage of M-SCAN \cite{batra2025mscan}.
\textbf{ConvNeXt} represents the most recent convolutional line, modernising the
CNN with design choices borrowed from transformers while retaining convolution
\cite{liu2022convnet}; Udomluck et al.\ include ConvNeXt-Tiny in their feature
extractor comparison and find it stable across classifiers
\cite{udomluck2026multitask}.

The architectural limitation of convolution is a bounded receptive field. Deep
stacks aggregate local responses into progressively broader ones, but the
network never explicitly models a relationship between two distant regions---for
instance, between the L1--L2 and L5--S1 levels, whose degenerative states are
biomechanically coupled.

\subsection{Vision Transformers and the Attention Bottleneck}
\label{sec:lr-vit}

Transformers dispense with convolution and rely on self-attention, in which
every element of the input is weighted against every other regardless of
distance \cite{vaswani2017attention}. The Vision Transformer applies this by
partitioning an image into fixed patches, flattening them and treating the
result as a sequence \cite{dosovitskiy2021image}. Global context is available
from the first layer, and there is no locality prior to overcome---but also no
locality prior to exploit, which is why vision transformers are markedly more
data-hungry than convolutional networks of comparable capacity.

The practical obstacle is cost. Global self-attention scales quadratically with
the number of patches, hence quadratically with image area. Clinical MRI demands
high spatial resolution to resolve structures a few millimetres across, so a
naive vision transformer over full-resolution lumbar images is computationally
prohibitive.

\subsection{The Swin Transformer}
\label{sec:lr-swin}

The Swin Transformer resolves this bottleneck through two mechanisms
\cite{liu2021swin}. \emph{Windowed self-attention} restricts attention to
non-overlapping local windows, making cost linear rather than quadratic in image
size; to restore cross-window communication, the window partition is
\emph{shifted} by half a window between consecutive layers, so information
propagates across boundaries over successive blocks. \emph{Hierarchical feature
maps} are produced by patch-merging layers that progressively reduce spatial
resolution while increasing channel depth, yielding a pyramid analogous to a
CNN's and making the backbone usable for dense prediction as well as
classification.

For lumbar grading the appeal is specific. Severity depends simultaneously on
fine local detail---the texture of a nucleus pulposus, the obliteration of
perineural fat---and on regional context, since a rootlet configuration is
interpretable only relative to the surrounding dural sac and the levels above
and below. A hierarchical, windowed attention mechanism is a plausible fit to
that structure, which is the substantive reason for selecting it as the
transformer representative in this thesis. Swin-derived architectures have been
applied to spinal segmentation, and Zheng et al.\ incorporate Swin components in
their quantitative degeneration network \cite{zheng2022deep}.

\subsection{Hybrid CNN--Transformer Architectures}
\label{sec:lr-hybrids}

A third position holds that the two families are complementary. SymTC merges
convolutional and transformer layers in a \emph{parallel dual-path} arrangement
rather than stacking them sequentially, and integrates a position embedding into
the self-attention module to improve spatial awareness; benchmarked against
fifteen existing models on both a private dataset and a released synthetic one,
it performed best for both vertebral bones and discs \cite{chen2024symtc}. Zheng
et al.\ similarly combine convolutional and Swin components
\cite{zheng2022deep}, and Yi et al.\ pair a 3D ResNet18 backbone with transformer
components across three hospitals \cite{yi2023deep}. Baur et al.\ explore a
different pairing, segmenting discs with a U-Net, reconstructing each as a 3D
point cloud and grading it with a graph neural network; performance was uneven
across levels---F1 0.71 at L2/3 against 0.46 overall---which the authors present
as a foundation requiring refinement rather than a competitive result
\cite{baur2024automated}. Hybrids complicate the
clean comparison this thesis undertakes, but they indicate where the field
expects the answer to lie: not in the wholesale replacement of convolution, but
in supplying it with a mechanism for long-range interaction.

\subsection{Published CNN-versus-Transformer Comparisons and Their Limits}
\label{sec:lr-cnn-vs-transformer}

Several studies bear directly on the architectural question, and their results
do not favour transformers as strongly as the general computer vision literature
might suggest.

Kowlagi et al.\ address it explicitly. Observing that recent gains had been
attributed to transformer architectures, they show that a well-tuned three-stage
convolutional pipeline outperforms the state-of-the-art transformer approaches
on Pfirrmann grading, evaluated on the Northern Finland Birth Cohort 1966---a
population cohort rather than a clinical convenience sample---with an ablation
study and subgroup breakdown, and public code \cite{kowlagi2022a}. Their
conclusion is that pipeline design and tuning outweigh backbone choice.

Lee et al.\ provide the closest published analogue to the present study,
developing a CNN-based and a transformer-based model for lumbar stenosis
classification on 564 MRIs with a 100-study external test set, labelled by four
radiologists with expert consensus as reference and read by eight participants
across four experience levels. Both models exceeded 94\% detection recall; on
dichotomised classification the CNN, transformer and human participants achieved
kappas of 0.99/0.99/0.97--0.98 for central canal, 0.98/0.94/0.81--0.94 for
lateral recesses and 0.98/0.95/0.91--0.95 for neural foramina. The CNN was
superior overall, the transformer similar to superior against clinicians
\cite{lee2026deep}.

Udomluck et al.\ compared VGG19, ConvNeXt-Tiny and DINOv2 features with classical
classifiers under external validation, finding conventional VGG19 features
strongest (accuracy 0.9556 with logistic regression) and DINOv2 features
generalising worst, dropping to 0.6222 with LightGBM
\cite{udomluck2026multitask}. Patil et al.\ reached a still more pointed
conclusion through ablation: radiomic features carried essentially the entire
signal for Pfirrmann grading, and frozen ImageNet-pretrained ResNet50 features
added no measurable discriminative value \cite{patil2026mribased}. Silveira et
al.\ likewise found their enhanced U-Net outperformed TransUNet on segmentation
\cite{silveira2025automated}.

Three methodological limitations recur across this evidence. First, the
compared architectures are rarely given identical preprocessing, localisation
and training budgets, so differences may reflect tuning effort rather than
representational capacity---exactly the confound Kowlagi et al.\ identify
\cite{kowlagi2022a}. Second, the input configuration is usually held fixed,
leaving unexamined the interaction between architecture and available context,
even though a transformer's advantage should plausibly be largest where global
context must substitute for missing views. Third, several comparisons rely on
frozen pretrained features rather than fine-tuned backbones
\cite{udomluck2026multitask,patil2026mribased}, which disadvantages transformers
disproportionately given their weaker inductive biases. These three gaps
motivate the experimental design of this thesis.

\subsection{Transfer Learning, Pretraining, and Self-Supervision}
\label{sec:lr-pretraining}

Medical datasets are small by the standards these architectures were designed
for, so initialisation matters. ImageNet pretraining is near-universal, though
Patil et al.'s ablation is a caution that its benefit is not automatic
\cite{patil2026mribased}. Domain-specific pretraining is better supported:
Jamaludin et al.\ obtained equivalent performance from far fewer labels through
longitudinal self-supervision \cite{jamaludin2017selfsupervised}, and Acharya and
Kansakar used contrastive pretraining on disc-centred crops to reduce sensitivity
to irrelevant appearance variation \cite{acharya2026disccentric,chen2020simple}.
Sequential transfer between related targets---initialising a subarticular model
from spinal canal weights---has also been reported to help on smaller datasets, a
strategy of particular relevance to transformers given their lack of
convolutional inductive bias.

% ===========================================================================
\section{Exploiting Spatial Context: Slices, Volumes, and Views}
\label{sec:lr-context}

\subsection{2D, 2.5D, and 3D Input Representations}
\label{sec:lr-dimensionality}

An MRI series is a volume, but severity is assigned per level, so a
representation must be chosen. \emph{2D} models classify individual slices and
are cheap but discard through-plane context; since a disc spans several slices
and a radiologist integrates across them, a single slice may not contain the
evidence. \emph{3D} models consume the volume directly \cite{yi2023deep} at
considerable memory cost---Acharya and Kansakar note that a standard
high-resolution series contains millions of voxels, making full-volume
processing impractical for many architectures \cite{acharya2026disccentric}.

\emph{2.5D} representations occupy the middle ground and have become the
default in this domain. A small stack of adjacent slices is treated as channels
or as a short sequence, capturing local through-plane context at roughly 2D
cost. SpineNet's input was a stack of nine slices per disc
\cite{jamaludin2017spinenet}; M-SCAN describes its approach explicitly as 2.5D,
selecting the three axial slices closest to each level and modelling them with
recurrent layers \cite{batra2025mscan}; the leading challenge solutions combined
2D CNNs with sequence models or attention-based multiple-instance learning over
the slice dimension.

\subsection{Cross-Sequence and Cross-Plane Fusion}
\label{sec:lr-fusion}

Where multiple sequences are available, they must be combined. The simplest
approach concatenates them as channels, which presumes spatial registration that
sagittal and axial acquisitions do not satisfy. More considered strategies
process each sequence separately and fuse at the feature level.

CrossSpine is the most explicit treatment. Nguyen et al.\ identify two failings
of baselines---single-stream designs cannot integrate complementary information
across sequences, and uniform treatment of all discs ignores level-specific
anatomy---and address the first with a cross-sequence attention mechanism that
fuses features at multiple spatial scales, letting each sequence attend to all
others, followed by a weighting mechanism that prioritises the most
diagnostically informative sequence. They report Macro F1 improved by more than
125\%, Mean AUPRC by 99\% and Mean AUROC by 36\% over their baseline
\cite{nguyen2026crossspine}. M-SCAN applies cross-attention to align sagittal and
axial feature maps before classification \cite{batra2025mscan}, and Park et al.\
pursue transformer-based multimodal fusion on the same benchmark
\cite{park2025multiview}. Zeng et al.\ extend fusion beyond imaging entirely,
combining ConvNeXt image features with GPT-4o-generated textual descriptions
encoded by RoBERTa through a Mamba fusion stage \cite{zeng2025gpt4lfs}.

\subsection{Geometric Correspondence Between Planes}
\label{sec:lr-geometry}

Fusing sagittal and axial data requires knowing which axial slice corresponds to
which sagittal level---a geometric problem solved through DICOM spatial
metadata rather than learned. Pan et al.\ derive the correspondence from the
geometric relationship between the two acquisitions, reporting 100\% accuracy in
identifying the disc in each axial slice \cite{pan2021automatically}. M-SCAN
projects 2D sagittal stenosis coordinates into 3D using patient orientation data
from the DICOM headers, then selects the three nearest axial slices for each of
the five levels \cite{batra2025mscan}.

This step is a hidden cost of multi-sequence modelling, and one that bears
directly on the trade-off examined here. It adds a preprocessing dependency, a
failure mode when metadata is absent or inconsistent, and an engineering burden
that a single-sequence pipeline simply does not incur.

\subsection{Multi-Sequence versus Single-Sequence Input}
\label{sec:lr-multiseq}

The clinical standard is unambiguous that all three sequences contribute
\cite{richards2026the,hallinan2021deep}, and the strongest reported results on
the RSNA benchmark use all three \cite{batra2025mscan}. The countervailing
consideration is cost: multi-sequence processing demands greater memory and
compute, cross-plane registration, and longer inference, and the resulting
infrastructure requirements exceed what many hospital environments provide.

The evidence on whether reduced input suffices is genuinely mixed, which is why
the question remains open. On one side, external validation of sagittal-only
SpineNetV2 found agreement dropping specifically for foraminal stenosis and
herniation---the pathologies requiring transverse context
\cite{nigru2024external}---and Wu et al.\ found ordinal Pfirrmann grading
degrading in upper lumbar discs while binary detection held up
\cite{wu2026external}. On the other, Acharya and Kansakar grade severity from
\emph{sagittal T2 alone}, reaching 78.1\% balanced accuracy with a
severe-to-normal misclassification rate of 2.13\%, on the argument that
anatomical focus rather than input breadth is the binding constraint
\cite{acharya2026disccentric}. Kowlagi et al.\ likewise operate on sagittal input
and outperform transformer baselines \cite{kowlagi2022a}.

Two observations sharpen the question. First, the comparison has not been made
under controlled conditions: the sagittal-only and multi-sequence results above
come from different architectures, datasets, localisation strategies and
evaluation protocols, so the contribution of sequence configuration cannot be
isolated. Second, the expected penalty is target-dependent---severe for
foraminal and subarticular grading, mild for central canal---so a single
aggregate accuracy figure would obscure precisely the structure of interest.
Both observations inform the design adopted in Chapter~3.

\subsection{Robustness to Missing or Incomplete Sequences}
\label{sec:lr-missing}

Clinical data is incomplete: sequences are omitted, truncated or corrupted by
artifact. The LumbarDISC inclusion criteria required all three sequences and
excluded studies with substantial artifact \cite{richards2026the}, so a model
trained on it has not been exposed to the degraded conditions of routine
practice. A model that requires three sequences fails outright when one is
absent, whereas a single-sequence model is unaffected---an argument for reduced
input configurations that is about robustness rather than efficiency, and one
that the literature on missing-modality learning addresses only in passing for
this application.

% ===========================================================================
\section{Class Imbalance and Ordinal Severity Modelling}
\label{sec:lr-imbalance}

\subsection{Severity Distributions in Lumbar MRI Datasets}
\label{sec:lr-distributions}

Class imbalance in this domain is not an artefact of sampling but a property of
the population: in any unselected clinical cohort, most disc levels are normal or
mildly degenerate, and severe compression is uncommon. The RSNA LumbarDISC
distributions make the scale explicit \cite{richards2026the}. Of 13{,}474 spinal
canal stenosis grades, 85.4\% are normal or mild, 8.8\% moderate and 5.9\%
severe. Of 26{,}919 neural foraminal grades, roughly 78\% are normal or mild,
18\% moderate and 4.5\% severe. Subarticular grading is the least skewed of the
three, with roughly 69\% normal or mild, 19\% moderate and 11\% severe across
26{,}285 grades.

The same pattern recurs elsewhere. Liawrungrueang et al.\ report a Pfirrmann
distribution of 220 grade~I, 530 grade~II, 170 grade~III, 160 grade~IV and just
20 grade~V across 1{,}000 images \cite{liawrungrueang2023automatic}; Niemeyer et
al.\ describe an approximately Gaussian grade distribution with grades~1 and~5
under-represented \cite{niemeyer2021a}. Two consequences follow. First, a model
trained under unweighted cross-entropy can achieve roughly 85\% accuracy on
spinal canal stenosis by predicting \emph{normal/mild} unconditionally, so
aggregate accuracy is close to uninformative on this task. Second, the ordering
of imbalance across targets is the inverse of the ordering of difficulty:
subarticular stenosis has the most balanced label distribution yet the poorest
human reliability \cite{lurie2008reliability}, confirming that its difficulty is
intrinsic rather than a consequence of scarce positive examples.

\subsection{Data-Level Strategies}
\label{sec:lr-data-level}

The most common data-level intervention is weighted random sampling, in which
the probability of drawing a minority-class example is inflated so that each
minibatch approaches a uniform severity distribution and gradients are not
dominated by the majority class. Aggressive augmentation of minority classes
complements this, since with few severe examples a network may memorise them
rather than learn the morphology of severe compression. Reported techniques
include contrast-limited adaptive histogram equalisation, random scaling, affine
rotation and elastic deformation, alongside mixing methods.

Tsai et al.\ provide the clearest single demonstration of augmentation's value
in the small-data regime, treating it as the subject of the study rather than an
implementation detail: mean average precision reached 92.4\% at 550 images with
augmentation, and augmentation was decisive in preventing both under- and
over-fitting \cite{tsai2021lumbar}. Ahmed et al.\ attack a related form of
imbalance---between anatomical structures rather than severity classes---through
preprocessing that rectifies class inconsistencies before training
\cite{ahmed2024pioneering}.

\subsection{Loss-Level Strategies}
\label{sec:lr-loss-level}

Class-weighted cross-entropy multiplies each class's contribution to the loss by
a scalar, forcing the optimiser to penalise minority-class errors more heavily.
The RSNA challenge encoded this directly into its evaluation metric, assigning
weights of 1, 2 and 4 to normal/mild, moderate and severe respectively, so that
aligning the training objective with the evaluation criterion is a matter of
adopting the same coefficients.

Focal loss goes further, introducing a modulating factor that down-weights
well-classified examples dynamically and concentrates gradient on hard,
borderline cases \cite{lin2017focal}. This is well matched to the present
problem, where the abundant normal cases are mostly easy and the difficulty
concentrates at the mild/moderate and moderate/severe boundaries---precisely
where human readers also disagree. Acharya and Kansakar combine weighted focal
loss with contrastive pretraining \cite{acharya2026disccentric}, and Islam~Sk et
al.\ propose an adaptive PID-Tversky loss that borrows feedback control to adjust
training penalties dynamically toward under-segmented minority instances
\cite{islamsk2026an}.

\subsection{Ordinal-Aware Objectives and Whether They Help}
\label{sec:lr-ordinal-losses}

Given the ordinal label structure of Section~\ref{sec:lr-ordinality}, it is
natural to expect ordinal-aware objectives to outperform categorical ones. The
evidence does not support this expectation as strongly as the intuition
suggests.

Niemeyer et al.\ conducted the most systematic test available, evaluating soft
kappa loss, ordinal cross-entropy and regression losses against standard
cross-entropy for Pfirrmann grading on 1{,}599 patients and 7{,}948 discs, with
extensive hyperparameter optimisation, and additionally trying class pooling and
class-weighted losses to counter imbalance. None of these improved classification
performance overall. The default cross-entropy classifier reached Cohen
$\kappa = 0.92$, average sensitivity 90.2\% and precision 92.5\%, with
predictions within one grade of ground truth in 99.2\% of cases and a mean
absolute error of 0.08 grades \cite{niemeyer2021a}. Patil et al.\ list ordinal
regression among the supplementary analyses they ran, and still report that
intermediate-grade discrimination remains the outstanding difficulty
\cite{patil2026mribased}.

Two readings are possible. Either the ordinal structure is already implicitly
captured---a network trained with cross-entropy on visually ordered classes will
naturally confuse adjacent grades more often than distant ones---or the specific
formulations tested were inadequate. The literature does not currently
distinguish these, which is why this thesis treats the choice of objective as an
empirical question rather than a settled one.

\subsection{Contrastive and Representation-Learning Approaches}
\label{sec:lr-contrastive}

A third strategy attacks imbalance in the representation rather than in the
sampling or the loss. Contrastive learning pulls examples of similar severity
together in the latent space and pushes dissimilar ones apart before any
classifier is fitted \cite{chen2020simple}, so the representation is shaped by
the structure of the problem rather than by class frequency.

Acharya and Kansakar apply this directly, combining contrastive pretraining with
disc-level fine-tuning over a single anatomically localised region of interest
per disc, on the reasoning that this forces attention onto intrinsic disc
features and reduces sensitivity to irrelevant appearance variation. An auxiliary
regression task handles disc localisation and a weighted focal loss addresses the
residual imbalance. The result---78.1\% balanced accuracy with the
severe-to-normal misclassification rate reduced to 2.13\% relative to supervised
training from scratch---is notable less for the headline accuracy than for what
it optimises \cite{acharya2026disccentric}. The lineage runs back to Jamaludin et
al.'s longitudinal self-supervision \cite{jamaludin2017selfsupervised}.

\subsection{The Clinical Asymmetry of Error}
\label{sec:lr-error-asymmetry}

The most important point in this section is that not all errors are equal, and
that standard metrics conceal the difference. Grading a severe stenosis as
normal risks a missed diagnosis and delayed decompression; grading a normal disc
as moderate risks an unnecessary review. Overall accuracy, macro-F1 and even
Cohen's $\kappa$ treat these identically.

Three studies take the asymmetry seriously. Acharya and Kansakar report the
severe-to-normal misclassification rate as a headline result alongside balanced
accuracy, arguing explicitly that conventional metrics do not capture the
differential clinical risk \cite{acharya2026disccentric}. Wu et al.\ characterise
SpineNetv2's error profile as specificity-oriented, with false negatives
exceeding false positives, and conclude that the system is suitable as a
confirmatory second reader but that reliance on its negative findings is not
recommended---a conclusion about deployment that follows from the shape of the
errors rather than their number \cite{wu2026external}. The RSNA weighted metric
encodes the same judgement into the benchmark itself.

Ishimoto et al.\ supply the complementary observation: collapsing four grades to
a severe-versus-rest dichotomy raised agreement from 65.7\% to 94.1\% with
$\kappa = 0.75$ \cite{ishimoto2020could}. If the clinically decisive question is
whether a level is severe, then performance on that question is considerably
better than the multi-class figures suggest, and evaluation should report both.
This chapter therefore recommends, and Chapter~3 adopts, reporting per-class
recall for the severe class alongside aggregate measures.

% ===========================================================================
\section{Datasets and Benchmarks}
\label{sec:lr-datasets}

\subsection{Private and Single-Institution Cohorts}
\label{sec:lr-private-data}

Most work before 2024 used private single-institution data, with cohort sizes
spanning three orders of magnitude---from 75 exams \cite{gao2022automatic} and
146 consecutive patients \cite{lehnen2021detection} through 200 studies
\cite{bharadwaj2023deep} and 542 images \cite{won2020spinal} to 4{,}075 patients
\cite{lu2018deepspine} and 15{,}254 axial images \cite{su2022automatic}.
Aggregate performance correlates only loosely with cohort size, because larger
private datasets are frequently drawn from a single scanner and protocol.

Two consequences follow, and both are structural rather than incidental. Results
are not comparable across studies, since each is measured against a different
label distribution and reference standard; and models cannot be independently
re-evaluated, since the data cannot be shared. Both systematic reviews of the
period identify this as the field's central methodological weakness
\cite{compte2023are,mendes2026deep}.

\subsection{Population and Epidemiological Cohorts}
\label{sec:lr-population-data}

A smaller body of work uses population cohorts, which differ from clinical
convenience samples in a way that matters for generalisation: participants are
not selected by referral, so the severity distribution reflects the population
rather than the imaged-because-symptomatic subset. The Northern Finland Birth
Cohort 1966, comprising lumbar MRI from participants imaged at age 46--47, is the
most used \cite{mcsweeney2023external,kowlagi2022a}; the Wakayama Spine Study
\cite{ishimoto2020could} and the Hong Kong Disc Degeneration Population-Based
Cohort \cite{cheung2022learningbased} serve similar roles. Kowlagi et al.\ are
explicit that evaluating on a population cohort rather than a clinical sample is
a deliberate methodological choice, reflecting the transferability of results to
downstream research \cite{kowlagi2022a}.

\subsection{Public Lumbar MRI Datasets}
\label{sec:lr-public-data}

Public data arrived late to this field relative to CT, where large annotated
vertebral collections have existed for years. SPIDER is the principal
segmentation resource: 447 sagittal T1 and T2 series from 218 patients across
four hospitals, with reference segmentations of vertebrae, discs and spinal
canal, per-level radiological gradings, baseline and nnU-Net reference
performance, and a continuous challenge on a sequestered split
\cite{vandergraaf2024lumbar}. Its influence is visible in the number of
subsequent methods benchmarked on it
\cite{ahmed2024pioneering,silveira2025automated,patil2026mribased,ghobrial2025deep}.
Smaller public collections support classification work
\cite{natalia2024lumbar,shahzadi2023nerve,nikpasand2024automated}.

\subsection{The RSNA 2024 Challenge and LumbarDISC}
\label{sec:lr-lumbardisc}

The RSNA Lumbar Degenerative Imaging Spine Classification dataset is the largest
and most diverse expertly annotated lumbar spondylosis collection assembled, and
it is the dataset used in this thesis \cite{richards2026the}. It comprises MRI
studies from 2{,}697 patients and 8{,}593 image series drawn from eight
institutions across six countries and five continents---a geographic spread
chosen deliberately so that models trained on it face genuine acquisition
heterogeneity rather than a single site's conventions.

Inclusion required a sagittal T2-like series (conventional spin echo T2, STIR or
Dixon), a sagittal T1 series and an axial T2 series covering at least three
lumbar disc levels, from outpatients aged 18 or over imaged for degenerative
disease. Studies with lumbosacral hardware, active non-degenerative pathology,
severe scoliosis or substantial artifact were excluded \cite{richards2026the}.
Annotation was performed by expert volunteer neuroradiologists and
musculoskeletal radiologists recruited through the ASNR
\cite{rsna2024rsna}. The associated challenge \cite{rsna2024rsna} made the
dataset the field's common benchmark and shifted attention decisively from
whole-image classification toward anatomically localised, per-level severity
grading.

\subsection{Label Structure, Distribution, and the Weighted Metric}
\label{sec:lr-label-structure}

The label structure is the defining feature of the benchmark and warrants precise
statement. Each study receives twenty-five graded outputs: five conditions
---spinal canal stenosis, left and right neural foraminal narrowing, left and
right subarticular stenosis---at each of five disc levels from L1/L2 to L5/S1.
Each is assigned one of three ordinal classes: Normal/Mild, Moderate or Severe.
Expert coordinate annotations accompany the grades, so anatomical localisation
can be learned in a supervised manner rather than inferred
\cite{rsna2024rsna,richards2026the}.

Three properties follow. The task is \emph{not} generic three-class image
classification, and results from studies that treat it as such are not
comparable. The severity distribution is severely imbalanced and differently so
for each condition (Section~\ref{sec:lr-distributions}). And the challenge metric
is a sample-weighted log loss with weights of 1, 2 and 4 across the three
classes, which encodes the clinical asymmetry of
Section~\ref{sec:lr-error-asymmetry} directly into the evaluation. A model tuned
for unweighted accuracy is therefore optimising a different objective from the
one the benchmark measures.

\subsection{Challenge Solutions and the Practices They Standardised}
\label{sec:lr-challenge-solutions}

The solutions converged on a common template, and that convergence is itself a
finding. Nearly all adopted multi-stage pipelines in which localisation precedes
classification: a detection model---commonly a U-Net, YOLO variant or Faster
R-CNN with a feature pyramid---identifies each disc level on the sagittal
series, and a second-stage classifier grades tightly cropped regions, typically
using 2.5D representations that combine a 2D backbone with a sequence model or
attention over adjacent slices.

M-SCAN exemplifies the approach and reports AUROC 0.971 on 1{,}975 studies,
using a U-Net with ResNet50 encoder to locate stenosis points on sagittal images,
projecting these into 3D via DICOM orientation metadata to select corresponding
axial slices, and applying cross-attention to align sagittal and axial features
\cite{batra2025mscan}. Liu et al.\ pair MobileNetV2 with U-Net for efficiency,
reporting 94.93\% accuracy under five-fold cross-validation
\cite{liu2025automated}; Park et al.\ apply transformer-based multimodal fusion
\cite{park2025multiview}; and further challenge-derived work continues to appear
\cite{patel2025automated,garciadecelis2025deep}.

Two cautions attach to these figures. Most are cross-validation results on the
public training split rather than external validation, so they measure fit to
LumbarDISC rather than transportability beyond it. And the top solutions rely on
large ensembles, which inflate reported accuracy relative to what a single
deployable model achieves---a gap directly relevant to the efficiency question
this thesis examines.

% ===========================================================================
\section{Evaluation, Validation, and Generalization}
\label{sec:lr-validation}

\subsection{Metrics and Their Appropriateness}
\label{sec:lr-metrics}

Reported metrics in this field are heterogeneous and frequently ill-matched to
the task. Overall accuracy remains the most common headline figure despite being
close to meaningless under the severity distributions of
Section~\ref{sec:lr-distributions}. Area under the ROC curve is widely reported
\cite{batra2025mscan,wang2024machine} but is defined for binary problems and must
be extended by one-versus-rest averaging for three-class grading, which conceals
which class is failing.

More appropriate choices appear in the better studies. Cohen's and Gwet's
$\kappa$ measure agreement above chance and are directly comparable with the
human reliability figures of Section~\ref{sec:lr-reliability}, which is the
correct reference point \cite{lim2022improved,su2022automatic,mcsweeney2023external}.
Balanced accuracy corrects for prevalence \cite{mcsweeney2023external,acharya2026disccentric}.
Mean absolute error in grades exploits the ordinal structure and distinguishes a
one-grade slip from a two-grade error \cite{niemeyer2021a,wu2026external}.
Ishimoto et al.\ report both exact-grade agreement and the collapsed
severe-versus-rest figure, making the difference visible rather than choosing
between them \cite{ishimoto2020could}.

The recommendation that follows is to report a small panel rather than a single
number: a prevalence-corrected aggregate, an agreement statistic comparable to
inter-reader values, per-class recall for the severe class, and the distribution
of error magnitudes across grade boundaries.

\subsection{Internal Validation and Optimistic Reporting}
\label{sec:lr-internal-validation}

A recurring pattern is that internally validated performance is high and
externally validated performance is materially lower. Several reported accuracies
in the reviewed literature exceed 95\%
\cite{shahzadi2023nerve,liawrungrueang2023automatic,liu2025automated,tabarestani2026a},
which sits uneasily beside inter-reader agreement of $\kappa = 0.49$ for
subarticular stenosis \cite{lurie2008reliability}. Where a model appears to
exceed the reproducibility of the process that generated its labels, the
explanation usually lies in the evaluation protocol: single-institution data,
cross-validation rather than held-out external testing, labels from a single
reader, or a task collapsed to a binary decision.

The point is not that these studies are wrong on their own terms, but that
their numbers do not transfer. Compte et al.\ state the general finding
directly: algorithms did not perform as well in replication or external
validation as in their development studies \cite{compte2023are}.

\subsection{External Validation and Domain Shift}
\label{sec:lr-generalization}

External validation on geographically and temporally separate cohorts is the
strongest available evidence of clinical viability, and a small number of
studies supply it.

McSweeney et al.\ evaluated SpineNet on 1{,}331 participants from the Northern
Finland Birth Cohort 1966, a dataset separate from its UK training data in both
geography and time. Balanced accuracy was 78\% for disc degeneration and 86\%
for Modic changes, with Pfirrmann reliability against expert radiologists of Lin
concordance 0.86 and Cohen $\kappa = 0.68$, and Modic detection at
$\kappa = 0.74$; both were essentially unchanged in a low-back-pain subset. The
informative detail is the error structure: 20.83\% of discs were graded
differently from human raters, but only 0.85\% differed by more than one grade
\cite{mcsweeney2023external}. Nigru et al.\ evaluated SpineNetV2 across eleven
features on 1{,}747 discs from 353 patients, finding agreement above 0.7 for most
targets but distinctly lower for foraminal stenosis and herniation, attributed
to the limits of sagittal input \cite{nigru2024external}. Wu et al.\ tested
SpineNetv2 on 491 patients and 2{,}455 discs, reporting overall accuracy of
83.5--97.5\% and superiority over a junior orthopaedic surgeon on canal stenosis,
spondylolisthesis and bilateral foraminal stenosis \cite{wu2026external}. Grob et
al.\ contribute a further independent validation on 882 scans
\cite{grob2022external}.

Studies that build external testing into their development are still uncommon
but instructive. Tumko et al.\ evaluated on 150 studies graded by a panel of
seven radiologists with majority voting and external adjudication
\cite{tumko2024a}; Su et al.\ used a 1{,}273-image external set from a different
hospital and observed accuracy falling by roughly 7--10 points relative to
internal testing \cite{su2022automatic}; Zhang et al.\ found classification
accuracy dropping from 87.70\% to 74.23\% externally \cite{zhang2023deep}; Zeng
et al.\ report an unusually small external drop, from 93.7\% to 92.2\%
\cite{zeng2025gpt4lfs}.

\subsection{Detection Versus Fine-Grained Grading}
\label{sec:lr-detection-vs-grading}

A consistent and under-appreciated finding is that these two capabilities
generalise differently. Zhang et al.\ provide the clearest illustration:
moving to an external institution, detection degraded only modestly---mean IoU
0.82 to 0.70, precision 98.4\% to 96.3\%, sensitivity 99.4\% to 97.8\%---while
classification accuracy fell from 87.70\% to 74.23\% \cite{zhang2023deep}. Wu et
al.\ report the same asymmetry in different terms: binary detection across five
pathologies held up well while ordinal Pfirrmann grading remained the limiting
factor, degrading further in older patients and upper lumbar discs
\cite{wu2026external}. Yilihamu et al.\ show the effect within a single study as
task granularity increases: segmentation IoU remained above 97\% externally,
while classification precision fell from 81.21\% to 74.50\% for an 18-category
problem yet held at 92.51\% to 90.07\% for four-category severity
\cite{yilihamu2025quantification}.

The implication is that where a system fails is predictable. Finding the anatomy
is robust; judging how bad it looks is not. Reporting a single aggregate figure
for a pipeline that performs both obscures this, and evaluation should separate
the two stages.

\subsection{Benchmarking Against Clinicians}
\label{sec:lr-clinician-benchmark}

Comparison against human readers is the most clinically meaningful benchmark, and
the better studies construct it carefully. Won et al.\ provide the cleanest
design: two experts each labelled the data, a model was trained on each expert's
labels, and the key result was that the difference between the two experts was
statistically significant while the difference between each expert and their
corresponding model was not \cite{won2020spinal}. Tumko et al.\ compared against
a seven-radiologist panel average and found the model matched or exceeded it on
all three stenosis types \cite{tumko2024a}. Lee et al.\ benchmarked against eight
participants spanning four experience levels \cite{lee2026deep}, and Wu et al.\
against orthopaedic surgeons of differing seniority \cite{wu2026external}.

Two cautions apply. Comparisons against junior readers or against surgeons rather
than radiologists set a lower bar than subspecialty radiological practice
\cite{ke2025integrating,wu2026external}. And readers in these studies typically
work from cropped images without clinical history, which is not how they read in
practice---so such comparisons measure a narrower task than clinical
interpretation.

\subsection{Longitudinal Stability and Reproducibility}
\label{sec:lr-longitudinal}

One-shot accuracy does not establish that a model is stable over repeat imaging,
which matters for monitoring progression. Murto et al.\ supply the only
longitudinal evidence in this collection, comparing SpineNet's Pfirrmann grading
against radiologists in 19 volunteers imaged at age 37 and again at 51.
Radiologist-to-SpineNet concordance was 0.54 at both timepoints, clearly below
the 0.83 and 0.78 observed between radiologists, and SpineNet systematically
assigned grade~1 to several discs and grade~5 to more discs than the readers did
\cite{murto2025comparison}. The cohort is small, but a systematic offset that is
stable across time is exactly the failure mode that one-shot evaluation cannot
detect. Cheung et al.\ approach the temporal dimension differently, predicting
five-year progression from baseline imaging across 1{,}343 participants
\cite{cheung2022learningbased}.

\subsection{Reporting Quality and Recurring Weaknesses}
\label{sec:lr-reporting}

The systematic reviews converge on a consistent diagnosis. Wang et al.\ pooled
twelve studies covering 15{,}044 patients, obtaining pooled sensitivity 0.84 and
specificity 0.87 with SROC AUC 0.92, but with heterogeneity of $I^2$ near 99\%
and four of twelve studies at high risk of bias; they concluded that no current
model is reliable enough for clinical practice, citing the need for accepted
reference standards, larger datasets, external validation and fuller reporting
\cite{wang2024machine}. Compte et al.\ found no significant difference between
algorithm families and identified the development-to-validation performance drop
as the field's central problem \cite{compte2023are}. Mendes et al.\ found that
most of 56 segmentation studies relied on private datasets and lacked external
validation \cite{mendes2026deep}. Liawrungrueang et al.\ report metrics spanning
71.5\% to 99\% across twenty studies---a range so wide as to indicate that the
underlying tasks are not comparable \cite{liawrungrueang2025artificial}.

Four weaknesses recur: single-institution data without external testing; ground
truth from a single reader with no adjudication; task definitions that vary
between studies while sharing a label vocabulary; and metric selection that
flatters imbalanced problems. The present work adopts the public LumbarDISC
benchmark, its panel-derived labels and its weighted metric specifically to
avoid the first, second and fourth of these.

% ===========================================================================
\section{Interpretability and Clinical Integration}
\label{sec:lr-interpretability}

\subsection{Saliency, Evidence Hotspots, and Grad-CAM}
\label{sec:lr-saliency}

The earliest interpretability contribution in this domain is SpineNet's evidence
hotspots, which map output-layer gradients back to the input to produce a heatmap
per radiological score, obtained without any localisation supervision
\cite{jamaludin2017spinenet}. Grad-CAM \cite{selvaraju2017gradcam} has since
become the default. Tabarestani et al.\ use it to confirm that their cascade
attends to clinically relevant regions \cite{tabarestani2026a}, and Silveira et
al.\ report that attention concentrated on vertebral bodies and nucleus pulposus
in spatial correspondence with manual segmentations \cite{silveira2025automated}.

The value of these methods should not be overstated. A saliency map shows where
a network responded, not why, and confirming that attention falls on plausible
anatomy is a weak check---it excludes gross reliance on background artefact but
does not establish that the model has learned the morphological criteria the
grading system specifies.

\subsection{Inherently Interpretable Models}
\label{sec:lr-interpretable}

A stronger position is to build models whose reasoning is legible by
construction. Bharadwaj et al.\ provide the most striking result in this
literature: alongside a Big Transfer classifier they trained a decision-tree
classifier operating on segmentation-derived features for central canal stenosis,
and the interpretable model \emph{outperformed} the black box, reaching
$\kappa = 0.80$ against 0.54 and approaching inter-reader agreement of 0.74--0.86
\cite{bharadwaj2023deep}. Quantitative morphometry belongs to the same family:
a measured dural sac area \cite{ghobrial2025deep} or a quantitative degeneration
criterion \cite{zheng2022deep} is inspectable in a way a logit is not. Patil et
al.\ make the related point that radiomic features carried the entire
discriminative signal for Pfirrmann grading while deep features added nothing
measurable, and that radiomics additionally yields interpretable shape and
intensity descriptors usable as imaging biomarkers \cite{patil2026mribased}.

Together these findings undercut the assumption that interpretability must be
purchased at the cost of accuracy---at least for targets whose clinical
definition is already explicitly morphological.

\subsection{Vision--Language Models and Report Generation}
\label{sec:lr-vlm}

The most recent direction outputs prose rather than labels. Islam~Sk et al.\
present an explainable vision--language framework that classifies severity,
segments the anatomy and generates radiologist-style reports, built on
biomedical foundation models, with a spatial patch cross-attention module that
preserves the anatomical hierarchies global pooling discards and an adaptive
PID-Tversky loss for class imbalance; they report 90.69\% classification
accuracy, macro-averaged Dice 0.9512 and CIDEr 92.80\%
\cite{islamsk2026an}. Zeng et al.\ take a narrower approach, fusing ConvNeXt
image features with GPT-4o-generated descriptions for foraminal stenosis
classification across three centres and seven scanners \cite{zeng2025gpt4lfs}.
These systems are early and their language metrics measure fluency rather than
clinical correctness, but they indicate where the field is heading.

\subsection{Reader-Assistance Studies and Workflow Impact}
\label{sec:lr-reader-assistance}

The most consequential evidence for clinical value comes not from model metrics
but from measuring what changes when readers use the model. Lim et al.\ conducted
the definitive study of this kind: eight radiologists with 2--13 years of
experience read the same studies with and without deep-learning assistance,
separated by a one-month washout, against a 32-year-veteran musculoskeletal
radiologist as reference. Interpretation time per study fell from 124--274
seconds to 47--71 seconds, and interobserver agreement was superior or equivalent
for every stenosis grading. The largest gain fell to general and in-training
radiologists on four-class foraminal stenosis, where $\kappa$ rose from 0.39 to
0.71 and 0.70 \cite{lim2022improved}. Gao et al.\ observed a comparable effect
for Modic assessment, where AI assistance raised agreement between a junior
radiologist and a senior neuroradiologist from $\kappa = 0.52$ to 0.58
\cite{gao2022automatic}.

The pattern in both is that assistance helps least-experienced readers most and
standardises rather than replaces judgement. This reframes the objective: the
target is not a model that outperforms the best radiologist, but one that raises
the floor.

\subsection{Uncertainty and Calibration}
\label{sec:lr-uncertainty}

Uncertainty estimation is conspicuously underdeveloped in this literature. Where
errors concentrate at grade boundaries and where the reference standard is itself
unreliable, a calibrated confidence would let borderline cases be referred rather
than silently graded---a natural fit for triage. Almost none of the reviewed work
reports calibration. Patil et al.\ are the exception, applying isotonic
calibration on a validation split and reporting Brier scores alongside
discrimination metrics \cite{patil2026mribased}. Since modern networks are
systematically overconfident \cite{guo2017calibration}, an uncalibrated severity
probability should not be interpreted as a clinical confidence, and this remains
an open opportunity rather than a solved problem.

% ===========================================================================
\section{Synthesis and Research Gap}
\label{sec:lr-gap}

\subsection{What the Literature Has Established}
\label{sec:lr-established}

Five findings are supported well enough to be treated as settled.

\textbf{Anatomical localisation precedes classification.} From DeepSPINE
\cite{lu2018deepspine} through the interpretable pipelines of Bharadwaj et al.\
\cite{bharadwaj2023deep} to essentially every leading RSNA challenge solution
\cite{batra2025mscan}, systems that localise the disc level before grading it
outperform those that classify whole images. Detection is close to solved, with
several studies reporting near-perfect level assignment
\cite{pan2021automatically,ghosh2012a}.

\textbf{Automated grading is competitive with human reliability, because human
reliability is low.} Inter-reader agreement is substantial for central canal
stenosis but only moderate for foraminal and poor for subarticular grading
\cite{lurie2008reliability}. Against that benchmark, models sit inside the band
of human disagreement \cite{won2020spinal,mcsweeney2023external,tumko2024a}.
This is a real result, but it is a statement about the modesty of the reference
standard as much as the capability of the models.

\textbf{Errors concentrate at adjacent grade boundaries.} Disagreement above one
grade is rare---0.85\% of discs in the largest external validation
\cite{mcsweeney2023external}---and collapsing to a severe-versus-rest decision
raises agreement substantially \cite{ishimoto2020could}. The difficulty lies in
boundary placement, not gross misclassification.

\textbf{Detection generalises across institutions; fine-grained grading does
not.} The asymmetry is consistent and large
\cite{zhang2023deep,wu2026external,yilihamu2025quantification}.

\textbf{Class imbalance must be addressed explicitly.} With 85.4\% of spinal
canal grades normal or mild \cite{richards2026the}, unweighted objectives yield
degenerate solutions, and weighted sampling, weighted or focal losses and
contrastive pretraining are all established mitigations
\cite{lin2017focal,acharya2026disccentric}.

\subsection{What Remains Contested}
\label{sec:lr-contested}

Three questions remain genuinely open, and the disagreement is substantive
rather than terminological.

\emph{Does architectural sophistication help?} The general computer vision
literature would predict transformers to outperform convolutional networks given
sufficient data, but the lumbar evidence does not bear this out. Kowlagi et al.\
report a tuned CNN pipeline beating transformer approaches \cite{kowlagi2022a};
Lee et al.\ find a CNN superior and a transformer merely comparable
\cite{lee2026deep}; Silveira et al.\ beat TransUNet with an enhanced U-Net
\cite{silveira2025automated}; Udomluck et al.\ find conventional VGG19 features
generalising better than DINOv2 \cite{udomluck2026multitask}; and Patil et al.\
find deep features adding nothing over radiomics \cite{patil2026mribased}. Set
against this, hybrid architectures report gains
\cite{chen2024symtc,zheng2022deep}, and cross-sequence attention reports large
ones \cite{nguyen2026crossspine}.

\emph{Does the ordinal structure of the labels admit exploitation?} Errors
demonstrably cluster at adjacent grades, yet the most systematic test of
ordinal-aware objectives found none improved on plain cross-entropy
\cite{niemeyer2021a}.

\emph{How much imaging is actually required?} Clinical practice and the strongest
benchmark results use all three sequences
\cite{richards2026the,batra2025mscan}, while sagittal-only systems report
competitive performance \cite{acharya2026disccentric,kowlagi2022a} and
sagittal-only systems also fail in a specific, predicted way on the lateral
compartments \cite{nigru2024external}.

\subsection{Gap 1: Architecture Comparison Under Controlled Conditions}
\label{sec:lr-gap-architecture}

The comparative studies above are informative but none isolates the variable of
interest. Three confounds recur.

First, \emph{unequal pipelines}. Compared architectures rarely receive identical
preprocessing, localisation and training budgets, so a reported difference may
reflect engineering effort rather than representational capacity---precisely the
argument Kowlagi et al.\ advance when attributing their CNN's advantage to
pipeline design rather than to convolution as such \cite{kowlagi2022a}. Given the
evidence of Section~\ref{sec:lr-localize-first} that localisation dominates
performance, a comparison in which one arm is better localised than the other
measures preprocessing.

Second, \emph{frozen features}. Several comparisons extract features from frozen
pretrained backbones rather than fine-tuning end to end
\cite{udomluck2026multitask,patil2026mribased}. This systematically
disadvantages transformers, whose weaker inductive biases make them more
dependent on in-domain adaptation.

Third, \emph{fixed input configuration}. Every comparison holds the sequence
configuration constant, so the interaction between architecture and available
context is untested. This matters because the theoretical case for global
attention is strongest exactly where context must be inferred rather than
observed---that is, where a view is missing.

A controlled comparison therefore requires a shared localisation stage, a shared
training and augmentation protocol, comparable capacity and tuning budgets, and
end-to-end fine-tuning of both backbones, evaluated on the same benchmark with
the same metric. No study in this review satisfies all four conditions.

\subsection{Gap 2: Sequence Configuration and the Accuracy--Efficiency Trade-off}
\label{sec:lr-gap-sequence}

The second gap concerns input rather than model. Multi-sequence pipelines carry
costs that are rarely quantified alongside their accuracy: greater memory and
compute, cross-plane geometric registration dependent on DICOM metadata
\cite{batra2025mscan,pan2021automatically}, longer inference, and brittleness
when a sequence is missing or degraded (Section~\ref{sec:lr-missing}). Against
these, the incremental diagnostic yield of each additional sequence has not been
measured under controlled conditions.

The literature supplies both sides of the argument but no direct test. Evidence
that the third sequence matters is indirect, inferred from the failure of
sagittal-only systems on foraminal and subarticular targets
\cite{nigru2024external,wu2026external}; evidence that it may not is likewise
indirect, resting on sagittal-only systems that perform respectably against
different baselines on different data
\cite{acharya2026disccentric,kowlagi2022a}. Because these come from different
architectures, cohorts, localisation strategies and metrics, the contribution of
sequence configuration cannot be separated from everything else that differs.

Two design requirements follow. The comparison must hold architecture and
pipeline constant while varying only the input configuration; and it must report
results per condition rather than in aggregate, since the expected penalty is
concentrated in the foraminal and subarticular targets and an aggregate figure
would average away the structure of interest.

\subsection{Research Questions and Objectives}
\label{sec:lr-questions}

The two gaps combine into a single factorial investigation. Crossing
architecture with input configuration answers both questions and, additionally,
tests whether they interact---whether a transformer's global attention can
partially compensate for absent multiplanar data, which neither gap addresses
alone. This yields the following research questions:

\begin{description}
  \item[RQ1] Under matched localisation, training and evaluation protocols, how
  do convolutional backbones (ResNet50, EfficientNet-B4) compare with a
  hierarchical vision transformer (Swin) for per-level, per-condition severity
  grading of lumbar degenerative disease?

  \item[RQ2] What is the cost in diagnostic performance of reducing the input
  from the full multi-sequence configuration (Sagittal T1, Sagittal T2/STIR,
  Axial T2) to a single-sequence configuration (Sagittal T2), and how is that
  cost distributed across the five graded conditions?

  \item[RQ3] Do architecture and sequence configuration interact---specifically,
  is the penalty for single-sequence input smaller for the transformer than for
  the convolutional backbones?

  \item[RQ4] Which class-imbalance strategies most effectively reduce
  clinically consequential errors, in particular severe cases graded as
  normal or mild?
\end{description}

The corresponding objectives are to implement a shared localisation and
preprocessing stage so that all arms receive identical input regions; to train
each architecture under both configurations with matched budgets; to evaluate
using the benchmark's weighted metric together with per-condition and per-class
reporting; and to quantify computational cost---parameters, memory and inference
time---so that any accuracy difference can be weighed against it.

\subsection{Scope and Delimitations}
\label{sec:lr-scope-delimit}

Four boundaries are stated explicitly. The work addresses \emph{severity
grading}, taking localisation as a shared and largely solved precondition
(Section~\ref{sec:lr-detection}) rather than a contribution. It evaluates on
\emph{LumbarDISC} \cite{richards2026the}; because that collection excludes
studies with hardware, non-degenerative pathology, severe scoliosis or
substantial artifact, conclusions do not extend to those populations, and no
independent external cohort is available to test transportability beyond it. It
compares a \emph{bounded set of architectures}---two convolutional and one
transformer---chosen as representative rather than exhaustive, and does not
evaluate hybrid or vision--language designs
\cite{chen2024symtc,islamsk2026an}. Finally, it reports model performance and
computational cost, not clinical impact: the reader-assistance evidence of
Section~\ref{sec:lr-reader-assistance} \cite{lim2022improved} shows that
workflow benefit requires a prospective reader study, which lies outside the
present scope.

Two constraints inherited from the literature bound the interpretation of any
result. Performance is capped by the reliability of the reference standard
(Section~\ref{sec:lr-label-ceiling}), so figures approaching perfect agreement
should be read as evidence of a protocol artefact rather than of clinical
readiness. And single-benchmark evaluation cannot establish generalisation; the
consistent development-to-validation drop documented in
Section~\ref{sec:lr-generalization} applies to this work as much as to any other.

% ===========================================================================
\section{Chapter Summary}
\label{sec:lr-summary}

This chapter reviewed the literature on automated assessment of lumbar
degenerative disease from MRI, from the clinical grading systems that define the
task through to the architectures that currently attempt it.

The clinical sections established that severity labels originate in
morphological grading systems---Pfirrmann for disc degeneration
\cite{pfirrmann2001magnetic}, Schizas and Lee for central canal
\cite{schizas2010qualitative,lee2011a}, Lee for foraminal
\cite{lee2010a}---whose inter-reader reliability ranges from substantial to poor
depending on the compartment \cite{lurie2008reliability}. This imposes a ceiling
on achievable performance and makes agreement statistics, rather than accuracy,
the appropriate measure. The labels are ordinal, errors concentrate at adjacent
boundaries, and the clinical cost of an error depends on its direction.

The methodological sections traced the field from handcrafted features through
SpineNet's multi-task convolutional grading \cite{jamaludin2017spinenet} to
contemporary multi-stage, multi-view pipelines \cite{batra2025mscan}. The most
robust design principle to emerge is that anatomical localisation should precede
classification \cite{lu2018deepspine,kowlagi2022a}. The most robust empirical
regularity is that detection transfers across institutions while fine-grained
grading does not \cite{zhang2023deep,wu2026external}.

Two questions remain unresolved. Whether attention-based backbones offer any
advantage over well-tuned convolutional ones is contested, with the balance of
lumbar-specific evidence currently favouring CNNs
\cite{kowlagi2022a,lee2026deep,silveira2025automated}, yet no study has compared
them under matched localisation, training and evaluation conditions. And whether
the full three-sequence protocol is necessary, or whether a single ubiquitous
sequence suffices, has been argued from both sides without a controlled test.
These two questions, and their interaction, constitute the gap this thesis
addresses. Chapter~3 sets out the experimental design through which they are
answered.

% ===========================================================================
\printbibliography[heading=bibintoc,title={References}]

\end{document}
 from your main file.
% ===========================================================================

\documentclass[12pt,a4paper]{report}

\usepackage{lmodern}          % scalable fonts; required for microtype expansion
\usepackage[T1]{fontenc}
\usepackage[utf8]{inputenc}
\usepackage[english]{babel}
\usepackage{microtype}
\usepackage[margin=2.5cm]{geometry}
\usepackage{setspace}
\onehalfspacing

\usepackage{amsmath,amssymb}
\usepackage{booktabs}
\usepackage{longtable}
\usepackage{array}
\usepackage{graphicx}
\usepackage[hidelinks]{hyperref}

\usepackage[
  backend=biber,
  style=numeric-comp,
  sorting=none,
  maxbibnames=6,
  minbibnames=3,
  giveninits=true,
  doi=true,
  url=false,
  isbn=false
]{biblatex}

\addbibresource{../lumbar_spine_mri_ai_literature_inventory.bib}
\addbibresource{references-extra.bib}

% The inventory .bib carries internal bookkeeping fields (provenance note, the
% per-paper summary, the local PDF path). These are working metadata, not
% bibliographic data, so they are suppressed from the printed reference list.
% Leaving them in also breaks compilation, since they contain raw underscores.
\AtEveryBibitem{%
  \clearfield{note}%
  \clearfield{comment}%
  \clearfield{annotation}%
  \clearfield{abstract}%
  \clearfield{file}%
}

\begin{document}

\chapter{Literature Review}
\label{ch:litreview}

% ===========================================================================
\section{Introduction}
\label{sec:lr-intro}

\subsection{Scope and Objectives of the Review}
\label{sec:lr-scope}

Lumbar degenerative disease is among the most common indications for magnetic
resonance imaging (MRI) worldwide, and the interpretation of those images is a
repetitive, subjective and time-consuming task that has attracted sustained
attention from the medical image analysis community \cite{lee2024artificial,zhang2024a}.
This chapter reviews that body of work. Its purpose is not to catalogue every
published model, but to establish what is genuinely settled, what remains
contested, and where a controlled experiment can still add knowledge.

The review is organised around a single guiding question: \emph{what determines
whether an automated system can grade lumbar degenerative severity reliably
enough to be clinically useful?} Four candidate answers recur in the
literature---the quality and definition of the reference labels, the degree to
which the anatomy is localised before it is classified, the choice of
representational backbone, and the configuration of imaging sequences supplied
to the model. Each is examined in turn, and the chapter closes by identifying
the two of these that remain inadequately resolved and that this thesis
therefore addresses.

Three specific objectives follow from this. First, to characterise the clinical
grading systems that supply the ground truth for every supervised model in this
field, and to quantify the reliability ceiling those systems impose. Second, to
trace the methodological evolution from handcrafted features to contemporary
convolutional and attention-based architectures, with particular attention to
comparative claims between the two. Third, to assess the evidence on how much
diagnostic information is carried by each MRI sequence and plane, and thus
whether the computational expense of multi-sequence processing is warranted.

\subsection{Search Strategy, Sources, and Inclusion Criteria}
\label{sec:lr-search}

The literature underpinning this chapter was assembled into a structured
inventory of 92 records spanning 2001 to 2026. Sources were identified through
PubMed, Scopus, IEEE~Xplore and arXiv, using combinations of the terms
\emph{lumbar spine}, \emph{MRI}, \emph{spinal stenosis}, \emph{disc
degeneration}, \emph{deep learning}, \emph{convolutional neural network},
\emph{transformer}, \emph{segmentation} and \emph{grading}. Reference lists of
the retrieved systematic reviews \cite{compte2023are,wang2024machine,liawrungrueang2025artificial,mendes2026deep}
were searched by hand for records the database queries had missed, and the
citation graphs of the most influential frameworks were traced forward.

Records were retained if they satisfied at least one of four criteria: they
define or validate a clinical grading standard for lumbar degenerative disease;
they present an automated method for detecting, segmenting or grading lumbar
pathology on MRI; they externally validate such a method; or they describe a
dataset or benchmark on which such methods are trained. Studies confined to
computed tomography or plain radiography were excluded unless they establish a
methodological principle transferable to MRI, as were studies of the cervical
or thoracic spine except where explicitly used for contrast.

Two caveats on the evidence base should be recorded at the outset. First, the
field's publication rate accelerated sharply after 2024, and a substantial
fraction of the most directly relevant work
\cite{batra2025mscan,acharya2026disccentric,nguyen2026crossspine,islamsk2026an}
exists as preprints that have not completed peer review; these are cited as
indicative of methodological direction rather than as settled results. Second,
performance figures reported across studies are frequently not comparable,
because they derive from different datasets, different label definitions and
different evaluation protocols. Wherever numerical results are quoted in this
chapter, the cohort and validation regime that produced them are stated
alongside, and cross-study comparisons are avoided unless the underlying
protocols genuinely align.

\subsection{Structure of the Chapter}
\label{sec:lr-structure}

Sections~\ref{sec:lr-clinical} to~\ref{sec:lr-grading} establish the clinical
substrate: the disease, the imaging protocol used to observe it, and the
grading systems that convert observation into ordinal labels. This material is
not background decoration. The severity classes a model is trained to predict
are defined entirely by these systems, and their documented unreliability
propagates directly into every reported accuracy figure in the chapters that
follow.

Section~\ref{sec:lr-evolution} traces the methodological history of the field,
and Section~\ref{sec:lr-localization} isolates anatomical localisation as the
design principle that most consistently separates effective systems from
ineffective ones. Sections~\ref{sec:lr-backbones} and~\ref{sec:lr-context} then
address the two axes on which this thesis experiments: the representational
backbone, and the spatial and sequential context supplied to it.
Section~\ref{sec:lr-imbalance} treats class imbalance and the ordinal structure
of the label space, and Section~\ref{sec:lr-datasets} the datasets and
benchmarks---principally the RSNA LumbarDISC collection---on which the field now
converges. Sections~\ref{sec:lr-validation} and~\ref{sec:lr-interpretability}
consider validation practice and interpretability respectively.
Section~\ref{sec:lr-gap} synthesises the whole and states the research gap.

% ===========================================================================
\section{Clinical Background: Lumbar Degenerative Disease}
\label{sec:lr-clinical}

\subsection{Epidemiology and Socioeconomic Burden}
\label{sec:lr-epidemiology}

Low back pain is the single largest contributor to years lived with disability
worldwide, and lumbar degenerative change is among its most frequently
implicated structural correlates \cite{compte2023are,zheng2022deep}. Lifetime
prevalence is reported as high as 84\%, and in the United States alone the
combined direct medical and indirect productivity costs are estimated between
\$50 and \$100 billion annually \cite{richards2026the}. Lumbar spinal stenosis
in particular is the most common indication for spinal surgery in patients over
65 \cite{schizas2010qualitative,bharadwaj2023deep}, and demographic ageing is
expected to increase both its incidence and the imaging workload it generates
\cite{lee2024artificial}.

The operational consequence for this thesis is one of volume. MRI referrals for
low back pain have grown faster than the radiology workforce able to report
them \cite{ghosh2014supervised,compte2023are}, and each lumbar study requires
the reader to assign severity at multiple anatomical levels across several
sequences---an average of 14 to 19 minutes of expert time per study
\cite{compte2023are}. It is this combination of high volume, repetitive
structure and documented inter-reader variability that makes the task a
plausible target for automation, rather than any claim that machines perceive
the underlying pathology more acutely than radiologists do.

\subsection{The Degenerative Cascade}
\label{sec:lr-cascade}

Degeneration of the lumbar spine is a progressive, multifactorial process
driven by ageing, mechanical loading, genetic predisposition and lifestyle
factors \cite{esposito2025mribased,zhao2025advances}. It begins in the
intervertebral disc, whose nucleus pulposus loses proteoglycan content and
therefore water from approximately the third decade of life onward. The
resulting dehydration reduces the disc's height, its capacity to distribute
axial load, and its signal intensity on T2-weighted imaging---the last of these
being precisely what the Pfirrmann scale measures \cite{pfirrmann2001magnetic}.

Disc height loss initiates a compensatory cascade through the rest of the motion
segment. The annulus fibrosus weakens and may permit bulging or frank herniation
of nuclear material; the ligamentum flavum buckles and hypertrophies; the
posterior facet joints, now bearing altered load, develop osteoarthritic change
and osteophytosis \cite{bharadwaj2023deep,nikpasand2024automated}. Changes in
the adjacent vertebral endplates and marrow---the Modic
phenotypes~\cite{modic1988degenerative}---frequently accompany this process and
have themselves been the target of automated detection
\cite{gao2022automatic,athertya2021classification}. The clinical significance of
the cascade is that these individually modest changes act cumulatively to
encroach on the spaces occupied by neural tissue.

\subsection{The Three Anatomical Forms of Stenosis}
\label{sec:lr-stenosis-types}

Encroachment is conventionally partitioned by anatomical compartment, and this
partition is directly mirrored in the label structure of the dataset used in
this thesis \cite{richards2026the}:

\begin{enumerate}
  \item \textbf{Spinal canal stenosis} is narrowing of the central canal, which
  contains the thecal sac and the cauda equina. It is best appreciated on axial
  T2-weighted images, where the relationship between cerebrospinal fluid and
  nerve rootlets is visible \cite{schizas2010qualitative,lee2011a}. Severe
  central stenosis characteristically produces neurogenic claudication.

  \item \textbf{Neural foraminal narrowing} is constriction of the lateral
  foramina through which the exiting nerve roots pass. It is assessed on
  parasagittal images, where the epidural fat surrounding the root provides the
  contrast that makes compression visible---T1 weighting being particularly
  suited to this \cite{lee2010a}. It produces dermatomal radicular symptoms.

  \item \textbf{Subarticular stenosis}, also termed lateral recess stenosis,
  affects the zone between the medial edge of the facet and the foramen, where
  it compresses the traversing root before it exits at the level below
  \cite{lurie2008reliability,andreisek2014consensus}.
\end{enumerate}

The distinction matters computationally as well as clinically, because the three
compartments are not equally visible in any single plane. Foraminal narrowing is
poorly assessed from axial images and subarticular stenosis poorly from
sagittal ones, a constraint that Section~\ref{sec:lr-context} shows to have
concrete consequences for model design. It is also the compartment that
correlates with which sequence a model must be given, and thus sits at the
centre of the multi-sequence question this thesis examines.

\subsection{Clinical Presentation and the Treatment Decision}
\label{sec:lr-treatment}

Severity grading is not an academic exercise; it partitions patients into
distinct management pathways. Conservative management---physiotherapy, analgesia,
epidural steroid injection---is appropriate for mild and most moderate disease,
whereas severe compression with corresponding symptoms may warrant decompressive
surgery or fusion \cite{schizas2010qualitative,mysliwiec2010msu}. The evidence
that morphological grade carries genuine prognostic weight is strongest for the
Schizas classification, where grade~C and~D patients were substantially more
likely to fail conservative treatment \cite{schizas2010qualitative}, and for the
MSU herniation grades, where size~2 and~3 lesions identify microdiscectomy
candidates \cite{mysliwiec2010msu}.

Two qualifications temper this. First, the correlation between imaging findings
and symptoms is imperfect: anatomical stenosis is demonstrable in a substantial
proportion of asymptomatic individuals \cite{lurie2008reliability,esposito2025mribased},
so imaging severity constrains rather than determines management. Second, and
more consequentially for automated systems, the clinically decisive boundary is
not evenly distributed across the ordinal scale. The distinction between
\emph{normal} and \emph{mild} rarely changes management, whereas the distinction
between \emph{moderate} and \emph{severe} frequently does. A model that
distributes its errors uniformly across grade boundaries is therefore not
clinically equivalent to one of the same overall accuracy that concentrates its
errors at the harmless end---a point developed in
Section~\ref{sec:lr-error-asymmetry}.

% ===========================================================================
\section{MRI Acquisition: Sequences, Planes, and Diagnostic Yield}
\label{sec:lr-mri}

\subsection{Why MRI is the Reference Modality}
\label{sec:lr-why-mri}

MRI is the reference standard for non-invasive assessment of the lumbar spine.
Computed tomography and plain radiography depict osseous anatomy well but
resolve soft tissue poorly, and both expose the patient to ionising radiation.
MRI resolves the intervertebral disc, thecal sac, individual nerve roots and
ligamentous structures with a contrast unavailable to the alternatives
\cite{zhao2025advances,liawrungrueang2025artificial}, and it is the disposition
of exactly these soft-tissue structures that defines every grading system
reviewed in Section~\ref{sec:lr-grading}. MRI has correspondingly displaced CT
myelography for stenosis assessment \cite{tumko2024a}.

\subsection{Complementary Information Across Sequences and Planes}
\label{sec:lr-sequences}

Lumbar protocols are inherently multiplanar and multisequence, because the spine
is a three-dimensional structure whose pathologies are not all visible in one
projection. Three acquisitions dominate clinical practice and constitute the
input available in the dataset used here \cite{richards2026the}.

\emph{Sagittal T2-weighted} imaging renders fluid hyperintense, so cerebrospinal
fluid in the thecal sac appears bright against darker ligament and bone. It is
the sequence on which disc hydration---and therefore Pfirrmann
grade---is judged \cite{pfirrmann2001magnetic}, and on which the
craniocaudal extent of canal compromise is surveyed. Short tau inversion
recovery (STIR) suppresses the signal from fat and thereby exposes marrow
oedema and inflammatory change, making it the sequence of choice for active
Modic type~1 change \cite{gao2022automatic}.

\emph{Sagittal T1-weighted} imaging renders fat hyperintense. Because the exiting
nerve root within the neural foramen is normally surrounded by epidural fat, T1
provides the contrast against which obliteration of that fat---the basis of the
Lee foraminal grading system---becomes visible \cite{lee2010a}.

\emph{Axial T2-weighted} imaging provides the transverse cross-section through
the disc level. It is the only view in which the morphology of the dural sac,
the arrangement of the cauda equina rootlets and the lateral recesses can be
directly assessed, and it therefore underpins both the Schizas and Lee central
canal systems \cite{schizas2010qualitative,lee2011a} and the assessment of
subarticular stenosis \cite{andreisek2014consensus}.

\subsection{Matching Pathology to Sequence}
\label{sec:lr-sequence-matching}

The correspondence between pathology and sequence is not a matter of
convention but of physics, and the empirical literature bears this out with
unusual clarity. The most direct evidence comes from external validation of
sagittal-only systems. Nigru et al.\ evaluated SpineNetV2 across eleven
radiological features on 1{,}747 lumbosacral discs and found agreement above
$\kappa = 0.7$ for most targets but distinctly lower agreement for foraminal
stenosis and disc herniation, attributing the shortfall specifically to the
absence of transverse anatomical context in sagittal input
\cite{nigru2024external}. Hallinan et al.\ approached the same constraint from
the opposite direction, deliberately assigning axial T2 to central canal and
lateral recess grading while reserving sagittal T1 for the neural foramina
\cite{hallinan2021deep}.

This has a direct bearing on the experiments in this thesis. If a
single-sequence model is to be evaluated, the choice of which sequence to
retain is not arbitrary: sagittal T2 is the most defensible candidate because it
is acquired in essentially every back-pain protocol and because it carries the
signal for canal stenosis and disc degeneration, but the literature predicts
that it should be weakest precisely on the foraminal and subarticular targets.
Whether a sufficiently expressive architecture can compensate for that
predicted weakness is an open question, and one this thesis is designed to
answer.

\subsection{Sources of Acquisition Variability}
\label{sec:lr-variability}

Models trained on one institution's images frequently degrade on another's, and
the proximate causes lie in acquisition. Scanner manufacturer and field strength
(typically 1.5 versus 3.0~T), slice thickness and spacing, in-plane resolution,
echo and repetition times, the presence or absence of fat suppression, and
patient positioning all vary between sites \cite{senzgamboa2023automatic,mendes2026deep}.
Axial acquisitions vary further in whether slices form a continuous stack or are
angled to lie parallel to each disc level---a difference the LumbarDISC
inclusion criteria explicitly accommodate \cite{richards2026the}, and one that
any model consuming axial data must tolerate.

S\'aenz~Gamboa et al.\ treat this variability as the research problem in its own
right rather than as a nuisance, evaluating segmentation across multiple centres
and acquisition parameters and finding that ensembles of U-Net variants were
required to maintain performance \cite{senzgamboa2023automatic}. Batra et al.\
make the complementary design argument: rather than assuming uniform histograms
and a fixed slice count per study, their pipeline is built to accept a variable
number of axial and sagittal slices, on the grounds that this reflects the
conditions under which clinical data actually arrives \cite{batra2025mscan}.
Variability is thus not merely a threat to validity but a specification
constraint, and Section~\ref{sec:lr-generalization} returns to how effectively
the field has met it.

% ===========================================================================
\section{Clinical Grading Systems as the Origin of Ground Truth}
\label{sec:lr-grading}

Every supervised model reviewed in this chapter learns to reproduce labels
generated by a human applying one of the systems described below. The
properties of those systems---their definitions, their granularity and above all
their reliability---therefore bound what any model trained on them can achieve.
This section treats them as the measurement instruments they are.

\subsection{Disc Degeneration: Pfirrmann and its Modified Extension}
\label{sec:lr-pfirrmann}

The Pfirrmann classification, introduced in 2001, remains the most widely
adopted scheme for intervertebral disc degeneration and is applied to sagittal
T2-weighted images \cite{pfirrmann2001magnetic}. It assigns one of five ordinal
grades on the basis of four morphological criteria: homogeneity of the nucleus
pulposus, distinction between nucleus and annulus, T2 signal intensity as a
proxy for hydration, and disc height. Grade~I denotes a homogeneous, bright,
normally proportioned disc; grade~V a collapsed disc of uniformly low signal.

Its virtues are simplicity and universality; its limitations are well
documented. It is a subjective visual judgement, it is insensitive to change
over short intervals, and it captures only intrinsic disc state, saying nothing
about Modic change, high-intensity zones or the downstream compression of neural
structures \cite{esposito2025mribased,zhao2025advances}. It also suffers a
ceiling effect in elderly cohorts, where most discs present as grade~IV or~V and
clinically meaningful variation is compressed into the top of the scale.
Griffith et al.\ addressed this with an eight-level modified scale supported by a
24-image reference panel, validated on 260 discs in 52 subjects of mean age 73;
intra-observer agreement was excellent (weighted $\kappa$ 0.79--0.91) and
inter-observer agreement substantial (0.65--0.67), with most disagreement
confined to a single grade \cite{griffith2007modified}.

Esposito et al.\ place these two systems in a wider context that is worth
recording: a scoping review of 569 studies identified \emph{93} distinct grading
systems for lumbar disc degeneration---63 subjective, 25 quantitative---of which
the Pfirrmann method accounted for just over half of all reported use
\cite{esposito2025mribased}. The apparent standardisation of the field is
therefore partly illusory, and comparisons of automated grading accuracy across
studies must confirm that the same instrument was being reproduced.

\subsection{Central Canal Stenosis: Schizas and Lee}
\label{sec:lr-central-grading}

Two morphological systems have largely displaced simple dural sac
cross-sectional area (DSCA) thresholds for central canal assessment, on the
grounds that area measurements correlate poorly with symptoms and overlap
substantially between symptomatic and asymptomatic individuals
\cite{schizas2010qualitative}.

The \textbf{Schizas classification} grades the morphology of the dural sac on
axial T2 images by the ratio of cerebrospinal fluid to nerve rootlets. Grades~A
and~B retain visible CSF---in~A the rootlets occupy part of the sac, in~B they
fill it but remain individually discernible, giving a characteristic grainy
appearance. Grades~C and~D show no discernible rootlets and no CSF signal,
distinguished from one another by whether posterior epidural fat is preserved
(C) or obliterated (D) \cite{schizas2010qualitative}. Applied to 95 subjects
across surgical, conservatively managed and low-back-pain groups, intra- and
inter-observer agreement were substantial and moderate respectively
($\kappa = 0.65$ and $0.44$). Critically, area measurement over-diagnosed
stenosis in 35 patients and under-diagnosed it in 12 relative to the
morphological grade, and grades~C and~D predicted failure of conservative
management \cite{schizas2010qualitative}.

The \textbf{Lee central canal system} offers a simpler four-grade alternative
based on the degree of separation of the cauda equina rootlets on axial T2:
grade~0 no stenosis; grade~1 obliteration of the anterior CSF space with all
rootlets still separated; grade~2 some rootlets aggregated; grade~3 none
separated \cite{lee2011a}. Assessed by four musculoskeletal radiologists across
81 patients, it achieved inter-reader ICC of 0.730--0.953 and intra-reader
$\kappa$ of 0.863--0.900, and grades tracked cross-sectional area and
anteroposterior diameter at nearly all levels---notably better agreement than
the Schizas system, at the cost of finer morphological detail.

\subsection{Neural Foraminal Stenosis: Lee and Sartoretti}
\label{sec:lr-foraminal-grading}

The \textbf{Lee foraminal system} grades parasagittal images by the pattern of
perineural fat obliteration and the morphology of the nerve root: grade~0
normal; grade~1 fat obliteration in two opposing directions; grade~2
obliteration in all four directions without morphological change to the root;
grade~3 nerve root collapse or distortion \cite{lee2010a}. Evaluated on 576
foramina in 96 patients from L3--L4 to L5--S1 by two radiologists, inter- and
intra-observer agreement were near-perfect ($\kappa$ approximately 0.80--1.0).
This is the direct clinical precursor of the left and right foraminal narrowing
targets in the benchmark used in this thesis. The Sartoretti classification
offers a six-point sagittal alternative (grades~A--F) indexed to progressive
contact between the exiting root and each foraminal boundary.

The near-perfect agreement Lee et al.\ report deserves scepticism when
generalised. It was obtained by two experienced readers applying explicit
criteria to a curated series; the multi-reader, multi-institution agreement
figures in Section~\ref{sec:lr-reliability} are considerably lower, and it is
the latter that describe the conditions under which large training corpora are
actually annotated.

\subsection{Subarticular and Lateral Recess Stenosis}
\label{sec:lr-subarticular-grading}

Subarticular stenosis is graded on axial T2 by the proportion of the lateral
recess compromised: mild at up to one third, moderate between one and two
thirds, severe beyond two thirds, with nerve root involvement described as
touching, displacing or compressing \cite{lurie2008reliability,andreisek2014consensus}.
Compared with the disc and central canal systems this is a notably less
crisply specified instrument, resting on a visually estimated ratio rather than
a discrete morphological pattern. That imprecision shows in the reliability
data: subarticular stenosis is consistently the least reproducible of the
targets \cite{lurie2008reliability}, and it is also the target on which the
severity distribution in LumbarDISC is least skewed
(Section~\ref{sec:lr-lumbardisc}). Its difficulty is thus intrinsic, not merely
a consequence of scarce severe examples.

\subsection{Disc Herniation: Nomenclature and the MSU Classification}
\label{sec:lr-herniation-grading}

Herniation is described qualitatively by the standardised nomenclature of
bulge, protrusion, extrusion and sequestration \cite{fardon2014lumbar}, and
quantitatively by the Michigan State University (MSU) classification, which
grades size against a transverse line joining the medial edges of the facet
joints on the axial slice of maximal herniation \cite{mysliwiec2010msu}.
Grade~1 extends less than half the distance to that line, grade~2 more than
half without crossing it, grade~3 completely beyond it. The MSU system supplies
the reference standard for several automated herniation studies
\cite{zhang2023deep}, and its size~1 versus size~2/3 boundary maps onto the
conservative versus surgical decision.

\subsection{Inter- and Intra-Observer Reliability}
\label{sec:lr-reliability}

The most consequential study for present purposes is that of Lurie et al., who
assessed standardised MRI interpretation of lumbar stenosis using 58 randomly
selected studies from the SPORT trial read by four independent readers, with 20
re-read after at least a month \cite{lurie2008reliability}. Inter-reader
agreement was substantial for central stenosis ($\kappa = 0.73$) but only
moderate for foraminal stenosis ($0.58$) and nerve root impingement ($0.51$),
and poorest for subarticular stenosis ($0.49$), with marked variation between
reader pairs. Quantitative measurement fared no better in absolute terms: mean
absolute inter-reader difference in thecal sac area was 128~mm$^2$, or 13\%.
Intra-reader agreement exceeded inter-reader agreement for every feature.

Two structural findings recur across the wider reliability literature. First,
agreement is systematically worse for the lateral compartments than for the
central canal---the same ordering that later appears in automated model
performance. Second, disagreement is overwhelmingly confined to adjacent
grades: Griffith et al.\ report 84\% of intra-observer and 91\% of
inter-observer disagreement spanning a single grade
\cite{griffith2007modified}, and the same pattern recurs in automated
evaluation \cite{mcsweeney2023external,udomluck2026multitask}. Earlier consensus
efforts confirm that the underlying definitions were themselves contested: the
Delphi survey of Mamisch et al.\ \cite{mamisch2012radiologic} and the consensus
conference reported by Andreisek et al.\ \cite{andreisek2014consensus} were both
convened precisely because radiologic criteria for stenosis varied between
institutions.

\subsection{Label Noise as a Performance Ceiling}
\label{sec:lr-label-ceiling}

The reliability figures above are not a peripheral caveat; they are a hard
constraint on the interpretation of every result in this field. If four expert
readers agree on subarticular severity at $\kappa = 0.49$, then a model trained
on one reader's labels and evaluated against another's cannot exceed that
agreement, and a reported accuracy approaching unity on such a target should
prompt scrutiny of the evaluation protocol rather than admiration. The
LumbarDISC dataset description states the point explicitly, noting that
published inter-rater agreement on these grades ranges from poor to substantial
depending on anatomical location \cite{richards2026the}.

The field has responded to this in two ways. The first is to make the reference
standard more robust by construction---majority voting across a panel with
adjudication of disagreements, as in Tumko et al.\ (seven radiologists)
\cite{tumko2024a} and Su et al.\ (two clinicians with expert adjudication)
\cite{su2022automatic}. The second is to reframe the evaluation question: rather
than asking whether a model matches a single reader, asking whether its
disagreement with readers falls inside the band of disagreement between readers.
Won et al.\ furnish the cleanest example, reporting that the difference between
their two experts was statistically significant while the difference between
each expert and the model trained on that expert's labels was not
\cite{won2020spinal}. Comparatively few studies adopt either safeguard, and
Section~\ref{sec:lr-reporting} returns to this as a systemic weakness.

\subsection{The Ordinal Structure of the Label Space}
\label{sec:lr-ordinality}

Severity labels are ordinal, not categorical: \emph{normal}, \emph{mild},
\emph{moderate} and \emph{severe} stand in a fixed order, and the distance
between them is neither uniform nor clinically symmetric. Treating them as
unordered classes under categorical cross-entropy---as the majority of the
reviewed work does---discards that structure, penalising a
severe-graded-as-normal error exactly as heavily as a moderate-graded-as-severe
one.

Two empirical observations follow. First, errors concentrate overwhelmingly at
adjacent grade boundaries. McSweeney et al.\ report that SpineNet disagreed with
human raters on 20.83\% of discs but that only 0.85\% of discs differed by more
than one grade \cite{mcsweeney2023external}; Udomluck et al.\ observe the same
concentration \cite{udomluck2026multitask}; Ishimoto et al.\ find exact-grade
agreement of 65.7\% rising to 94.1\% when collapsed to a severe-versus-rest
dichotomy \cite{ishimoto2020could}. Much of the apparent error in ordinal
grading is therefore boundary error rather than gross misclassification.

Second, and less expectedly, explicitly ordinal objectives have not reliably
improved on this. Niemeyer et al.\ tested soft-kappa loss, ordinal
cross-entropy and regression losses against plain cross-entropy for Pfirrmann
grading on 1{,}599 patients and 7{,}948 discs, together with class pooling and
class-weighted losses, and found that none improved overall classification
performance; the default cross-entropy classifier reached $\kappa = 0.92$ with
predictions within one grade of ground truth in 99.2\% of cases
\cite{niemeyer2021a}. Patil et al.\ likewise found performance strongest at the
extremes of the Pfirrmann scale and on the binary low-versus-high distinction,
with intermediate grades remaining the difficulty \cite{patil2026mribased}. The
ordinal structure of the problem is thus real and measurable, but the question
of how to exploit it is genuinely open---a point that informs the loss functions
adopted in Chapter~3.

% ===========================================================================
\section{Evolution of Automated Lumbar MRI Analysis}
\label{sec:lr-evolution}

\subsection{Handcrafted Features and Classical Machine Learning}
\label{sec:lr-classical}

Before deep learning, automated analysis of lumbar MRI depended on features
designed by hand. Investigators specified mathematical descriptors of local
appearance and passed them to a conventional classifier. Ghosh et al.\ combined
Histogram of Oriented Gradients (HOG) descriptors with a support vector machine
and a set of geometric heuristics, reporting 99\% disc localisation accuracy
across 53 clinical cases and 318 discs, and---unusually for the period---emitting
a tight bounding box for each disc rather than a single interior point, so that
a downstream system could operate on the region directly
\cite{ghosh2012a}. Athertya and Kumar took a comparable approach to vertebral
degeneration, extracting local binary patterns, Hu moments and gray-level
co-occurrence matrix (GLCM) statistics to classify Modic changes, endplate
defects and focal changes, reaching 85.91\% accuracy for the extent of Modic
change on 100 patients \cite{athertya2021classification}. Related efforts applied
hybrid classical models to disc abnormality detection
\cite{unal2015detection,castromateos2016intervertebral} and semi-automatic
computer-aided classification of degenerative disease \cite{ruizespaa2015semiautomatic},
while Huber et al.\ pursued texture analysis with decision trees as a
quantitative alternative to visual stenosis grading \cite{huber2019qualitative}.
He et al.\ applied a synchronised superpixel representation to define consistent
regions of interest across levels before classifying neural foraminal stenosis,
an early attack on one of the exact conditions graded in this thesis
\cite{he2016automated}.

These systems established feasibility but generalised poorly. Handcrafted
descriptors are sensitive to scanner hardware, field strength, slice thickness
and patient positioning, and the anatomical variability of the lumbar spine
defeats fixed geometric assumptions. One finding from this era nevertheless
retains value: Athertya and Kumar's feature-sensitivity analysis showed that
GLCM entropy alone sufficed to detect focal changes, while endplate-defect
classification depended almost entirely on the first two Hu moments
\cite{athertya2021classification}. Such explicit attribution of predictive power
to individual features is something the deep models that displaced these methods
recovered only later, and imperfectly, through the interpretability techniques of
Section~\ref{sec:lr-interpretability}.

\subsection{The First End-to-End Deep Learning Systems}
\label{sec:lr-first-dl}

Convolutional neural networks removed the need for explicit feature design by
learning hierarchical representations directly from pixels \cite{he2016deep}. In
lumbar imaging the decisive contribution was SpineNet
\cite{jamaludin2017spinenet}. Trained on 2{,}009 patients and 12{,}018 disc
volumes from sagittal T2 images, it predicted six radiological scores
simultaneously---Pfirrmann grade, disc narrowing, upper and lower endplate
defects, and upper and lower marrow changes---from a shared convolutional trunk
that branched into task-specific heads. Two design decisions proved durable.
First, multi-task learning let one architecture serve several targets, reducing
the free parameters relative to training separate networks and exploiting the
correlation between co-occurring pathologies. Second, the system required only
disc-level class labels: no voxel-wise segmentation and no slice-level
localisation, with the evidence hotspots emerging implicitly from classification
training. Reported performance was state-of-the-art and near-human across all
six scores.

The same group demonstrated shortly afterwards that longitudinal clinical scans
could serve as free supervision. A Siamese network was pre-trained with a
contrastive loss on whether two scans belonged to the same patient---information
available from DICOM metadata without knowing the patient's identity---together
with an auxiliary vertebral-level prediction task, then fine-tuned for
degeneration grading. On 1{,}016 subjects, 423 with follow-up imaging, the
pre-trained network outperformed training from scratch and reached equivalent
performance from far fewer annotated examples \cite{jamaludin2017selfsupervised}.
This anticipated by several years the self-supervised pretraining strategies now
routine in medical imaging, and it is the direct ancestor of the contrastive
approach discussed in Section~\ref{sec:lr-contrastive}.

SpineNet's principal limitation was structural rather than architectural: it
consumed preprocessed sagittal disc volumes, and could therefore not assess
pathologies that require the transverse plane. SpineNetV2 extended coverage to
eleven radiological features and moved to a multi-stage design, using a U-Net to
detect and label vertebral bodies and discs before a ResNet classifier graded
each \cite{windsor2022spinenetv2}. The reliance on sagittal input nonetheless
persisted, and Section~\ref{sec:lr-generalization} shows this to be the
consistent locus of its external-validation failures.

\subsection{From Binary Detection to Multi-Level, Multi-Pathology Grading}
\label{sec:lr-task-evolution}

The task definition itself has evolved, and the trajectory matters more than any
individual model. Early deep learning work asked binary questions: is stenosis
present or absent \cite{altun2023lssvgg16}; is a herniation visible in this image
\cite{tsai2021lumbar}. Such formulations are tractable but clinically thin,
because management depends on severity and on which level is affected.

Three shifts followed. The first was to \emph{multi-class severity}: Won et al.\
graded stenosis status on 542 L4--5 axial images against two independent experts
\cite{won2020spinal}; Niemeyer et al.\ predicted the full five-point Pfirrmann
scale across 7{,}948 discs \cite{niemeyer2021a}. The second was to
\emph{multi-pathology} output, in which one model reports several findings at
once: Lehnen et al.\ evaluated disc herniation, bulging, canal stenosis,
nerve-root compression and spondylolisthesis across 888 lumbar segments from a
single CNN \cite{lehnen2021detection}; Su et al.\ graded herniation, central
canal stenosis and nerve-root compression from a shared ResNet-50 trunk with
three classification heads \cite{su2022automatic}. The third was to
\emph{per-level, per-side} prediction, in which the unit of analysis is not the
study or the image but the individual disc level and, for the lateral
compartments, the individual side. Hallinan et al.\ exemplify the mature form,
grading central canal, lateral recess and neural foraminal stenosis at each
level with the sequence appropriate to each \cite{hallinan2021deep}. Single-stage
detectors have since been adapted to deliver the same output more cheaply: Wang
et al.\ modified a YOLOv5 network to detect regions of interest and grade
central, lateral recess and foraminal stenosis simultaneously
\cite{wang2025a}, and a related design grades Pfirrmann degeneration, herniation
and high-intensity zones from sagittal and axial input in one pass
\cite{wang2025automated}. Minin et al.\ pursued the same objective for Pfirrmann
grading with an explicitly open model, motivated by the inaccessibility of
proprietary graders \cite{minin2025spinescan}.

This third formulation is precisely that of the benchmark used in this thesis,
which requires twenty-five graded outputs per study: five conditions across five
disc levels \cite{rsna2024rsna,richards2026the}. The progression is therefore not
merely historical. It explains why performance figures from the binary era
cannot be compared with contemporary results, and why a model that classifies a
whole image is solving a different and substantially easier problem than one
that must attribute severity to a specific anatomical location.

\subsection{Chronological Synthesis}
\label{sec:lr-chronology}

Four phases are discernible. Until roughly 2016 the field was dominated by
handcrafted features and classical classifiers, with localisation as the central
problem \cite{ghosh2012a,ghosh2014supervised}. From 2017 to 2020, CNNs displaced
these methods and multi-task grading became feasible
\cite{jamaludin2017spinenet,lu2018deepspine,han2018spinegan}. From 2021 to 2023
the emphasis moved to clinical credibility: external validation
\cite{mcsweeney2023external}, interpretability \cite{bharadwaj2023deep,gao2022automatic},
reader-assistance evaluation \cite{lim2022improved} and multi-centre testing
\cite{yi2023deep,su2022automatic}. From 2024 onward, the release of a large
public benchmark \cite{richards2026the} redirected effort toward anatomically
localised, multi-view severity grading, alongside the arrival of transformer and
vision-language architectures
\cite{batra2025mscan,nguyen2026crossspine,islamsk2026an,zeng2025gpt4lfs}.

The through-line is a steady increase in the specificity of what is being
predicted, accompanied---until recently---by a decrease in the rigour with which
it is validated. Both systematic reviews of the period make the same complaint:
Compte et al.\ found that algorithms performed markedly worse in replication or
external validation than in their development studies \cite{compte2023are}, and
Wang et al.\ concluded that despite pooled sensitivity of 0.84 and specificity of
0.87 across twelve studies and 15{,}044 patients, no model was reliable enough
for clinical practice \cite{wang2024machine}. Section~\ref{sec:lr-validation}
takes up this tension directly.

% ===========================================================================
\section{Anatomical Localization as a Design Principle}
\label{sec:lr-localization}

\subsection{Disc and Vertebra Detection, Labelling, and Level Assignment}
\label{sec:lr-detection}

Before severity can be attributed to L4--L5, the system must know which structure
is L4--L5. Level assignment is thus a precondition of per-level grading, and it
has been treated as a problem in its own right since the earliest work
\cite{ghosh2012a,lootus2015automated}. Pan et al.\ illustrate the mature
solution: a Faster R-CNN locates vertebral bodies L1--S1 on sagittal images, the
geometric relationship between sagittal and axial acquisitions identifies the
corresponding disc in each axial slice, and the disc region of interest is then
extracted---each of the three stages reported at 100\% accuracy, with the residual
difficulty falling entirely on the subsequent classification
\cite{pan2021automatically}. Detection has, in short, largely been solved;
grading has not.

\subsection{Semantic and Instance Segmentation of Lumbar Structures}
\label{sec:lr-segmentation}

Dense segmentation provides a stronger anatomical substrate than bounding boxes.
The U-Net architecture \cite{ronneberger2015unet} and its descendants dominate,
as a dedicated systematic review of deep-learning-assisted disc segmentation
confirms \cite{wang2024deep}.
Ghosh and Chaudhary segmented vertebrae, disc, dural sac and background jointly
using cascaded random forests over image and sparsely sampled context features,
achieving Dice 0.87 and 0.84 for dural sac and disc across 212 clinical cases
\cite{ghosh2014supervised}. Han et al.\ introduced adversarial training for
multi-structure segmentation \cite{han2018spinegan}, S\'aenz-Gamboa et al.\
established CNN semantic segmentation of the structural elements surrounding the
spinal cord \cite{senzgamboa2021automatic}, and Al-Kafri et al.\ applied
SegNet to boundary delineation for stenosis assessment across 515 symptomatic
patients, contributing two purpose-built metrics---confidence and
consistency---for assessing the quality of the ground truth itself rather than
only the model \cite{alkafri2019boundary}.

Subsequent work has largely consisted of architectural refinement. Wang et al.\
added multilevel feature-map attention and residual blocks to an Attention U-Net
\cite{wang2022automatic}; Ahmed et al.\ targeted class imbalance between
anatomical structures with a custom combined loss and modified upsampling
\cite{ahmed2024pioneering}; Silveira et al.\ added an Inception module for
multi-scale extraction and a dual-output mechanism, reporting mIoU 0.8974 and
F1 0.9444 on SPIDER and---notably---outperforming TransUNet
\cite{silveira2025automated}. M\"oller et al.\ extended coverage to the whole
spine with SPINEPS, segmenting fourteen structures including the posterior
elements, and established a useful general result: segmenting semantically first
and instance-wise second outperformed training directly on instance segmentation,
beating an nnU-Net baseline \cite{isensee2021nnunet} on every structure and metric
\cite{mller2025spinepsautomatic}. Chen et al.\ merged convolutional and attention
pathways in parallel with a novel position embedding, outperforming fifteen
competing models \cite{chen2024symtc}.

\subsection{Public Segmentation Resources and Benchmarks}
\label{sec:lr-segmentation-data}

Progress in this subfield has been unusually well served by shared data. The
SPIDER dataset provides 447 sagittal T1 and T2 series from 218 patients across
four hospitals, with reference segmentations of vertebrae, discs and spinal
canal plus per-level radiological gradings, together with reference performance
for a baseline algorithm and nnU-Net and a continuous public challenge on a
sequestered test split \cite{vandergraaf2024lumbar}. Several of the segmentation
models above are trained and benchmarked on it
\cite{ahmed2024pioneering,silveira2025automated,patil2026mribased}, which makes
their results mutually comparable in a way that is rare elsewhere in this
literature.

\subsection{Quantitative Morphometry}
\label{sec:lr-morphometry}

Segmentation enables measurement, and measurement offers an escape from ordinal
subjectivity. The dural sac cross-sectional area is the canonical example.
Ghobrial and Roth compared U-Net, Attention U-Net and MultiResUNet for automated
DSCA quantification on 683 scans with 50 held back for external validation;
MultiResUNet performed best, with Pearson correlation 0.9917 and mean absolute
error 23.70~mm$^2$, holding up externally at 99.95\% accuracy and F1 0.9393
\cite{ghobrial2025deep}. Zheng et al.\ pursued the same objective for disc
degeneration, segmenting degeneration-related regions with a Swin-transformer
network and computing signal-intensity and geometric features, then proposing a
quantitative criterion to replace subjective Pfirrmann grading on the basis of
strong correlation with grade across a large population \cite{zheng2022deep}.
Related work measures the lordosis angle from segmentation masks
\cite{ji2026mdvmunet}.

A caution is warranted. Ghobrial and Roth's model relies on T1 weighting alone
and a modest dataset, as the authors acknowledge \cite{ghobrial2025deep}, and
Schizas' original finding---that area measurement over- and under-diagnosed
stenosis relative to morphological grade in a substantial minority of patients
\cite{schizas2010qualitative}---applies equally whether the area is traced by
hand or by network. Automating a measurement does not repair a weak correlation
between that measurement and the clinical state.

\subsection{Evidence that Localisation Precedes Classification}
\label{sec:lr-localize-first}

The strongest recurring methodological finding in this literature is that
explicit anatomical localisation before classification improves grading. The
argument was made structurally by DeepSPINE, which chained U-Net vertebral
segmentation and disc-level designation to a multi-task grading network across
4{,}075 patients and more than 22{,}000 disc-level annotations
\cite{lu2018deepspine}, and it has been corroborated repeatedly since.

Four lines of evidence converge. Bharadwaj et al.\ segmented dural sac and disc
with V-Net, localised facet and foramen by geometric rule, and only then
classified, reaching Dice 0.93 and 0.94 with foramen and facet localisation IoU
0.72 and 0.83 \cite{bharadwaj2023deep}. \v{S}u\v{s}ter\v{s}i\v{c} et al.\ showed
that segmenting, cropping and enhancing the region of interest before
classification produced 0.87 and 0.91 accuracy on axial and sagittal views for a
multi-class herniation problem \cite{sustersic2022a}. Kowlagi et al.\ made the
comparative case most directly, arguing that a well-tuned three-stage pipeline of
segmentation, localisation and classification allows convolutional networks to
outperform transformer approaches on Pfirrmann grading in a population cohort
\cite{kowlagi2022a}. Acharya and Kansakar tested the principle in its strongest
form, hypothesising that anatomical focus rather than architecture is the binding
constraint, and grading severity from a single localised region of interest per
disc \cite{acharya2026disccentric}.

The implication for the present work is direct. If localisation dominates
architecture in determining performance, then a comparison between backbones is
only meaningful when both receive identically localised input---otherwise the
comparison measures preprocessing rather than representation. This is a
methodological requirement that Section~\ref{sec:lr-cnn-vs-transformer} shows
much of the existing comparative literature fails to satisfy.

\subsection{Reduced-Annotation and Unsupervised Alternatives}
\label{sec:lr-weak-supervision}

Dense annotation is expensive, and several strands of work reduce the
requirement. Kuang et al.\ eliminated manual labels entirely for vertebral
segmentation: sub-optimal masks from a rule-based method were combined through a
voting mechanism to supervise CNN training, with a regional strategy restricting
attention to a specific vertebral region, validated across 243 volunteers
\cite{kuang2020mrisegflow}. SpineNet's evidence hotspots represent the
complementary weakly supervised route, obtaining localisation from
classification labels alone \cite{jamaludin2017spinenet}, while the SPIDER
annotation protocol adopted a human-in-the-loop compromise, training on a small
subset to produce initial masks that were then corrected and fed back
\cite{vandergraaf2024lumbar}. Self-supervised pretraining on unlabelled or
metadata-labelled data \cite{jamaludin2017selfsupervised} completes the picture.
Together these establish that the annotation burden, long cited as the principal
obstacle to clinical deployment, is substantially negotiable.

% ===========================================================================
\section{Backbone Architectures for Severity Classification}
\label{sec:lr-backbones}

\subsection{Convolutional Networks}
\label{sec:lr-cnns}

Convolutional networks exploit two properties well matched to imaging: local
connectivity, since a filter responds to a neighbourhood rather than the whole
field; and translation equivariance, since the same filter is applied at every
position. Stacking convolutions builds a hierarchy in which early layers respond
to edges and texture and later layers to composite anatomical structure.

\textbf{ResNet} addressed the degradation encountered when networks are made
deeper, introducing residual blocks whose skip connections allow a layer's input
to bypass subsequent convolutions and be added to their output, providing an
unimpeded path for gradients and permitting stable training at fifty layers and
beyond \cite{he2016deep}. In lumbar imaging ResNet is ubiquitous: as the
classification head in multi-task grading \cite{su2022automatic}, as the
detection backbone in staged pipelines \cite{zhang2023deep}, as the encoder in
segmentation networks \cite{batra2025mscan}, and as the classifier in
SpineNetV2 \cite{windsor2022spinenetv2}.

\textbf{EfficientNet} attacked a different problem---how to allocate a fixed
computational budget---by scaling depth, width and input resolution together
under fixed coefficients derived from a neural architecture search, rather than
scaling one dimension in isolation \cite{tan2019efficientnet}. The resulting
models reach comparable accuracy at substantially lower parameter and FLOP cost,
which matters for the deployment scenarios that motivate this thesis.
EfficientNet appears as the classifier in Kowlagi et al.'s Pfirrmann pipeline
\cite{kowlagi2022a} and in the final stage of M-SCAN \cite{batra2025mscan}.
\textbf{ConvNeXt} represents the most recent convolutional line, modernising the
CNN with design choices borrowed from transformers while retaining convolution
\cite{liu2022convnet}; Udomluck et al.\ include ConvNeXt-Tiny in their feature
extractor comparison and find it stable across classifiers
\cite{udomluck2026multitask}.

The architectural limitation of convolution is a bounded receptive field. Deep
stacks aggregate local responses into progressively broader ones, but the
network never explicitly models a relationship between two distant regions---for
instance, between the L1--L2 and L5--S1 levels, whose degenerative states are
biomechanically coupled.

\subsection{Vision Transformers and the Attention Bottleneck}
\label{sec:lr-vit}

Transformers dispense with convolution and rely on self-attention, in which
every element of the input is weighted against every other regardless of
distance \cite{vaswani2017attention}. The Vision Transformer applies this by
partitioning an image into fixed patches, flattening them and treating the
result as a sequence \cite{dosovitskiy2021image}. Global context is available
from the first layer, and there is no locality prior to overcome---but also no
locality prior to exploit, which is why vision transformers are markedly more
data-hungry than convolutional networks of comparable capacity.

The practical obstacle is cost. Global self-attention scales quadratically with
the number of patches, hence quadratically with image area. Clinical MRI demands
high spatial resolution to resolve structures a few millimetres across, so a
naive vision transformer over full-resolution lumbar images is computationally
prohibitive.

\subsection{The Swin Transformer}
\label{sec:lr-swin}

The Swin Transformer resolves this bottleneck through two mechanisms
\cite{liu2021swin}. \emph{Windowed self-attention} restricts attention to
non-overlapping local windows, making cost linear rather than quadratic in image
size; to restore cross-window communication, the window partition is
\emph{shifted} by half a window between consecutive layers, so information
propagates across boundaries over successive blocks. \emph{Hierarchical feature
maps} are produced by patch-merging layers that progressively reduce spatial
resolution while increasing channel depth, yielding a pyramid analogous to a
CNN's and making the backbone usable for dense prediction as well as
classification.

For lumbar grading the appeal is specific. Severity depends simultaneously on
fine local detail---the texture of a nucleus pulposus, the obliteration of
perineural fat---and on regional context, since a rootlet configuration is
interpretable only relative to the surrounding dural sac and the levels above
and below. A hierarchical, windowed attention mechanism is a plausible fit to
that structure, which is the substantive reason for selecting it as the
transformer representative in this thesis. Swin-derived architectures have been
applied to spinal segmentation, and Zheng et al.\ incorporate Swin components in
their quantitative degeneration network \cite{zheng2022deep}.

\subsection{Hybrid CNN--Transformer Architectures}
\label{sec:lr-hybrids}

A third position holds that the two families are complementary. SymTC merges
convolutional and transformer layers in a \emph{parallel dual-path} arrangement
rather than stacking them sequentially, and integrates a position embedding into
the self-attention module to improve spatial awareness; benchmarked against
fifteen existing models on both a private dataset and a released synthetic one,
it performed best for both vertebral bones and discs \cite{chen2024symtc}. Zheng
et al.\ similarly combine convolutional and Swin components
\cite{zheng2022deep}, and Yi et al.\ pair a 3D ResNet18 backbone with transformer
components across three hospitals \cite{yi2023deep}. Baur et al.\ explore a
different pairing, segmenting discs with a U-Net, reconstructing each as a 3D
point cloud and grading it with a graph neural network; performance was uneven
across levels---F1 0.71 at L2/3 against 0.46 overall---which the authors present
as a foundation requiring refinement rather than a competitive result
\cite{baur2024automated}. Hybrids complicate the
clean comparison this thesis undertakes, but they indicate where the field
expects the answer to lie: not in the wholesale replacement of convolution, but
in supplying it with a mechanism for long-range interaction.

\subsection{Published CNN-versus-Transformer Comparisons and Their Limits}
\label{sec:lr-cnn-vs-transformer}

Several studies bear directly on the architectural question, and their results
do not favour transformers as strongly as the general computer vision literature
might suggest.

Kowlagi et al.\ address it explicitly. Observing that recent gains had been
attributed to transformer architectures, they show that a well-tuned three-stage
convolutional pipeline outperforms the state-of-the-art transformer approaches
on Pfirrmann grading, evaluated on the Northern Finland Birth Cohort 1966---a
population cohort rather than a clinical convenience sample---with an ablation
study and subgroup breakdown, and public code \cite{kowlagi2022a}. Their
conclusion is that pipeline design and tuning outweigh backbone choice.

Lee et al.\ provide the closest published analogue to the present study,
developing a CNN-based and a transformer-based model for lumbar stenosis
classification on 564 MRIs with a 100-study external test set, labelled by four
radiologists with expert consensus as reference and read by eight participants
across four experience levels. Both models exceeded 94\% detection recall; on
dichotomised classification the CNN, transformer and human participants achieved
kappas of 0.99/0.99/0.97--0.98 for central canal, 0.98/0.94/0.81--0.94 for
lateral recesses and 0.98/0.95/0.91--0.95 for neural foramina. The CNN was
superior overall, the transformer similar to superior against clinicians
\cite{lee2026deep}.

Udomluck et al.\ compared VGG19, ConvNeXt-Tiny and DINOv2 features with classical
classifiers under external validation, finding conventional VGG19 features
strongest (accuracy 0.9556 with logistic regression) and DINOv2 features
generalising worst, dropping to 0.6222 with LightGBM
\cite{udomluck2026multitask}. Patil et al.\ reached a still more pointed
conclusion through ablation: radiomic features carried essentially the entire
signal for Pfirrmann grading, and frozen ImageNet-pretrained ResNet50 features
added no measurable discriminative value \cite{patil2026mribased}. Silveira et
al.\ likewise found their enhanced U-Net outperformed TransUNet on segmentation
\cite{silveira2025automated}.

Three methodological limitations recur across this evidence. First, the
compared architectures are rarely given identical preprocessing, localisation
and training budgets, so differences may reflect tuning effort rather than
representational capacity---exactly the confound Kowlagi et al.\ identify
\cite{kowlagi2022a}. Second, the input configuration is usually held fixed,
leaving unexamined the interaction between architecture and available context,
even though a transformer's advantage should plausibly be largest where global
context must substitute for missing views. Third, several comparisons rely on
frozen pretrained features rather than fine-tuned backbones
\cite{udomluck2026multitask,patil2026mribased}, which disadvantages transformers
disproportionately given their weaker inductive biases. These three gaps
motivate the experimental design of this thesis.

\subsection{Transfer Learning, Pretraining, and Self-Supervision}
\label{sec:lr-pretraining}

Medical datasets are small by the standards these architectures were designed
for, so initialisation matters. ImageNet pretraining is near-universal, though
Patil et al.'s ablation is a caution that its benefit is not automatic
\cite{patil2026mribased}. Domain-specific pretraining is better supported:
Jamaludin et al.\ obtained equivalent performance from far fewer labels through
longitudinal self-supervision \cite{jamaludin2017selfsupervised}, and Acharya and
Kansakar used contrastive pretraining on disc-centred crops to reduce sensitivity
to irrelevant appearance variation \cite{acharya2026disccentric,chen2020simple}.
Sequential transfer between related targets---initialising a subarticular model
from spinal canal weights---has also been reported to help on smaller datasets, a
strategy of particular relevance to transformers given their lack of
convolutional inductive bias.

% ===========================================================================
\section{Exploiting Spatial Context: Slices, Volumes, and Views}
\label{sec:lr-context}

\subsection{2D, 2.5D, and 3D Input Representations}
\label{sec:lr-dimensionality}

An MRI series is a volume, but severity is assigned per level, so a
representation must be chosen. \emph{2D} models classify individual slices and
are cheap but discard through-plane context; since a disc spans several slices
and a radiologist integrates across them, a single slice may not contain the
evidence. \emph{3D} models consume the volume directly \cite{yi2023deep} at
considerable memory cost---Acharya and Kansakar note that a standard
high-resolution series contains millions of voxels, making full-volume
processing impractical for many architectures \cite{acharya2026disccentric}.

\emph{2.5D} representations occupy the middle ground and have become the
default in this domain. A small stack of adjacent slices is treated as channels
or as a short sequence, capturing local through-plane context at roughly 2D
cost. SpineNet's input was a stack of nine slices per disc
\cite{jamaludin2017spinenet}; M-SCAN describes its approach explicitly as 2.5D,
selecting the three axial slices closest to each level and modelling them with
recurrent layers \cite{batra2025mscan}; the leading challenge solutions combined
2D CNNs with sequence models or attention-based multiple-instance learning over
the slice dimension.

\subsection{Cross-Sequence and Cross-Plane Fusion}
\label{sec:lr-fusion}

Where multiple sequences are available, they must be combined. The simplest
approach concatenates them as channels, which presumes spatial registration that
sagittal and axial acquisitions do not satisfy. More considered strategies
process each sequence separately and fuse at the feature level.

CrossSpine is the most explicit treatment. Nguyen et al.\ identify two failings
of baselines---single-stream designs cannot integrate complementary information
across sequences, and uniform treatment of all discs ignores level-specific
anatomy---and address the first with a cross-sequence attention mechanism that
fuses features at multiple spatial scales, letting each sequence attend to all
others, followed by a weighting mechanism that prioritises the most
diagnostically informative sequence. They report Macro F1 improved by more than
125\%, Mean AUPRC by 99\% and Mean AUROC by 36\% over their baseline
\cite{nguyen2026crossspine}. M-SCAN applies cross-attention to align sagittal and
axial feature maps before classification \cite{batra2025mscan}, and Park et al.\
pursue transformer-based multimodal fusion on the same benchmark
\cite{park2025multiview}. Zeng et al.\ extend fusion beyond imaging entirely,
combining ConvNeXt image features with GPT-4o-generated textual descriptions
encoded by RoBERTa through a Mamba fusion stage \cite{zeng2025gpt4lfs}.

\subsection{Geometric Correspondence Between Planes}
\label{sec:lr-geometry}

Fusing sagittal and axial data requires knowing which axial slice corresponds to
which sagittal level---a geometric problem solved through DICOM spatial
metadata rather than learned. Pan et al.\ derive the correspondence from the
geometric relationship between the two acquisitions, reporting 100\% accuracy in
identifying the disc in each axial slice \cite{pan2021automatically}. M-SCAN
projects 2D sagittal stenosis coordinates into 3D using patient orientation data
from the DICOM headers, then selects the three nearest axial slices for each of
the five levels \cite{batra2025mscan}.

This step is a hidden cost of multi-sequence modelling, and one that bears
directly on the trade-off examined here. It adds a preprocessing dependency, a
failure mode when metadata is absent or inconsistent, and an engineering burden
that a single-sequence pipeline simply does not incur.

\subsection{Multi-Sequence versus Single-Sequence Input}
\label{sec:lr-multiseq}

The clinical standard is unambiguous that all three sequences contribute
\cite{richards2026the,hallinan2021deep}, and the strongest reported results on
the RSNA benchmark use all three \cite{batra2025mscan}. The countervailing
consideration is cost: multi-sequence processing demands greater memory and
compute, cross-plane registration, and longer inference, and the resulting
infrastructure requirements exceed what many hospital environments provide.

The evidence on whether reduced input suffices is genuinely mixed, which is why
the question remains open. On one side, external validation of sagittal-only
SpineNetV2 found agreement dropping specifically for foraminal stenosis and
herniation---the pathologies requiring transverse context
\cite{nigru2024external}---and Wu et al.\ found ordinal Pfirrmann grading
degrading in upper lumbar discs while binary detection held up
\cite{wu2026external}. On the other, Acharya and Kansakar grade severity from
\emph{sagittal T2 alone}, reaching 78.1\% balanced accuracy with a
severe-to-normal misclassification rate of 2.13\%, on the argument that
anatomical focus rather than input breadth is the binding constraint
\cite{acharya2026disccentric}. Kowlagi et al.\ likewise operate on sagittal input
and outperform transformer baselines \cite{kowlagi2022a}.

Two observations sharpen the question. First, the comparison has not been made
under controlled conditions: the sagittal-only and multi-sequence results above
come from different architectures, datasets, localisation strategies and
evaluation protocols, so the contribution of sequence configuration cannot be
isolated. Second, the expected penalty is target-dependent---severe for
foraminal and subarticular grading, mild for central canal---so a single
aggregate accuracy figure would obscure precisely the structure of interest.
Both observations inform the design adopted in Chapter~3.

\subsection{Robustness to Missing or Incomplete Sequences}
\label{sec:lr-missing}

Clinical data is incomplete: sequences are omitted, truncated or corrupted by
artifact. The LumbarDISC inclusion criteria required all three sequences and
excluded studies with substantial artifact \cite{richards2026the}, so a model
trained on it has not been exposed to the degraded conditions of routine
practice. A model that requires three sequences fails outright when one is
absent, whereas a single-sequence model is unaffected---an argument for reduced
input configurations that is about robustness rather than efficiency, and one
that the literature on missing-modality learning addresses only in passing for
this application.

% ===========================================================================
\section{Class Imbalance and Ordinal Severity Modelling}
\label{sec:lr-imbalance}

\subsection{Severity Distributions in Lumbar MRI Datasets}
\label{sec:lr-distributions}

Class imbalance in this domain is not an artefact of sampling but a property of
the population: in any unselected clinical cohort, most disc levels are normal or
mildly degenerate, and severe compression is uncommon. The RSNA LumbarDISC
distributions make the scale explicit \cite{richards2026the}. Of 13{,}474 spinal
canal stenosis grades, 85.4\% are normal or mild, 8.8\% moderate and 5.9\%
severe. Of 26{,}919 neural foraminal grades, roughly 78\% are normal or mild,
18\% moderate and 4.5\% severe. Subarticular grading is the least skewed of the
three, with roughly 69\% normal or mild, 19\% moderate and 11\% severe across
26{,}285 grades.

The same pattern recurs elsewhere. Liawrungrueang et al.\ report a Pfirrmann
distribution of 220 grade~I, 530 grade~II, 170 grade~III, 160 grade~IV and just
20 grade~V across 1{,}000 images \cite{liawrungrueang2023automatic}; Niemeyer et
al.\ describe an approximately Gaussian grade distribution with grades~1 and~5
under-represented \cite{niemeyer2021a}. Two consequences follow. First, a model
trained under unweighted cross-entropy can achieve roughly 85\% accuracy on
spinal canal stenosis by predicting \emph{normal/mild} unconditionally, so
aggregate accuracy is close to uninformative on this task. Second, the ordering
of imbalance across targets is the inverse of the ordering of difficulty:
subarticular stenosis has the most balanced label distribution yet the poorest
human reliability \cite{lurie2008reliability}, confirming that its difficulty is
intrinsic rather than a consequence of scarce positive examples.

\subsection{Data-Level Strategies}
\label{sec:lr-data-level}

The most common data-level intervention is weighted random sampling, in which
the probability of drawing a minority-class example is inflated so that each
minibatch approaches a uniform severity distribution and gradients are not
dominated by the majority class. Aggressive augmentation of minority classes
complements this, since with few severe examples a network may memorise them
rather than learn the morphology of severe compression. Reported techniques
include contrast-limited adaptive histogram equalisation, random scaling, affine
rotation and elastic deformation, alongside mixing methods.

Tsai et al.\ provide the clearest single demonstration of augmentation's value
in the small-data regime, treating it as the subject of the study rather than an
implementation detail: mean average precision reached 92.4\% at 550 images with
augmentation, and augmentation was decisive in preventing both under- and
over-fitting \cite{tsai2021lumbar}. Ahmed et al.\ attack a related form of
imbalance---between anatomical structures rather than severity classes---through
preprocessing that rectifies class inconsistencies before training
\cite{ahmed2024pioneering}.

\subsection{Loss-Level Strategies}
\label{sec:lr-loss-level}

Class-weighted cross-entropy multiplies each class's contribution to the loss by
a scalar, forcing the optimiser to penalise minority-class errors more heavily.
The RSNA challenge encoded this directly into its evaluation metric, assigning
weights of 1, 2 and 4 to normal/mild, moderate and severe respectively, so that
aligning the training objective with the evaluation criterion is a matter of
adopting the same coefficients.

Focal loss goes further, introducing a modulating factor that down-weights
well-classified examples dynamically and concentrates gradient on hard,
borderline cases \cite{lin2017focal}. This is well matched to the present
problem, where the abundant normal cases are mostly easy and the difficulty
concentrates at the mild/moderate and moderate/severe boundaries---precisely
where human readers also disagree. Acharya and Kansakar combine weighted focal
loss with contrastive pretraining \cite{acharya2026disccentric}, and Islam~Sk et
al.\ propose an adaptive PID-Tversky loss that borrows feedback control to adjust
training penalties dynamically toward under-segmented minority instances
\cite{islamsk2026an}.

\subsection{Ordinal-Aware Objectives and Whether They Help}
\label{sec:lr-ordinal-losses}

Given the ordinal label structure of Section~\ref{sec:lr-ordinality}, it is
natural to expect ordinal-aware objectives to outperform categorical ones. The
evidence does not support this expectation as strongly as the intuition
suggests.

Niemeyer et al.\ conducted the most systematic test available, evaluating soft
kappa loss, ordinal cross-entropy and regression losses against standard
cross-entropy for Pfirrmann grading on 1{,}599 patients and 7{,}948 discs, with
extensive hyperparameter optimisation, and additionally trying class pooling and
class-weighted losses to counter imbalance. None of these improved classification
performance overall. The default cross-entropy classifier reached Cohen
$\kappa = 0.92$, average sensitivity 90.2\% and precision 92.5\%, with
predictions within one grade of ground truth in 99.2\% of cases and a mean
absolute error of 0.08 grades \cite{niemeyer2021a}. Patil et al.\ list ordinal
regression among the supplementary analyses they ran, and still report that
intermediate-grade discrimination remains the outstanding difficulty
\cite{patil2026mribased}.

Two readings are possible. Either the ordinal structure is already implicitly
captured---a network trained with cross-entropy on visually ordered classes will
naturally confuse adjacent grades more often than distant ones---or the specific
formulations tested were inadequate. The literature does not currently
distinguish these, which is why this thesis treats the choice of objective as an
empirical question rather than a settled one.

\subsection{Contrastive and Representation-Learning Approaches}
\label{sec:lr-contrastive}

A third strategy attacks imbalance in the representation rather than in the
sampling or the loss. Contrastive learning pulls examples of similar severity
together in the latent space and pushes dissimilar ones apart before any
classifier is fitted \cite{chen2020simple}, so the representation is shaped by
the structure of the problem rather than by class frequency.

Acharya and Kansakar apply this directly, combining contrastive pretraining with
disc-level fine-tuning over a single anatomically localised region of interest
per disc, on the reasoning that this forces attention onto intrinsic disc
features and reduces sensitivity to irrelevant appearance variation. An auxiliary
regression task handles disc localisation and a weighted focal loss addresses the
residual imbalance. The result---78.1\% balanced accuracy with the
severe-to-normal misclassification rate reduced to 2.13\% relative to supervised
training from scratch---is notable less for the headline accuracy than for what
it optimises \cite{acharya2026disccentric}. The lineage runs back to Jamaludin et
al.'s longitudinal self-supervision \cite{jamaludin2017selfsupervised}.

\subsection{The Clinical Asymmetry of Error}
\label{sec:lr-error-asymmetry}

The most important point in this section is that not all errors are equal, and
that standard metrics conceal the difference. Grading a severe stenosis as
normal risks a missed diagnosis and delayed decompression; grading a normal disc
as moderate risks an unnecessary review. Overall accuracy, macro-F1 and even
Cohen's $\kappa$ treat these identically.

Three studies take the asymmetry seriously. Acharya and Kansakar report the
severe-to-normal misclassification rate as a headline result alongside balanced
accuracy, arguing explicitly that conventional metrics do not capture the
differential clinical risk \cite{acharya2026disccentric}. Wu et al.\ characterise
SpineNetv2's error profile as specificity-oriented, with false negatives
exceeding false positives, and conclude that the system is suitable as a
confirmatory second reader but that reliance on its negative findings is not
recommended---a conclusion about deployment that follows from the shape of the
errors rather than their number \cite{wu2026external}. The RSNA weighted metric
encodes the same judgement into the benchmark itself.

Ishimoto et al.\ supply the complementary observation: collapsing four grades to
a severe-versus-rest dichotomy raised agreement from 65.7\% to 94.1\% with
$\kappa = 0.75$ \cite{ishimoto2020could}. If the clinically decisive question is
whether a level is severe, then performance on that question is considerably
better than the multi-class figures suggest, and evaluation should report both.
This chapter therefore recommends, and Chapter~3 adopts, reporting per-class
recall for the severe class alongside aggregate measures.

% ===========================================================================
\section{Datasets and Benchmarks}
\label{sec:lr-datasets}

\subsection{Private and Single-Institution Cohorts}
\label{sec:lr-private-data}

Most work before 2024 used private single-institution data, with cohort sizes
spanning three orders of magnitude---from 75 exams \cite{gao2022automatic} and
146 consecutive patients \cite{lehnen2021detection} through 200 studies
\cite{bharadwaj2023deep} and 542 images \cite{won2020spinal} to 4{,}075 patients
\cite{lu2018deepspine} and 15{,}254 axial images \cite{su2022automatic}.
Aggregate performance correlates only loosely with cohort size, because larger
private datasets are frequently drawn from a single scanner and protocol.

Two consequences follow, and both are structural rather than incidental. Results
are not comparable across studies, since each is measured against a different
label distribution and reference standard; and models cannot be independently
re-evaluated, since the data cannot be shared. Both systematic reviews of the
period identify this as the field's central methodological weakness
\cite{compte2023are,mendes2026deep}.

\subsection{Population and Epidemiological Cohorts}
\label{sec:lr-population-data}

A smaller body of work uses population cohorts, which differ from clinical
convenience samples in a way that matters for generalisation: participants are
not selected by referral, so the severity distribution reflects the population
rather than the imaged-because-symptomatic subset. The Northern Finland Birth
Cohort 1966, comprising lumbar MRI from participants imaged at age 46--47, is the
most used \cite{mcsweeney2023external,kowlagi2022a}; the Wakayama Spine Study
\cite{ishimoto2020could} and the Hong Kong Disc Degeneration Population-Based
Cohort \cite{cheung2022learningbased} serve similar roles. Kowlagi et al.\ are
explicit that evaluating on a population cohort rather than a clinical sample is
a deliberate methodological choice, reflecting the transferability of results to
downstream research \cite{kowlagi2022a}.

\subsection{Public Lumbar MRI Datasets}
\label{sec:lr-public-data}

Public data arrived late to this field relative to CT, where large annotated
vertebral collections have existed for years. SPIDER is the principal
segmentation resource: 447 sagittal T1 and T2 series from 218 patients across
four hospitals, with reference segmentations of vertebrae, discs and spinal
canal, per-level radiological gradings, baseline and nnU-Net reference
performance, and a continuous challenge on a sequestered split
\cite{vandergraaf2024lumbar}. Its influence is visible in the number of
subsequent methods benchmarked on it
\cite{ahmed2024pioneering,silveira2025automated,patil2026mribased,ghobrial2025deep}.
Smaller public collections support classification work
\cite{natalia2024lumbar,shahzadi2023nerve,nikpasand2024automated}.

\subsection{The RSNA 2024 Challenge and LumbarDISC}
\label{sec:lr-lumbardisc}

The RSNA Lumbar Degenerative Imaging Spine Classification dataset is the largest
and most diverse expertly annotated lumbar spondylosis collection assembled, and
it is the dataset used in this thesis \cite{richards2026the}. It comprises MRI
studies from 2{,}697 patients and 8{,}593 image series drawn from eight
institutions across six countries and five continents---a geographic spread
chosen deliberately so that models trained on it face genuine acquisition
heterogeneity rather than a single site's conventions.

Inclusion required a sagittal T2-like series (conventional spin echo T2, STIR or
Dixon), a sagittal T1 series and an axial T2 series covering at least three
lumbar disc levels, from outpatients aged 18 or over imaged for degenerative
disease. Studies with lumbosacral hardware, active non-degenerative pathology,
severe scoliosis or substantial artifact were excluded \cite{richards2026the}.
Annotation was performed by expert volunteer neuroradiologists and
musculoskeletal radiologists recruited through the ASNR
\cite{rsna2024rsna}. The associated challenge \cite{rsna2024rsna} made the
dataset the field's common benchmark and shifted attention decisively from
whole-image classification toward anatomically localised, per-level severity
grading.

\subsection{Label Structure, Distribution, and the Weighted Metric}
\label{sec:lr-label-structure}

The label structure is the defining feature of the benchmark and warrants precise
statement. Each study receives twenty-five graded outputs: five conditions
---spinal canal stenosis, left and right neural foraminal narrowing, left and
right subarticular stenosis---at each of five disc levels from L1/L2 to L5/S1.
Each is assigned one of three ordinal classes: Normal/Mild, Moderate or Severe.
Expert coordinate annotations accompany the grades, so anatomical localisation
can be learned in a supervised manner rather than inferred
\cite{rsna2024rsna,richards2026the}.

Three properties follow. The task is \emph{not} generic three-class image
classification, and results from studies that treat it as such are not
comparable. The severity distribution is severely imbalanced and differently so
for each condition (Section~\ref{sec:lr-distributions}). And the challenge metric
is a sample-weighted log loss with weights of 1, 2 and 4 across the three
classes, which encodes the clinical asymmetry of
Section~\ref{sec:lr-error-asymmetry} directly into the evaluation. A model tuned
for unweighted accuracy is therefore optimising a different objective from the
one the benchmark measures.

\subsection{Challenge Solutions and the Practices They Standardised}
\label{sec:lr-challenge-solutions}

The solutions converged on a common template, and that convergence is itself a
finding. Nearly all adopted multi-stage pipelines in which localisation precedes
classification: a detection model---commonly a U-Net, YOLO variant or Faster
R-CNN with a feature pyramid---identifies each disc level on the sagittal
series, and a second-stage classifier grades tightly cropped regions, typically
using 2.5D representations that combine a 2D backbone with a sequence model or
attention over adjacent slices.

M-SCAN exemplifies the approach and reports AUROC 0.971 on 1{,}975 studies,
using a U-Net with ResNet50 encoder to locate stenosis points on sagittal images,
projecting these into 3D via DICOM orientation metadata to select corresponding
axial slices, and applying cross-attention to align sagittal and axial features
\cite{batra2025mscan}. Liu et al.\ pair MobileNetV2 with U-Net for efficiency,
reporting 94.93\% accuracy under five-fold cross-validation
\cite{liu2025automated}; Park et al.\ apply transformer-based multimodal fusion
\cite{park2025multiview}; and further challenge-derived work continues to appear
\cite{patel2025automated,garciadecelis2025deep}.

Two cautions attach to these figures. Most are cross-validation results on the
public training split rather than external validation, so they measure fit to
LumbarDISC rather than transportability beyond it. And the top solutions rely on
large ensembles, which inflate reported accuracy relative to what a single
deployable model achieves---a gap directly relevant to the efficiency question
this thesis examines.

% ===========================================================================
\section{Evaluation, Validation, and Generalization}
\label{sec:lr-validation}

\subsection{Metrics and Their Appropriateness}
\label{sec:lr-metrics}

Reported metrics in this field are heterogeneous and frequently ill-matched to
the task. Overall accuracy remains the most common headline figure despite being
close to meaningless under the severity distributions of
Section~\ref{sec:lr-distributions}. Area under the ROC curve is widely reported
\cite{batra2025mscan,wang2024machine} but is defined for binary problems and must
be extended by one-versus-rest averaging for three-class grading, which conceals
which class is failing.

More appropriate choices appear in the better studies. Cohen's and Gwet's
$\kappa$ measure agreement above chance and are directly comparable with the
human reliability figures of Section~\ref{sec:lr-reliability}, which is the
correct reference point \cite{lim2022improved,su2022automatic,mcsweeney2023external}.
Balanced accuracy corrects for prevalence \cite{mcsweeney2023external,acharya2026disccentric}.
Mean absolute error in grades exploits the ordinal structure and distinguishes a
one-grade slip from a two-grade error \cite{niemeyer2021a,wu2026external}.
Ishimoto et al.\ report both exact-grade agreement and the collapsed
severe-versus-rest figure, making the difference visible rather than choosing
between them \cite{ishimoto2020could}.

The recommendation that follows is to report a small panel rather than a single
number: a prevalence-corrected aggregate, an agreement statistic comparable to
inter-reader values, per-class recall for the severe class, and the distribution
of error magnitudes across grade boundaries.

\subsection{Internal Validation and Optimistic Reporting}
\label{sec:lr-internal-validation}

A recurring pattern is that internally validated performance is high and
externally validated performance is materially lower. Several reported accuracies
in the reviewed literature exceed 95\%
\cite{shahzadi2023nerve,liawrungrueang2023automatic,liu2025automated,tabarestani2026a},
which sits uneasily beside inter-reader agreement of $\kappa = 0.49$ for
subarticular stenosis \cite{lurie2008reliability}. Where a model appears to
exceed the reproducibility of the process that generated its labels, the
explanation usually lies in the evaluation protocol: single-institution data,
cross-validation rather than held-out external testing, labels from a single
reader, or a task collapsed to a binary decision.

The point is not that these studies are wrong on their own terms, but that
their numbers do not transfer. Compte et al.\ state the general finding
directly: algorithms did not perform as well in replication or external
validation as in their development studies \cite{compte2023are}.

\subsection{External Validation and Domain Shift}
\label{sec:lr-generalization}

External validation on geographically and temporally separate cohorts is the
strongest available evidence of clinical viability, and a small number of
studies supply it.

McSweeney et al.\ evaluated SpineNet on 1{,}331 participants from the Northern
Finland Birth Cohort 1966, a dataset separate from its UK training data in both
geography and time. Balanced accuracy was 78\% for disc degeneration and 86\%
for Modic changes, with Pfirrmann reliability against expert radiologists of Lin
concordance 0.86 and Cohen $\kappa = 0.68$, and Modic detection at
$\kappa = 0.74$; both were essentially unchanged in a low-back-pain subset. The
informative detail is the error structure: 20.83\% of discs were graded
differently from human raters, but only 0.85\% differed by more than one grade
\cite{mcsweeney2023external}. Nigru et al.\ evaluated SpineNetV2 across eleven
features on 1{,}747 discs from 353 patients, finding agreement above 0.7 for most
targets but distinctly lower for foraminal stenosis and herniation, attributed
to the limits of sagittal input \cite{nigru2024external}. Wu et al.\ tested
SpineNetv2 on 491 patients and 2{,}455 discs, reporting overall accuracy of
83.5--97.5\% and superiority over a junior orthopaedic surgeon on canal stenosis,
spondylolisthesis and bilateral foraminal stenosis \cite{wu2026external}. Grob et
al.\ contribute a further independent validation on 882 scans
\cite{grob2022external}.

Studies that build external testing into their development are still uncommon
but instructive. Tumko et al.\ evaluated on 150 studies graded by a panel of
seven radiologists with majority voting and external adjudication
\cite{tumko2024a}; Su et al.\ used a 1{,}273-image external set from a different
hospital and observed accuracy falling by roughly 7--10 points relative to
internal testing \cite{su2022automatic}; Zhang et al.\ found classification
accuracy dropping from 87.70\% to 74.23\% externally \cite{zhang2023deep}; Zeng
et al.\ report an unusually small external drop, from 93.7\% to 92.2\%
\cite{zeng2025gpt4lfs}.

\subsection{Detection Versus Fine-Grained Grading}
\label{sec:lr-detection-vs-grading}

A consistent and under-appreciated finding is that these two capabilities
generalise differently. Zhang et al.\ provide the clearest illustration:
moving to an external institution, detection degraded only modestly---mean IoU
0.82 to 0.70, precision 98.4\% to 96.3\%, sensitivity 99.4\% to 97.8\%---while
classification accuracy fell from 87.70\% to 74.23\% \cite{zhang2023deep}. Wu et
al.\ report the same asymmetry in different terms: binary detection across five
pathologies held up well while ordinal Pfirrmann grading remained the limiting
factor, degrading further in older patients and upper lumbar discs
\cite{wu2026external}. Yilihamu et al.\ show the effect within a single study as
task granularity increases: segmentation IoU remained above 97\% externally,
while classification precision fell from 81.21\% to 74.50\% for an 18-category
problem yet held at 92.51\% to 90.07\% for four-category severity
\cite{yilihamu2025quantification}.

The implication is that where a system fails is predictable. Finding the anatomy
is robust; judging how bad it looks is not. Reporting a single aggregate figure
for a pipeline that performs both obscures this, and evaluation should separate
the two stages.

\subsection{Benchmarking Against Clinicians}
\label{sec:lr-clinician-benchmark}

Comparison against human readers is the most clinically meaningful benchmark, and
the better studies construct it carefully. Won et al.\ provide the cleanest
design: two experts each labelled the data, a model was trained on each expert's
labels, and the key result was that the difference between the two experts was
statistically significant while the difference between each expert and their
corresponding model was not \cite{won2020spinal}. Tumko et al.\ compared against
a seven-radiologist panel average and found the model matched or exceeded it on
all three stenosis types \cite{tumko2024a}. Lee et al.\ benchmarked against eight
participants spanning four experience levels \cite{lee2026deep}, and Wu et al.\
against orthopaedic surgeons of differing seniority \cite{wu2026external}.

Two cautions apply. Comparisons against junior readers or against surgeons rather
than radiologists set a lower bar than subspecialty radiological practice
\cite{ke2025integrating,wu2026external}. And readers in these studies typically
work from cropped images without clinical history, which is not how they read in
practice---so such comparisons measure a narrower task than clinical
interpretation.

\subsection{Longitudinal Stability and Reproducibility}
\label{sec:lr-longitudinal}

One-shot accuracy does not establish that a model is stable over repeat imaging,
which matters for monitoring progression. Murto et al.\ supply the only
longitudinal evidence in this collection, comparing SpineNet's Pfirrmann grading
against radiologists in 19 volunteers imaged at age 37 and again at 51.
Radiologist-to-SpineNet concordance was 0.54 at both timepoints, clearly below
the 0.83 and 0.78 observed between radiologists, and SpineNet systematically
assigned grade~1 to several discs and grade~5 to more discs than the readers did
\cite{murto2025comparison}. The cohort is small, but a systematic offset that is
stable across time is exactly the failure mode that one-shot evaluation cannot
detect. Cheung et al.\ approach the temporal dimension differently, predicting
five-year progression from baseline imaging across 1{,}343 participants
\cite{cheung2022learningbased}.

\subsection{Reporting Quality and Recurring Weaknesses}
\label{sec:lr-reporting}

The systematic reviews converge on a consistent diagnosis. Wang et al.\ pooled
twelve studies covering 15{,}044 patients, obtaining pooled sensitivity 0.84 and
specificity 0.87 with SROC AUC 0.92, but with heterogeneity of $I^2$ near 99\%
and four of twelve studies at high risk of bias; they concluded that no current
model is reliable enough for clinical practice, citing the need for accepted
reference standards, larger datasets, external validation and fuller reporting
\cite{wang2024machine}. Compte et al.\ found no significant difference between
algorithm families and identified the development-to-validation performance drop
as the field's central problem \cite{compte2023are}. Mendes et al.\ found that
most of 56 segmentation studies relied on private datasets and lacked external
validation \cite{mendes2026deep}. Liawrungrueang et al.\ report metrics spanning
71.5\% to 99\% across twenty studies---a range so wide as to indicate that the
underlying tasks are not comparable \cite{liawrungrueang2025artificial}.

Four weaknesses recur: single-institution data without external testing; ground
truth from a single reader with no adjudication; task definitions that vary
between studies while sharing a label vocabulary; and metric selection that
flatters imbalanced problems. The present work adopts the public LumbarDISC
benchmark, its panel-derived labels and its weighted metric specifically to
avoid the first, second and fourth of these.

% ===========================================================================
\section{Interpretability and Clinical Integration}
\label{sec:lr-interpretability}

\subsection{Saliency, Evidence Hotspots, and Grad-CAM}
\label{sec:lr-saliency}

The earliest interpretability contribution in this domain is SpineNet's evidence
hotspots, which map output-layer gradients back to the input to produce a heatmap
per radiological score, obtained without any localisation supervision
\cite{jamaludin2017spinenet}. Grad-CAM \cite{selvaraju2017gradcam} has since
become the default. Tabarestani et al.\ use it to confirm that their cascade
attends to clinically relevant regions \cite{tabarestani2026a}, and Silveira et
al.\ report that attention concentrated on vertebral bodies and nucleus pulposus
in spatial correspondence with manual segmentations \cite{silveira2025automated}.

The value of these methods should not be overstated. A saliency map shows where
a network responded, not why, and confirming that attention falls on plausible
anatomy is a weak check---it excludes gross reliance on background artefact but
does not establish that the model has learned the morphological criteria the
grading system specifies.

\subsection{Inherently Interpretable Models}
\label{sec:lr-interpretable}

A stronger position is to build models whose reasoning is legible by
construction. Bharadwaj et al.\ provide the most striking result in this
literature: alongside a Big Transfer classifier they trained a decision-tree
classifier operating on segmentation-derived features for central canal stenosis,
and the interpretable model \emph{outperformed} the black box, reaching
$\kappa = 0.80$ against 0.54 and approaching inter-reader agreement of 0.74--0.86
\cite{bharadwaj2023deep}. Quantitative morphometry belongs to the same family:
a measured dural sac area \cite{ghobrial2025deep} or a quantitative degeneration
criterion \cite{zheng2022deep} is inspectable in a way a logit is not. Patil et
al.\ make the related point that radiomic features carried the entire
discriminative signal for Pfirrmann grading while deep features added nothing
measurable, and that radiomics additionally yields interpretable shape and
intensity descriptors usable as imaging biomarkers \cite{patil2026mribased}.

Together these findings undercut the assumption that interpretability must be
purchased at the cost of accuracy---at least for targets whose clinical
definition is already explicitly morphological.

\subsection{Vision--Language Models and Report Generation}
\label{sec:lr-vlm}

The most recent direction outputs prose rather than labels. Islam~Sk et al.\
present an explainable vision--language framework that classifies severity,
segments the anatomy and generates radiologist-style reports, built on
biomedical foundation models, with a spatial patch cross-attention module that
preserves the anatomical hierarchies global pooling discards and an adaptive
PID-Tversky loss for class imbalance; they report 90.69\% classification
accuracy, macro-averaged Dice 0.9512 and CIDEr 92.80\%
\cite{islamsk2026an}. Zeng et al.\ take a narrower approach, fusing ConvNeXt
image features with GPT-4o-generated descriptions for foraminal stenosis
classification across three centres and seven scanners \cite{zeng2025gpt4lfs}.
These systems are early and their language metrics measure fluency rather than
clinical correctness, but they indicate where the field is heading.

\subsection{Reader-Assistance Studies and Workflow Impact}
\label{sec:lr-reader-assistance}

The most consequential evidence for clinical value comes not from model metrics
but from measuring what changes when readers use the model. Lim et al.\ conducted
the definitive study of this kind: eight radiologists with 2--13 years of
experience read the same studies with and without deep-learning assistance,
separated by a one-month washout, against a 32-year-veteran musculoskeletal
radiologist as reference. Interpretation time per study fell from 124--274
seconds to 47--71 seconds, and interobserver agreement was superior or equivalent
for every stenosis grading. The largest gain fell to general and in-training
radiologists on four-class foraminal stenosis, where $\kappa$ rose from 0.39 to
0.71 and 0.70 \cite{lim2022improved}. Gao et al.\ observed a comparable effect
for Modic assessment, where AI assistance raised agreement between a junior
radiologist and a senior neuroradiologist from $\kappa = 0.52$ to 0.58
\cite{gao2022automatic}.

The pattern in both is that assistance helps least-experienced readers most and
standardises rather than replaces judgement. This reframes the objective: the
target is not a model that outperforms the best radiologist, but one that raises
the floor.

\subsection{Uncertainty and Calibration}
\label{sec:lr-uncertainty}

Uncertainty estimation is conspicuously underdeveloped in this literature. Where
errors concentrate at grade boundaries and where the reference standard is itself
unreliable, a calibrated confidence would let borderline cases be referred rather
than silently graded---a natural fit for triage. Almost none of the reviewed work
reports calibration. Patil et al.\ are the exception, applying isotonic
calibration on a validation split and reporting Brier scores alongside
discrimination metrics \cite{patil2026mribased}. Since modern networks are
systematically overconfident \cite{guo2017calibration}, an uncalibrated severity
probability should not be interpreted as a clinical confidence, and this remains
an open opportunity rather than a solved problem.

% ===========================================================================
\section{Synthesis and Research Gap}
\label{sec:lr-gap}

\subsection{What the Literature Has Established}
\label{sec:lr-established}

Five findings are supported well enough to be treated as settled.

\textbf{Anatomical localisation precedes classification.} From DeepSPINE
\cite{lu2018deepspine} through the interpretable pipelines of Bharadwaj et al.\
\cite{bharadwaj2023deep} to essentially every leading RSNA challenge solution
\cite{batra2025mscan}, systems that localise the disc level before grading it
outperform those that classify whole images. Detection is close to solved, with
several studies reporting near-perfect level assignment
\cite{pan2021automatically,ghosh2012a}.

\textbf{Automated grading is competitive with human reliability, because human
reliability is low.} Inter-reader agreement is substantial for central canal
stenosis but only moderate for foraminal and poor for subarticular grading
\cite{lurie2008reliability}. Against that benchmark, models sit inside the band
of human disagreement \cite{won2020spinal,mcsweeney2023external,tumko2024a}.
This is a real result, but it is a statement about the modesty of the reference
standard as much as the capability of the models.

\textbf{Errors concentrate at adjacent grade boundaries.} Disagreement above one
grade is rare---0.85\% of discs in the largest external validation
\cite{mcsweeney2023external}---and collapsing to a severe-versus-rest decision
raises agreement substantially \cite{ishimoto2020could}. The difficulty lies in
boundary placement, not gross misclassification.

\textbf{Detection generalises across institutions; fine-grained grading does
not.} The asymmetry is consistent and large
\cite{zhang2023deep,wu2026external,yilihamu2025quantification}.

\textbf{Class imbalance must be addressed explicitly.} With 85.4\% of spinal
canal grades normal or mild \cite{richards2026the}, unweighted objectives yield
degenerate solutions, and weighted sampling, weighted or focal losses and
contrastive pretraining are all established mitigations
\cite{lin2017focal,acharya2026disccentric}.

\subsection{What Remains Contested}
\label{sec:lr-contested}

Three questions remain genuinely open, and the disagreement is substantive
rather than terminological.

\emph{Does architectural sophistication help?} The general computer vision
literature would predict transformers to outperform convolutional networks given
sufficient data, but the lumbar evidence does not bear this out. Kowlagi et al.\
report a tuned CNN pipeline beating transformer approaches \cite{kowlagi2022a};
Lee et al.\ find a CNN superior and a transformer merely comparable
\cite{lee2026deep}; Silveira et al.\ beat TransUNet with an enhanced U-Net
\cite{silveira2025automated}; Udomluck et al.\ find conventional VGG19 features
generalising better than DINOv2 \cite{udomluck2026multitask}; and Patil et al.\
find deep features adding nothing over radiomics \cite{patil2026mribased}. Set
against this, hybrid architectures report gains
\cite{chen2024symtc,zheng2022deep}, and cross-sequence attention reports large
ones \cite{nguyen2026crossspine}.

\emph{Does the ordinal structure of the labels admit exploitation?} Errors
demonstrably cluster at adjacent grades, yet the most systematic test of
ordinal-aware objectives found none improved on plain cross-entropy
\cite{niemeyer2021a}.

\emph{How much imaging is actually required?} Clinical practice and the strongest
benchmark results use all three sequences
\cite{richards2026the,batra2025mscan}, while sagittal-only systems report
competitive performance \cite{acharya2026disccentric,kowlagi2022a} and
sagittal-only systems also fail in a specific, predicted way on the lateral
compartments \cite{nigru2024external}.

\subsection{Gap 1: Architecture Comparison Under Controlled Conditions}
\label{sec:lr-gap-architecture}

The comparative studies above are informative but none isolates the variable of
interest. Three confounds recur.

First, \emph{unequal pipelines}. Compared architectures rarely receive identical
preprocessing, localisation and training budgets, so a reported difference may
reflect engineering effort rather than representational capacity---precisely the
argument Kowlagi et al.\ advance when attributing their CNN's advantage to
pipeline design rather than to convolution as such \cite{kowlagi2022a}. Given the
evidence of Section~\ref{sec:lr-localize-first} that localisation dominates
performance, a comparison in which one arm is better localised than the other
measures preprocessing.

Second, \emph{frozen features}. Several comparisons extract features from frozen
pretrained backbones rather than fine-tuning end to end
\cite{udomluck2026multitask,patil2026mribased}. This systematically
disadvantages transformers, whose weaker inductive biases make them more
dependent on in-domain adaptation.

Third, \emph{fixed input configuration}. Every comparison holds the sequence
configuration constant, so the interaction between architecture and available
context is untested. This matters because the theoretical case for global
attention is strongest exactly where context must be inferred rather than
observed---that is, where a view is missing.

A controlled comparison therefore requires a shared localisation stage, a shared
training and augmentation protocol, comparable capacity and tuning budgets, and
end-to-end fine-tuning of both backbones, evaluated on the same benchmark with
the same metric. No study in this review satisfies all four conditions.

\subsection{Gap 2: Sequence Configuration and the Accuracy--Efficiency Trade-off}
\label{sec:lr-gap-sequence}

The second gap concerns input rather than model. Multi-sequence pipelines carry
costs that are rarely quantified alongside their accuracy: greater memory and
compute, cross-plane geometric registration dependent on DICOM metadata
\cite{batra2025mscan,pan2021automatically}, longer inference, and brittleness
when a sequence is missing or degraded (Section~\ref{sec:lr-missing}). Against
these, the incremental diagnostic yield of each additional sequence has not been
measured under controlled conditions.

The literature supplies both sides of the argument but no direct test. Evidence
that the third sequence matters is indirect, inferred from the failure of
sagittal-only systems on foraminal and subarticular targets
\cite{nigru2024external,wu2026external}; evidence that it may not is likewise
indirect, resting on sagittal-only systems that perform respectably against
different baselines on different data
\cite{acharya2026disccentric,kowlagi2022a}. Because these come from different
architectures, cohorts, localisation strategies and metrics, the contribution of
sequence configuration cannot be separated from everything else that differs.

Two design requirements follow. The comparison must hold architecture and
pipeline constant while varying only the input configuration; and it must report
results per condition rather than in aggregate, since the expected penalty is
concentrated in the foraminal and subarticular targets and an aggregate figure
would average away the structure of interest.

\subsection{Research Questions and Objectives}
\label{sec:lr-questions}

The two gaps combine into a single factorial investigation. Crossing
architecture with input configuration answers both questions and, additionally,
tests whether they interact---whether a transformer's global attention can
partially compensate for absent multiplanar data, which neither gap addresses
alone. This yields the following research questions:

\begin{description}
  \item[RQ1] Under matched localisation, training and evaluation protocols, how
  do convolutional backbones (ResNet50, EfficientNet-B4) compare with a
  hierarchical vision transformer (Swin) for per-level, per-condition severity
  grading of lumbar degenerative disease?

  \item[RQ2] What is the cost in diagnostic performance of reducing the input
  from the full multi-sequence configuration (Sagittal T1, Sagittal T2/STIR,
  Axial T2) to a single-sequence configuration (Sagittal T2), and how is that
  cost distributed across the five graded conditions?

  \item[RQ3] Do architecture and sequence configuration interact---specifically,
  is the penalty for single-sequence input smaller for the transformer than for
  the convolutional backbones?

  \item[RQ4] Which class-imbalance strategies most effectively reduce
  clinically consequential errors, in particular severe cases graded as
  normal or mild?
\end{description}

The corresponding objectives are to implement a shared localisation and
preprocessing stage so that all arms receive identical input regions; to train
each architecture under both configurations with matched budgets; to evaluate
using the benchmark's weighted metric together with per-condition and per-class
reporting; and to quantify computational cost---parameters, memory and inference
time---so that any accuracy difference can be weighed against it.

\subsection{Scope and Delimitations}
\label{sec:lr-scope-delimit}

Four boundaries are stated explicitly. The work addresses \emph{severity
grading}, taking localisation as a shared and largely solved precondition
(Section~\ref{sec:lr-detection}) rather than a contribution. It evaluates on
\emph{LumbarDISC} \cite{richards2026the}; because that collection excludes
studies with hardware, non-degenerative pathology, severe scoliosis or
substantial artifact, conclusions do not extend to those populations, and no
independent external cohort is available to test transportability beyond it. It
compares a \emph{bounded set of architectures}---two convolutional and one
transformer---chosen as representative rather than exhaustive, and does not
evaluate hybrid or vision--language designs
\cite{chen2024symtc,islamsk2026an}. Finally, it reports model performance and
computational cost, not clinical impact: the reader-assistance evidence of
Section~\ref{sec:lr-reader-assistance} \cite{lim2022improved} shows that
workflow benefit requires a prospective reader study, which lies outside the
present scope.

Two constraints inherited from the literature bound the interpretation of any
result. Performance is capped by the reliability of the reference standard
(Section~\ref{sec:lr-label-ceiling}), so figures approaching perfect agreement
should be read as evidence of a protocol artefact rather than of clinical
readiness. And single-benchmark evaluation cannot establish generalisation; the
consistent development-to-validation drop documented in
Section~\ref{sec:lr-generalization} applies to this work as much as to any other.

% ===========================================================================
\section{Chapter Summary}
\label{sec:lr-summary}

This chapter reviewed the literature on automated assessment of lumbar
degenerative disease from MRI, from the clinical grading systems that define the
task through to the architectures that currently attempt it.

The clinical sections established that severity labels originate in
morphological grading systems---Pfirrmann for disc degeneration
\cite{pfirrmann2001magnetic}, Schizas and Lee for central canal
\cite{schizas2010qualitative,lee2011a}, Lee for foraminal
\cite{lee2010a}---whose inter-reader reliability ranges from substantial to poor
depending on the compartment \cite{lurie2008reliability}. This imposes a ceiling
on achievable performance and makes agreement statistics, rather than accuracy,
the appropriate measure. The labels are ordinal, errors concentrate at adjacent
boundaries, and the clinical cost of an error depends on its direction.

The methodological sections traced the field from handcrafted features through
SpineNet's multi-task convolutional grading \cite{jamaludin2017spinenet} to
contemporary multi-stage, multi-view pipelines \cite{batra2025mscan}. The most
robust design principle to emerge is that anatomical localisation should precede
classification \cite{lu2018deepspine,kowlagi2022a}. The most robust empirical
regularity is that detection transfers across institutions while fine-grained
grading does not \cite{zhang2023deep,wu2026external}.

Two questions remain unresolved. Whether attention-based backbones offer any
advantage over well-tuned convolutional ones is contested, with the balance of
lumbar-specific evidence currently favouring CNNs
\cite{kowlagi2022a,lee2026deep,silveira2025automated}, yet no study has compared
them under matched localisation, training and evaluation conditions. And whether
the full three-sequence protocol is necessary, or whether a single ubiquitous
sequence suffices, has been argued from both sides without a controlled test.
These two questions, and their interaction, constitute the gap this thesis
addresses. Chapter~3 sets out the experimental design through which they are
answered.

% ===========================================================================
\printbibliography[heading=bibintoc,title={References}]

\end{document}
 from your main file.
% ===========================================================================

\documentclass[12pt,a4paper]{report}

\usepackage{lmodern}          % scalable fonts; required for microtype expansion
\usepackage[T1]{fontenc}
\usepackage[utf8]{inputenc}
\usepackage[english]{babel}
\usepackage{microtype}
\usepackage[margin=2.5cm]{geometry}
\usepackage{setspace}
\onehalfspacing

\usepackage{amsmath,amssymb}
\usepackage{booktabs}
\usepackage{longtable}
\usepackage{array}
\usepackage{graphicx}
\usepackage[hidelinks]{hyperref}

\usepackage[
  backend=biber,
  style=numeric-comp,
  sorting=none,
  maxbibnames=6,
  minbibnames=3,
  giveninits=true,
  doi=true,
  url=false,
  isbn=false
]{biblatex}

\addbibresource{../lumbar_spine_mri_ai_literature_inventory.bib}
\addbibresource{references-extra.bib}

% The inventory .bib carries internal bookkeeping fields (provenance note, the
% per-paper summary, the local PDF path). These are working metadata, not
% bibliographic data, so they are suppressed from the printed reference list.
% Leaving them in also breaks compilation, since they contain raw underscores.
\AtEveryBibitem{%
  \clearfield{note}%
  \clearfield{comment}%
  \clearfield{annotation}%
  \clearfield{abstract}%
  \clearfield{file}%
}

\begin{document}

\chapter{Literature Review}
\label{ch:litreview}

% ===========================================================================
\section{Introduction}
\label{sec:lr-intro}

\subsection{Scope and Objectives of the Review}
\label{sec:lr-scope}

Lumbar degenerative disease is among the most common indications for magnetic
resonance imaging (MRI) worldwide, and the interpretation of those images is a
repetitive, subjective and time-consuming task that has attracted sustained
attention from the medical image analysis community \cite{lee2024artificial,zhang2024a}.
This chapter reviews that body of work. Its purpose is not to catalogue every
published model, but to establish what is genuinely settled, what remains
contested, and where a controlled experiment can still add knowledge.

The review is organised around a single guiding question: \emph{what determines
whether an automated system can grade lumbar degenerative severity reliably
enough to be clinically useful?} Four candidate answers recur in the
literature---the quality and definition of the reference labels, the degree to
which the anatomy is localised before it is classified, the choice of
representational backbone, and the configuration of imaging sequences supplied
to the model. Each is examined in turn, and the chapter closes by identifying
the two of these that remain inadequately resolved and that this thesis
therefore addresses.

Three specific objectives follow from this. First, to characterise the clinical
grading systems that supply the ground truth for every supervised model in this
field, and to quantify the reliability ceiling those systems impose. Second, to
trace the methodological evolution from handcrafted features to contemporary
convolutional and attention-based architectures, with particular attention to
comparative claims between the two. Third, to assess the evidence on how much
diagnostic information is carried by each MRI sequence and plane, and thus
whether the computational expense of multi-sequence processing is warranted.

\subsection{Search Strategy, Sources, and Inclusion Criteria}
\label{sec:lr-search}

The literature underpinning this chapter was assembled into a structured
inventory of 92 records spanning 2001 to 2026. Sources were identified through
PubMed, Scopus, IEEE~Xplore and arXiv, using combinations of the terms
\emph{lumbar spine}, \emph{MRI}, \emph{spinal stenosis}, \emph{disc
degeneration}, \emph{deep learning}, \emph{convolutional neural network},
\emph{transformer}, \emph{segmentation} and \emph{grading}. Reference lists of
the retrieved systematic reviews \cite{compte2023are,wang2024machine,liawrungrueang2025artificial,mendes2026deep}
were searched by hand for records the database queries had missed, and the
citation graphs of the most influential frameworks were traced forward.

Records were retained if they satisfied at least one of four criteria: they
define or validate a clinical grading standard for lumbar degenerative disease;
they present an automated method for detecting, segmenting or grading lumbar
pathology on MRI; they externally validate such a method; or they describe a
dataset or benchmark on which such methods are trained. Studies confined to
computed tomography or plain radiography were excluded unless they establish a
methodological principle transferable to MRI, as were studies of the cervical
or thoracic spine except where explicitly used for contrast.

Two caveats on the evidence base should be recorded at the outset. First, the
field's publication rate accelerated sharply after 2024, and a substantial
fraction of the most directly relevant work
\cite{batra2025mscan,acharya2026disccentric,nguyen2026crossspine,islamsk2026an}
exists as preprints that have not completed peer review; these are cited as
indicative of methodological direction rather than as settled results. Second,
performance figures reported across studies are frequently not comparable,
because they derive from different datasets, different label definitions and
different evaluation protocols. Wherever numerical results are quoted in this
chapter, the cohort and validation regime that produced them are stated
alongside, and cross-study comparisons are avoided unless the underlying
protocols genuinely align.

\subsection{Structure of the Chapter}
\label{sec:lr-structure}

Sections~\ref{sec:lr-clinical} to~\ref{sec:lr-grading} establish the clinical
substrate: the disease, the imaging protocol used to observe it, and the
grading systems that convert observation into ordinal labels. This material is
not background decoration. The severity classes a model is trained to predict
are defined entirely by these systems, and their documented unreliability
propagates directly into every reported accuracy figure in the chapters that
follow.

Section~\ref{sec:lr-evolution} traces the methodological history of the field,
and Section~\ref{sec:lr-localization} isolates anatomical localisation as the
design principle that most consistently separates effective systems from
ineffective ones. Sections~\ref{sec:lr-backbones} and~\ref{sec:lr-context} then
address the two axes on which this thesis experiments: the representational
backbone, and the spatial and sequential context supplied to it.
Section~\ref{sec:lr-imbalance} treats class imbalance and the ordinal structure
of the label space, and Section~\ref{sec:lr-datasets} the datasets and
benchmarks---principally the RSNA LumbarDISC collection---on which the field now
converges. Sections~\ref{sec:lr-validation} and~\ref{sec:lr-interpretability}
consider validation practice and interpretability respectively.
Section~\ref{sec:lr-gap} synthesises the whole and states the research gap.

% ===========================================================================
\section{Clinical Background: Lumbar Degenerative Disease}
\label{sec:lr-clinical}

\subsection{Epidemiology and Socioeconomic Burden}
\label{sec:lr-epidemiology}

Low back pain is the single largest contributor to years lived with disability
worldwide, and lumbar degenerative change is among its most frequently
implicated structural correlates \cite{compte2023are,zheng2022deep}. Lifetime
prevalence is reported as high as 84\%, and in the United States alone the
combined direct medical and indirect productivity costs are estimated between
\$50 and \$100 billion annually \cite{richards2026the}. Lumbar spinal stenosis
in particular is the most common indication for spinal surgery in patients over
65 \cite{schizas2010qualitative,bharadwaj2023deep}, and demographic ageing is
expected to increase both its incidence and the imaging workload it generates
\cite{lee2024artificial}.

The operational consequence for this thesis is one of volume. MRI referrals for
low back pain have grown faster than the radiology workforce able to report
them \cite{ghosh2014supervised,compte2023are}, and each lumbar study requires
the reader to assign severity at multiple anatomical levels across several
sequences---an average of 14 to 19 minutes of expert time per study
\cite{compte2023are}. It is this combination of high volume, repetitive
structure and documented inter-reader variability that makes the task a
plausible target for automation, rather than any claim that machines perceive
the underlying pathology more acutely than radiologists do.

\subsection{The Degenerative Cascade}
\label{sec:lr-cascade}

Degeneration of the lumbar spine is a progressive, multifactorial process
driven by ageing, mechanical loading, genetic predisposition and lifestyle
factors \cite{esposito2025mribased,zhao2025advances}. It begins in the
intervertebral disc, whose nucleus pulposus loses proteoglycan content and
therefore water from approximately the third decade of life onward. The
resulting dehydration reduces the disc's height, its capacity to distribute
axial load, and its signal intensity on T2-weighted imaging---the last of these
being precisely what the Pfirrmann scale measures \cite{pfirrmann2001magnetic}.

Disc height loss initiates a compensatory cascade through the rest of the motion
segment. The annulus fibrosus weakens and may permit bulging or frank herniation
of nuclear material; the ligamentum flavum buckles and hypertrophies; the
posterior facet joints, now bearing altered load, develop osteoarthritic change
and osteophytosis \cite{bharadwaj2023deep,nikpasand2024automated}. Changes in
the adjacent vertebral endplates and marrow---the Modic
phenotypes~\cite{modic1988degenerative}---frequently accompany this process and
have themselves been the target of automated detection
\cite{gao2022automatic,athertya2021classification}. The clinical significance of
the cascade is that these individually modest changes act cumulatively to
encroach on the spaces occupied by neural tissue.

\subsection{The Three Anatomical Forms of Stenosis}
\label{sec:lr-stenosis-types}

Encroachment is conventionally partitioned by anatomical compartment, and this
partition is directly mirrored in the label structure of the dataset used in
this thesis \cite{richards2026the}:

\begin{enumerate}
  \item \textbf{Spinal canal stenosis} is narrowing of the central canal, which
  contains the thecal sac and the cauda equina. It is best appreciated on axial
  T2-weighted images, where the relationship between cerebrospinal fluid and
  nerve rootlets is visible \cite{schizas2010qualitative,lee2011a}. Severe
  central stenosis characteristically produces neurogenic claudication.

  \item \textbf{Neural foraminal narrowing} is constriction of the lateral
  foramina through which the exiting nerve roots pass. It is assessed on
  parasagittal images, where the epidural fat surrounding the root provides the
  contrast that makes compression visible---T1 weighting being particularly
  suited to this \cite{lee2010a}. It produces dermatomal radicular symptoms.

  \item \textbf{Subarticular stenosis}, also termed lateral recess stenosis,
  affects the zone between the medial edge of the facet and the foramen, where
  it compresses the traversing root before it exits at the level below
  \cite{lurie2008reliability,andreisek2014consensus}.
\end{enumerate}

The distinction matters computationally as well as clinically, because the three
compartments are not equally visible in any single plane. Foraminal narrowing is
poorly assessed from axial images and subarticular stenosis poorly from
sagittal ones, a constraint that Section~\ref{sec:lr-context} shows to have
concrete consequences for model design. It is also the compartment that
correlates with which sequence a model must be given, and thus sits at the
centre of the multi-sequence question this thesis examines.

\subsection{Clinical Presentation and the Treatment Decision}
\label{sec:lr-treatment}

Severity grading is not an academic exercise; it partitions patients into
distinct management pathways. Conservative management---physiotherapy, analgesia,
epidural steroid injection---is appropriate for mild and most moderate disease,
whereas severe compression with corresponding symptoms may warrant decompressive
surgery or fusion \cite{schizas2010qualitative,mysliwiec2010msu}. The evidence
that morphological grade carries genuine prognostic weight is strongest for the
Schizas classification, where grade~C and~D patients were substantially more
likely to fail conservative treatment \cite{schizas2010qualitative}, and for the
MSU herniation grades, where size~2 and~3 lesions identify microdiscectomy
candidates \cite{mysliwiec2010msu}.

Two qualifications temper this. First, the correlation between imaging findings
and symptoms is imperfect: anatomical stenosis is demonstrable in a substantial
proportion of asymptomatic individuals \cite{lurie2008reliability,esposito2025mribased},
so imaging severity constrains rather than determines management. Second, and
more consequentially for automated systems, the clinically decisive boundary is
not evenly distributed across the ordinal scale. The distinction between
\emph{normal} and \emph{mild} rarely changes management, whereas the distinction
between \emph{moderate} and \emph{severe} frequently does. A model that
distributes its errors uniformly across grade boundaries is therefore not
clinically equivalent to one of the same overall accuracy that concentrates its
errors at the harmless end---a point developed in
Section~\ref{sec:lr-error-asymmetry}.

% ===========================================================================
\section{MRI Acquisition: Sequences, Planes, and Diagnostic Yield}
\label{sec:lr-mri}

\subsection{Why MRI is the Reference Modality}
\label{sec:lr-why-mri}

MRI is the reference standard for non-invasive assessment of the lumbar spine.
Computed tomography and plain radiography depict osseous anatomy well but
resolve soft tissue poorly, and both expose the patient to ionising radiation.
MRI resolves the intervertebral disc, thecal sac, individual nerve roots and
ligamentous structures with a contrast unavailable to the alternatives
\cite{zhao2025advances,liawrungrueang2025artificial}, and it is the disposition
of exactly these soft-tissue structures that defines every grading system
reviewed in Section~\ref{sec:lr-grading}. MRI has correspondingly displaced CT
myelography for stenosis assessment \cite{tumko2024a}.

\subsection{Complementary Information Across Sequences and Planes}
\label{sec:lr-sequences}

Lumbar protocols are inherently multiplanar and multisequence, because the spine
is a three-dimensional structure whose pathologies are not all visible in one
projection. Three acquisitions dominate clinical practice and constitute the
input available in the dataset used here \cite{richards2026the}.

\emph{Sagittal T2-weighted} imaging renders fluid hyperintense, so cerebrospinal
fluid in the thecal sac appears bright against darker ligament and bone. It is
the sequence on which disc hydration---and therefore Pfirrmann
grade---is judged \cite{pfirrmann2001magnetic}, and on which the
craniocaudal extent of canal compromise is surveyed. Short tau inversion
recovery (STIR) suppresses the signal from fat and thereby exposes marrow
oedema and inflammatory change, making it the sequence of choice for active
Modic type~1 change \cite{gao2022automatic}.

\emph{Sagittal T1-weighted} imaging renders fat hyperintense. Because the exiting
nerve root within the neural foramen is normally surrounded by epidural fat, T1
provides the contrast against which obliteration of that fat---the basis of the
Lee foraminal grading system---becomes visible \cite{lee2010a}.

\emph{Axial T2-weighted} imaging provides the transverse cross-section through
the disc level. It is the only view in which the morphology of the dural sac,
the arrangement of the cauda equina rootlets and the lateral recesses can be
directly assessed, and it therefore underpins both the Schizas and Lee central
canal systems \cite{schizas2010qualitative,lee2011a} and the assessment of
subarticular stenosis \cite{andreisek2014consensus}.

\subsection{Matching Pathology to Sequence}
\label{sec:lr-sequence-matching}

The correspondence between pathology and sequence is not a matter of
convention but of physics, and the empirical literature bears this out with
unusual clarity. The most direct evidence comes from external validation of
sagittal-only systems. Nigru et al.\ evaluated SpineNetV2 across eleven
radiological features on 1{,}747 lumbosacral discs and found agreement above
$\kappa = 0.7$ for most targets but distinctly lower agreement for foraminal
stenosis and disc herniation, attributing the shortfall specifically to the
absence of transverse anatomical context in sagittal input
\cite{nigru2024external}. Hallinan et al.\ approached the same constraint from
the opposite direction, deliberately assigning axial T2 to central canal and
lateral recess grading while reserving sagittal T1 for the neural foramina
\cite{hallinan2021deep}.

This has a direct bearing on the experiments in this thesis. If a
single-sequence model is to be evaluated, the choice of which sequence to
retain is not arbitrary: sagittal T2 is the most defensible candidate because it
is acquired in essentially every back-pain protocol and because it carries the
signal for canal stenosis and disc degeneration, but the literature predicts
that it should be weakest precisely on the foraminal and subarticular targets.
Whether a sufficiently expressive architecture can compensate for that
predicted weakness is an open question, and one this thesis is designed to
answer.

\subsection{Sources of Acquisition Variability}
\label{sec:lr-variability}

Models trained on one institution's images frequently degrade on another's, and
the proximate causes lie in acquisition. Scanner manufacturer and field strength
(typically 1.5 versus 3.0~T), slice thickness and spacing, in-plane resolution,
echo and repetition times, the presence or absence of fat suppression, and
patient positioning all vary between sites \cite{senzgamboa2023automatic,mendes2026deep}.
Axial acquisitions vary further in whether slices form a continuous stack or are
angled to lie parallel to each disc level---a difference the LumbarDISC
inclusion criteria explicitly accommodate \cite{richards2026the}, and one that
any model consuming axial data must tolerate.

S\'aenz~Gamboa et al.\ treat this variability as the research problem in its own
right rather than as a nuisance, evaluating segmentation across multiple centres
and acquisition parameters and finding that ensembles of U-Net variants were
required to maintain performance \cite{senzgamboa2023automatic}. Batra et al.\
make the complementary design argument: rather than assuming uniform histograms
and a fixed slice count per study, their pipeline is built to accept a variable
number of axial and sagittal slices, on the grounds that this reflects the
conditions under which clinical data actually arrives \cite{batra2025mscan}.
Variability is thus not merely a threat to validity but a specification
constraint, and Section~\ref{sec:lr-generalization} returns to how effectively
the field has met it.

% ===========================================================================
\section{Clinical Grading Systems as the Origin of Ground Truth}
\label{sec:lr-grading}

Every supervised model reviewed in this chapter learns to reproduce labels
generated by a human applying one of the systems described below. The
properties of those systems---their definitions, their granularity and above all
their reliability---therefore bound what any model trained on them can achieve.
This section treats them as the measurement instruments they are.

\subsection{Disc Degeneration: Pfirrmann and its Modified Extension}
\label{sec:lr-pfirrmann}

The Pfirrmann classification, introduced in 2001, remains the most widely
adopted scheme for intervertebral disc degeneration and is applied to sagittal
T2-weighted images \cite{pfirrmann2001magnetic}. It assigns one of five ordinal
grades on the basis of four morphological criteria: homogeneity of the nucleus
pulposus, distinction between nucleus and annulus, T2 signal intensity as a
proxy for hydration, and disc height. Grade~I denotes a homogeneous, bright,
normally proportioned disc; grade~V a collapsed disc of uniformly low signal.

Its virtues are simplicity and universality; its limitations are well
documented. It is a subjective visual judgement, it is insensitive to change
over short intervals, and it captures only intrinsic disc state, saying nothing
about Modic change, high-intensity zones or the downstream compression of neural
structures \cite{esposito2025mribased,zhao2025advances}. It also suffers a
ceiling effect in elderly cohorts, where most discs present as grade~IV or~V and
clinically meaningful variation is compressed into the top of the scale.
Griffith et al.\ addressed this with an eight-level modified scale supported by a
24-image reference panel, validated on 260 discs in 52 subjects of mean age 73;
intra-observer agreement was excellent (weighted $\kappa$ 0.79--0.91) and
inter-observer agreement substantial (0.65--0.67), with most disagreement
confined to a single grade \cite{griffith2007modified}.

Esposito et al.\ place these two systems in a wider context that is worth
recording: a scoping review of 569 studies identified \emph{93} distinct grading
systems for lumbar disc degeneration---63 subjective, 25 quantitative---of which
the Pfirrmann method accounted for just over half of all reported use
\cite{esposito2025mribased}. The apparent standardisation of the field is
therefore partly illusory, and comparisons of automated grading accuracy across
studies must confirm that the same instrument was being reproduced.

\subsection{Central Canal Stenosis: Schizas and Lee}
\label{sec:lr-central-grading}

Two morphological systems have largely displaced simple dural sac
cross-sectional area (DSCA) thresholds for central canal assessment, on the
grounds that area measurements correlate poorly with symptoms and overlap
substantially between symptomatic and asymptomatic individuals
\cite{schizas2010qualitative}.

The \textbf{Schizas classification} grades the morphology of the dural sac on
axial T2 images by the ratio of cerebrospinal fluid to nerve rootlets. Grades~A
and~B retain visible CSF---in~A the rootlets occupy part of the sac, in~B they
fill it but remain individually discernible, giving a characteristic grainy
appearance. Grades~C and~D show no discernible rootlets and no CSF signal,
distinguished from one another by whether posterior epidural fat is preserved
(C) or obliterated (D) \cite{schizas2010qualitative}. Applied to 95 subjects
across surgical, conservatively managed and low-back-pain groups, intra- and
inter-observer agreement were substantial and moderate respectively
($\kappa = 0.65$ and $0.44$). Critically, area measurement over-diagnosed
stenosis in 35 patients and under-diagnosed it in 12 relative to the
morphological grade, and grades~C and~D predicted failure of conservative
management \cite{schizas2010qualitative}.

The \textbf{Lee central canal system} offers a simpler four-grade alternative
based on the degree of separation of the cauda equina rootlets on axial T2:
grade~0 no stenosis; grade~1 obliteration of the anterior CSF space with all
rootlets still separated; grade~2 some rootlets aggregated; grade~3 none
separated \cite{lee2011a}. Assessed by four musculoskeletal radiologists across
81 patients, it achieved inter-reader ICC of 0.730--0.953 and intra-reader
$\kappa$ of 0.863--0.900, and grades tracked cross-sectional area and
anteroposterior diameter at nearly all levels---notably better agreement than
the Schizas system, at the cost of finer morphological detail.

\subsection{Neural Foraminal Stenosis: Lee and Sartoretti}
\label{sec:lr-foraminal-grading}

The \textbf{Lee foraminal system} grades parasagittal images by the pattern of
perineural fat obliteration and the morphology of the nerve root: grade~0
normal; grade~1 fat obliteration in two opposing directions; grade~2
obliteration in all four directions without morphological change to the root;
grade~3 nerve root collapse or distortion \cite{lee2010a}. Evaluated on 576
foramina in 96 patients from L3--L4 to L5--S1 by two radiologists, inter- and
intra-observer agreement were near-perfect ($\kappa$ approximately 0.80--1.0).
This is the direct clinical precursor of the left and right foraminal narrowing
targets in the benchmark used in this thesis. The Sartoretti classification
offers a six-point sagittal alternative (grades~A--F) indexed to progressive
contact between the exiting root and each foraminal boundary.

The near-perfect agreement Lee et al.\ report deserves scepticism when
generalised. It was obtained by two experienced readers applying explicit
criteria to a curated series; the multi-reader, multi-institution agreement
figures in Section~\ref{sec:lr-reliability} are considerably lower, and it is
the latter that describe the conditions under which large training corpora are
actually annotated.

\subsection{Subarticular and Lateral Recess Stenosis}
\label{sec:lr-subarticular-grading}

Subarticular stenosis is graded on axial T2 by the proportion of the lateral
recess compromised: mild at up to one third, moderate between one and two
thirds, severe beyond two thirds, with nerve root involvement described as
touching, displacing or compressing \cite{lurie2008reliability,andreisek2014consensus}.
Compared with the disc and central canal systems this is a notably less
crisply specified instrument, resting on a visually estimated ratio rather than
a discrete morphological pattern. That imprecision shows in the reliability
data: subarticular stenosis is consistently the least reproducible of the
targets \cite{lurie2008reliability}, and it is also the target on which the
severity distribution in LumbarDISC is least skewed
(Section~\ref{sec:lr-lumbardisc}). Its difficulty is thus intrinsic, not merely
a consequence of scarce severe examples.

\subsection{Disc Herniation: Nomenclature and the MSU Classification}
\label{sec:lr-herniation-grading}

Herniation is described qualitatively by the standardised nomenclature of
bulge, protrusion, extrusion and sequestration \cite{fardon2014lumbar}, and
quantitatively by the Michigan State University (MSU) classification, which
grades size against a transverse line joining the medial edges of the facet
joints on the axial slice of maximal herniation \cite{mysliwiec2010msu}.
Grade~1 extends less than half the distance to that line, grade~2 more than
half without crossing it, grade~3 completely beyond it. The MSU system supplies
the reference standard for several automated herniation studies
\cite{zhang2023deep}, and its size~1 versus size~2/3 boundary maps onto the
conservative versus surgical decision.

\subsection{Inter- and Intra-Observer Reliability}
\label{sec:lr-reliability}

The most consequential study for present purposes is that of Lurie et al., who
assessed standardised MRI interpretation of lumbar stenosis using 58 randomly
selected studies from the SPORT trial read by four independent readers, with 20
re-read after at least a month \cite{lurie2008reliability}. Inter-reader
agreement was substantial for central stenosis ($\kappa = 0.73$) but only
moderate for foraminal stenosis ($0.58$) and nerve root impingement ($0.51$),
and poorest for subarticular stenosis ($0.49$), with marked variation between
reader pairs. Quantitative measurement fared no better in absolute terms: mean
absolute inter-reader difference in thecal sac area was 128~mm$^2$, or 13\%.
Intra-reader agreement exceeded inter-reader agreement for every feature.

Two structural findings recur across the wider reliability literature. First,
agreement is systematically worse for the lateral compartments than for the
central canal---the same ordering that later appears in automated model
performance. Second, disagreement is overwhelmingly confined to adjacent
grades: Griffith et al.\ report 84\% of intra-observer and 91\% of
inter-observer disagreement spanning a single grade
\cite{griffith2007modified}, and the same pattern recurs in automated
evaluation \cite{mcsweeney2023external,udomluck2026multitask}. Earlier consensus
efforts confirm that the underlying definitions were themselves contested: the
Delphi survey of Mamisch et al.\ \cite{mamisch2012radiologic} and the consensus
conference reported by Andreisek et al.\ \cite{andreisek2014consensus} were both
convened precisely because radiologic criteria for stenosis varied between
institutions.

\subsection{Label Noise as a Performance Ceiling}
\label{sec:lr-label-ceiling}

The reliability figures above are not a peripheral caveat; they are a hard
constraint on the interpretation of every result in this field. If four expert
readers agree on subarticular severity at $\kappa = 0.49$, then a model trained
on one reader's labels and evaluated against another's cannot exceed that
agreement, and a reported accuracy approaching unity on such a target should
prompt scrutiny of the evaluation protocol rather than admiration. The
LumbarDISC dataset description states the point explicitly, noting that
published inter-rater agreement on these grades ranges from poor to substantial
depending on anatomical location \cite{richards2026the}.

The field has responded to this in two ways. The first is to make the reference
standard more robust by construction---majority voting across a panel with
adjudication of disagreements, as in Tumko et al.\ (seven radiologists)
\cite{tumko2024a} and Su et al.\ (two clinicians with expert adjudication)
\cite{su2022automatic}. The second is to reframe the evaluation question: rather
than asking whether a model matches a single reader, asking whether its
disagreement with readers falls inside the band of disagreement between readers.
Won et al.\ furnish the cleanest example, reporting that the difference between
their two experts was statistically significant while the difference between
each expert and the model trained on that expert's labels was not
\cite{won2020spinal}. Comparatively few studies adopt either safeguard, and
Section~\ref{sec:lr-reporting} returns to this as a systemic weakness.

\subsection{The Ordinal Structure of the Label Space}
\label{sec:lr-ordinality}

Severity labels are ordinal, not categorical: \emph{normal}, \emph{mild},
\emph{moderate} and \emph{severe} stand in a fixed order, and the distance
between them is neither uniform nor clinically symmetric. Treating them as
unordered classes under categorical cross-entropy---as the majority of the
reviewed work does---discards that structure, penalising a
severe-graded-as-normal error exactly as heavily as a moderate-graded-as-severe
one.

Two empirical observations follow. First, errors concentrate overwhelmingly at
adjacent grade boundaries. McSweeney et al.\ report that SpineNet disagreed with
human raters on 20.83\% of discs but that only 0.85\% of discs differed by more
than one grade \cite{mcsweeney2023external}; Udomluck et al.\ observe the same
concentration \cite{udomluck2026multitask}; Ishimoto et al.\ find exact-grade
agreement of 65.7\% rising to 94.1\% when collapsed to a severe-versus-rest
dichotomy \cite{ishimoto2020could}. Much of the apparent error in ordinal
grading is therefore boundary error rather than gross misclassification.

Second, and less expectedly, explicitly ordinal objectives have not reliably
improved on this. Niemeyer et al.\ tested soft-kappa loss, ordinal
cross-entropy and regression losses against plain cross-entropy for Pfirrmann
grading on 1{,}599 patients and 7{,}948 discs, together with class pooling and
class-weighted losses, and found that none improved overall classification
performance; the default cross-entropy classifier reached $\kappa = 0.92$ with
predictions within one grade of ground truth in 99.2\% of cases
\cite{niemeyer2021a}. Patil et al.\ likewise found performance strongest at the
extremes of the Pfirrmann scale and on the binary low-versus-high distinction,
with intermediate grades remaining the difficulty \cite{patil2026mribased}. The
ordinal structure of the problem is thus real and measurable, but the question
of how to exploit it is genuinely open---a point that informs the loss functions
adopted in Chapter~3.

% ===========================================================================
\section{Evolution of Automated Lumbar MRI Analysis}
\label{sec:lr-evolution}

\subsection{Handcrafted Features and Classical Machine Learning}
\label{sec:lr-classical}

Before deep learning, automated analysis of lumbar MRI depended on features
designed by hand. Investigators specified mathematical descriptors of local
appearance and passed them to a conventional classifier. Ghosh et al.\ combined
Histogram of Oriented Gradients (HOG) descriptors with a support vector machine
and a set of geometric heuristics, reporting 99\% disc localisation accuracy
across 53 clinical cases and 318 discs, and---unusually for the period---emitting
a tight bounding box for each disc rather than a single interior point, so that
a downstream system could operate on the region directly
\cite{ghosh2012a}. Athertya and Kumar took a comparable approach to vertebral
degeneration, extracting local binary patterns, Hu moments and gray-level
co-occurrence matrix (GLCM) statistics to classify Modic changes, endplate
defects and focal changes, reaching 85.91\% accuracy for the extent of Modic
change on 100 patients \cite{athertya2021classification}. Related efforts applied
hybrid classical models to disc abnormality detection
\cite{unal2015detection,castromateos2016intervertebral} and semi-automatic
computer-aided classification of degenerative disease \cite{ruizespaa2015semiautomatic},
while Huber et al.\ pursued texture analysis with decision trees as a
quantitative alternative to visual stenosis grading \cite{huber2019qualitative}.
He et al.\ applied a synchronised superpixel representation to define consistent
regions of interest across levels before classifying neural foraminal stenosis,
an early attack on one of the exact conditions graded in this thesis
\cite{he2016automated}.

These systems established feasibility but generalised poorly. Handcrafted
descriptors are sensitive to scanner hardware, field strength, slice thickness
and patient positioning, and the anatomical variability of the lumbar spine
defeats fixed geometric assumptions. One finding from this era nevertheless
retains value: Athertya and Kumar's feature-sensitivity analysis showed that
GLCM entropy alone sufficed to detect focal changes, while endplate-defect
classification depended almost entirely on the first two Hu moments
\cite{athertya2021classification}. Such explicit attribution of predictive power
to individual features is something the deep models that displaced these methods
recovered only later, and imperfectly, through the interpretability techniques of
Section~\ref{sec:lr-interpretability}.

\subsection{The First End-to-End Deep Learning Systems}
\label{sec:lr-first-dl}

Convolutional neural networks removed the need for explicit feature design by
learning hierarchical representations directly from pixels \cite{he2016deep}. In
lumbar imaging the decisive contribution was SpineNet
\cite{jamaludin2017spinenet}. Trained on 2{,}009 patients and 12{,}018 disc
volumes from sagittal T2 images, it predicted six radiological scores
simultaneously---Pfirrmann grade, disc narrowing, upper and lower endplate
defects, and upper and lower marrow changes---from a shared convolutional trunk
that branched into task-specific heads. Two design decisions proved durable.
First, multi-task learning let one architecture serve several targets, reducing
the free parameters relative to training separate networks and exploiting the
correlation between co-occurring pathologies. Second, the system required only
disc-level class labels: no voxel-wise segmentation and no slice-level
localisation, with the evidence hotspots emerging implicitly from classification
training. Reported performance was state-of-the-art and near-human across all
six scores.

The same group demonstrated shortly afterwards that longitudinal clinical scans
could serve as free supervision. A Siamese network was pre-trained with a
contrastive loss on whether two scans belonged to the same patient---information
available from DICOM metadata without knowing the patient's identity---together
with an auxiliary vertebral-level prediction task, then fine-tuned for
degeneration grading. On 1{,}016 subjects, 423 with follow-up imaging, the
pre-trained network outperformed training from scratch and reached equivalent
performance from far fewer annotated examples \cite{jamaludin2017selfsupervised}.
This anticipated by several years the self-supervised pretraining strategies now
routine in medical imaging, and it is the direct ancestor of the contrastive
approach discussed in Section~\ref{sec:lr-contrastive}.

SpineNet's principal limitation was structural rather than architectural: it
consumed preprocessed sagittal disc volumes, and could therefore not assess
pathologies that require the transverse plane. SpineNetV2 extended coverage to
eleven radiological features and moved to a multi-stage design, using a U-Net to
detect and label vertebral bodies and discs before a ResNet classifier graded
each \cite{windsor2022spinenetv2}. The reliance on sagittal input nonetheless
persisted, and Section~\ref{sec:lr-generalization} shows this to be the
consistent locus of its external-validation failures.

\subsection{From Binary Detection to Multi-Level, Multi-Pathology Grading}
\label{sec:lr-task-evolution}

The task definition itself has evolved, and the trajectory matters more than any
individual model. Early deep learning work asked binary questions: is stenosis
present or absent \cite{altun2023lssvgg16}; is a herniation visible in this image
\cite{tsai2021lumbar}. Such formulations are tractable but clinically thin,
because management depends on severity and on which level is affected.

Three shifts followed. The first was to \emph{multi-class severity}: Won et al.\
graded stenosis status on 542 L4--5 axial images against two independent experts
\cite{won2020spinal}; Niemeyer et al.\ predicted the full five-point Pfirrmann
scale across 7{,}948 discs \cite{niemeyer2021a}. The second was to
\emph{multi-pathology} output, in which one model reports several findings at
once: Lehnen et al.\ evaluated disc herniation, bulging, canal stenosis,
nerve-root compression and spondylolisthesis across 888 lumbar segments from a
single CNN \cite{lehnen2021detection}; Su et al.\ graded herniation, central
canal stenosis and nerve-root compression from a shared ResNet-50 trunk with
three classification heads \cite{su2022automatic}. The third was to
\emph{per-level, per-side} prediction, in which the unit of analysis is not the
study or the image but the individual disc level and, for the lateral
compartments, the individual side. Hallinan et al.\ exemplify the mature form,
grading central canal, lateral recess and neural foraminal stenosis at each
level with the sequence appropriate to each \cite{hallinan2021deep}. Single-stage
detectors have since been adapted to deliver the same output more cheaply: Wang
et al.\ modified a YOLOv5 network to detect regions of interest and grade
central, lateral recess and foraminal stenosis simultaneously
\cite{wang2025a}, and a related design grades Pfirrmann degeneration, herniation
and high-intensity zones from sagittal and axial input in one pass
\cite{wang2025automated}. Minin et al.\ pursued the same objective for Pfirrmann
grading with an explicitly open model, motivated by the inaccessibility of
proprietary graders \cite{minin2025spinescan}.

This third formulation is precisely that of the benchmark used in this thesis,
which requires twenty-five graded outputs per study: five conditions across five
disc levels \cite{rsna2024rsna,richards2026the}. The progression is therefore not
merely historical. It explains why performance figures from the binary era
cannot be compared with contemporary results, and why a model that classifies a
whole image is solving a different and substantially easier problem than one
that must attribute severity to a specific anatomical location.

\subsection{Chronological Synthesis}
\label{sec:lr-chronology}

Four phases are discernible. Until roughly 2016 the field was dominated by
handcrafted features and classical classifiers, with localisation as the central
problem \cite{ghosh2012a,ghosh2014supervised}. From 2017 to 2020, CNNs displaced
these methods and multi-task grading became feasible
\cite{jamaludin2017spinenet,lu2018deepspine,han2018spinegan}. From 2021 to 2023
the emphasis moved to clinical credibility: external validation
\cite{mcsweeney2023external}, interpretability \cite{bharadwaj2023deep,gao2022automatic},
reader-assistance evaluation \cite{lim2022improved} and multi-centre testing
\cite{yi2023deep,su2022automatic}. From 2024 onward, the release of a large
public benchmark \cite{richards2026the} redirected effort toward anatomically
localised, multi-view severity grading, alongside the arrival of transformer and
vision-language architectures
\cite{batra2025mscan,nguyen2026crossspine,islamsk2026an,zeng2025gpt4lfs}.

The through-line is a steady increase in the specificity of what is being
predicted, accompanied---until recently---by a decrease in the rigour with which
it is validated. Both systematic reviews of the period make the same complaint:
Compte et al.\ found that algorithms performed markedly worse in replication or
external validation than in their development studies \cite{compte2023are}, and
Wang et al.\ concluded that despite pooled sensitivity of 0.84 and specificity of
0.87 across twelve studies and 15{,}044 patients, no model was reliable enough
for clinical practice \cite{wang2024machine}. Section~\ref{sec:lr-validation}
takes up this tension directly.

% ===========================================================================
\section{Anatomical Localization as a Design Principle}
\label{sec:lr-localization}

\subsection{Disc and Vertebra Detection, Labelling, and Level Assignment}
\label{sec:lr-detection}

Before severity can be attributed to L4--L5, the system must know which structure
is L4--L5. Level assignment is thus a precondition of per-level grading, and it
has been treated as a problem in its own right since the earliest work
\cite{ghosh2012a,lootus2015automated}. Pan et al.\ illustrate the mature
solution: a Faster R-CNN locates vertebral bodies L1--S1 on sagittal images, the
geometric relationship between sagittal and axial acquisitions identifies the
corresponding disc in each axial slice, and the disc region of interest is then
extracted---each of the three stages reported at 100\% accuracy, with the residual
difficulty falling entirely on the subsequent classification
\cite{pan2021automatically}. Detection has, in short, largely been solved;
grading has not.

\subsection{Semantic and Instance Segmentation of Lumbar Structures}
\label{sec:lr-segmentation}

Dense segmentation provides a stronger anatomical substrate than bounding boxes.
The U-Net architecture \cite{ronneberger2015unet} and its descendants dominate,
as a dedicated systematic review of deep-learning-assisted disc segmentation
confirms \cite{wang2024deep}.
Ghosh and Chaudhary segmented vertebrae, disc, dural sac and background jointly
using cascaded random forests over image and sparsely sampled context features,
achieving Dice 0.87 and 0.84 for dural sac and disc across 212 clinical cases
\cite{ghosh2014supervised}. Han et al.\ introduced adversarial training for
multi-structure segmentation \cite{han2018spinegan}, S\'aenz-Gamboa et al.\
established CNN semantic segmentation of the structural elements surrounding the
spinal cord \cite{senzgamboa2021automatic}, and Al-Kafri et al.\ applied
SegNet to boundary delineation for stenosis assessment across 515 symptomatic
patients, contributing two purpose-built metrics---confidence and
consistency---for assessing the quality of the ground truth itself rather than
only the model \cite{alkafri2019boundary}.

Subsequent work has largely consisted of architectural refinement. Wang et al.\
added multilevel feature-map attention and residual blocks to an Attention U-Net
\cite{wang2022automatic}; Ahmed et al.\ targeted class imbalance between
anatomical structures with a custom combined loss and modified upsampling
\cite{ahmed2024pioneering}; Silveira et al.\ added an Inception module for
multi-scale extraction and a dual-output mechanism, reporting mIoU 0.8974 and
F1 0.9444 on SPIDER and---notably---outperforming TransUNet
\cite{silveira2025automated}. M\"oller et al.\ extended coverage to the whole
spine with SPINEPS, segmenting fourteen structures including the posterior
elements, and established a useful general result: segmenting semantically first
and instance-wise second outperformed training directly on instance segmentation,
beating an nnU-Net baseline \cite{isensee2021nnunet} on every structure and metric
\cite{mller2025spinepsautomatic}. Chen et al.\ merged convolutional and attention
pathways in parallel with a novel position embedding, outperforming fifteen
competing models \cite{chen2024symtc}.

\subsection{Public Segmentation Resources and Benchmarks}
\label{sec:lr-segmentation-data}

Progress in this subfield has been unusually well served by shared data. The
SPIDER dataset provides 447 sagittal T1 and T2 series from 218 patients across
four hospitals, with reference segmentations of vertebrae, discs and spinal
canal plus per-level radiological gradings, together with reference performance
for a baseline algorithm and nnU-Net and a continuous public challenge on a
sequestered test split \cite{vandergraaf2024lumbar}. Several of the segmentation
models above are trained and benchmarked on it
\cite{ahmed2024pioneering,silveira2025automated,patil2026mribased}, which makes
their results mutually comparable in a way that is rare elsewhere in this
literature.

\subsection{Quantitative Morphometry}
\label{sec:lr-morphometry}

Segmentation enables measurement, and measurement offers an escape from ordinal
subjectivity. The dural sac cross-sectional area is the canonical example.
Ghobrial and Roth compared U-Net, Attention U-Net and MultiResUNet for automated
DSCA quantification on 683 scans with 50 held back for external validation;
MultiResUNet performed best, with Pearson correlation 0.9917 and mean absolute
error 23.70~mm$^2$, holding up externally at 99.95\% accuracy and F1 0.9393
\cite{ghobrial2025deep}. Zheng et al.\ pursued the same objective for disc
degeneration, segmenting degeneration-related regions with a Swin-transformer
network and computing signal-intensity and geometric features, then proposing a
quantitative criterion to replace subjective Pfirrmann grading on the basis of
strong correlation with grade across a large population \cite{zheng2022deep}.
Related work measures the lordosis angle from segmentation masks
\cite{ji2026mdvmunet}.

A caution is warranted. Ghobrial and Roth's model relies on T1 weighting alone
and a modest dataset, as the authors acknowledge \cite{ghobrial2025deep}, and
Schizas' original finding---that area measurement over- and under-diagnosed
stenosis relative to morphological grade in a substantial minority of patients
\cite{schizas2010qualitative}---applies equally whether the area is traced by
hand or by network. Automating a measurement does not repair a weak correlation
between that measurement and the clinical state.

\subsection{Evidence that Localisation Precedes Classification}
\label{sec:lr-localize-first}

The strongest recurring methodological finding in this literature is that
explicit anatomical localisation before classification improves grading. The
argument was made structurally by DeepSPINE, which chained U-Net vertebral
segmentation and disc-level designation to a multi-task grading network across
4{,}075 patients and more than 22{,}000 disc-level annotations
\cite{lu2018deepspine}, and it has been corroborated repeatedly since.

Four lines of evidence converge. Bharadwaj et al.\ segmented dural sac and disc
with V-Net, localised facet and foramen by geometric rule, and only then
classified, reaching Dice 0.93 and 0.94 with foramen and facet localisation IoU
0.72 and 0.83 \cite{bharadwaj2023deep}. \v{S}u\v{s}ter\v{s}i\v{c} et al.\ showed
that segmenting, cropping and enhancing the region of interest before
classification produced 0.87 and 0.91 accuracy on axial and sagittal views for a
multi-class herniation problem \cite{sustersic2022a}. Kowlagi et al.\ made the
comparative case most directly, arguing that a well-tuned three-stage pipeline of
segmentation, localisation and classification allows convolutional networks to
outperform transformer approaches on Pfirrmann grading in a population cohort
\cite{kowlagi2022a}. Acharya and Kansakar tested the principle in its strongest
form, hypothesising that anatomical focus rather than architecture is the binding
constraint, and grading severity from a single localised region of interest per
disc \cite{acharya2026disccentric}.

The implication for the present work is direct. If localisation dominates
architecture in determining performance, then a comparison between backbones is
only meaningful when both receive identically localised input---otherwise the
comparison measures preprocessing rather than representation. This is a
methodological requirement that Section~\ref{sec:lr-cnn-vs-transformer} shows
much of the existing comparative literature fails to satisfy.

\subsection{Reduced-Annotation and Unsupervised Alternatives}
\label{sec:lr-weak-supervision}

Dense annotation is expensive, and several strands of work reduce the
requirement. Kuang et al.\ eliminated manual labels entirely for vertebral
segmentation: sub-optimal masks from a rule-based method were combined through a
voting mechanism to supervise CNN training, with a regional strategy restricting
attention to a specific vertebral region, validated across 243 volunteers
\cite{kuang2020mrisegflow}. SpineNet's evidence hotspots represent the
complementary weakly supervised route, obtaining localisation from
classification labels alone \cite{jamaludin2017spinenet}, while the SPIDER
annotation protocol adopted a human-in-the-loop compromise, training on a small
subset to produce initial masks that were then corrected and fed back
\cite{vandergraaf2024lumbar}. Self-supervised pretraining on unlabelled or
metadata-labelled data \cite{jamaludin2017selfsupervised} completes the picture.
Together these establish that the annotation burden, long cited as the principal
obstacle to clinical deployment, is substantially negotiable.

% ===========================================================================
\section{Backbone Architectures for Severity Classification}
\label{sec:lr-backbones}

\subsection{Convolutional Networks}
\label{sec:lr-cnns}

Convolutional networks exploit two properties well matched to imaging: local
connectivity, since a filter responds to a neighbourhood rather than the whole
field; and translation equivariance, since the same filter is applied at every
position. Stacking convolutions builds a hierarchy in which early layers respond
to edges and texture and later layers to composite anatomical structure.

\textbf{ResNet} addressed the degradation encountered when networks are made
deeper, introducing residual blocks whose skip connections allow a layer's input
to bypass subsequent convolutions and be added to their output, providing an
unimpeded path for gradients and permitting stable training at fifty layers and
beyond \cite{he2016deep}. In lumbar imaging ResNet is ubiquitous: as the
classification head in multi-task grading \cite{su2022automatic}, as the
detection backbone in staged pipelines \cite{zhang2023deep}, as the encoder in
segmentation networks \cite{batra2025mscan}, and as the classifier in
SpineNetV2 \cite{windsor2022spinenetv2}.

\textbf{EfficientNet} attacked a different problem---how to allocate a fixed
computational budget---by scaling depth, width and input resolution together
under fixed coefficients derived from a neural architecture search, rather than
scaling one dimension in isolation \cite{tan2019efficientnet}. The resulting
models reach comparable accuracy at substantially lower parameter and FLOP cost,
which matters for the deployment scenarios that motivate this thesis.
EfficientNet appears as the classifier in Kowlagi et al.'s Pfirrmann pipeline
\cite{kowlagi2022a} and in the final stage of M-SCAN \cite{batra2025mscan}.
\textbf{ConvNeXt} represents the most recent convolutional line, modernising the
CNN with design choices borrowed from transformers while retaining convolution
\cite{liu2022convnet}; Udomluck et al.\ include ConvNeXt-Tiny in their feature
extractor comparison and find it stable across classifiers
\cite{udomluck2026multitask}.

The architectural limitation of convolution is a bounded receptive field. Deep
stacks aggregate local responses into progressively broader ones, but the
network never explicitly models a relationship between two distant regions---for
instance, between the L1--L2 and L5--S1 levels, whose degenerative states are
biomechanically coupled.

\subsection{Vision Transformers and the Attention Bottleneck}
\label{sec:lr-vit}

Transformers dispense with convolution and rely on self-attention, in which
every element of the input is weighted against every other regardless of
distance \cite{vaswani2017attention}. The Vision Transformer applies this by
partitioning an image into fixed patches, flattening them and treating the
result as a sequence \cite{dosovitskiy2021image}. Global context is available
from the first layer, and there is no locality prior to overcome---but also no
locality prior to exploit, which is why vision transformers are markedly more
data-hungry than convolutional networks of comparable capacity.

The practical obstacle is cost. Global self-attention scales quadratically with
the number of patches, hence quadratically with image area. Clinical MRI demands
high spatial resolution to resolve structures a few millimetres across, so a
naive vision transformer over full-resolution lumbar images is computationally
prohibitive.

\subsection{The Swin Transformer}
\label{sec:lr-swin}

The Swin Transformer resolves this bottleneck through two mechanisms
\cite{liu2021swin}. \emph{Windowed self-attention} restricts attention to
non-overlapping local windows, making cost linear rather than quadratic in image
size; to restore cross-window communication, the window partition is
\emph{shifted} by half a window between consecutive layers, so information
propagates across boundaries over successive blocks. \emph{Hierarchical feature
maps} are produced by patch-merging layers that progressively reduce spatial
resolution while increasing channel depth, yielding a pyramid analogous to a
CNN's and making the backbone usable for dense prediction as well as
classification.

For lumbar grading the appeal is specific. Severity depends simultaneously on
fine local detail---the texture of a nucleus pulposus, the obliteration of
perineural fat---and on regional context, since a rootlet configuration is
interpretable only relative to the surrounding dural sac and the levels above
and below. A hierarchical, windowed attention mechanism is a plausible fit to
that structure, which is the substantive reason for selecting it as the
transformer representative in this thesis. Swin-derived architectures have been
applied to spinal segmentation, and Zheng et al.\ incorporate Swin components in
their quantitative degeneration network \cite{zheng2022deep}.

\subsection{Hybrid CNN--Transformer Architectures}
\label{sec:lr-hybrids}

A third position holds that the two families are complementary. SymTC merges
convolutional and transformer layers in a \emph{parallel dual-path} arrangement
rather than stacking them sequentially, and integrates a position embedding into
the self-attention module to improve spatial awareness; benchmarked against
fifteen existing models on both a private dataset and a released synthetic one,
it performed best for both vertebral bones and discs \cite{chen2024symtc}. Zheng
et al.\ similarly combine convolutional and Swin components
\cite{zheng2022deep}, and Yi et al.\ pair a 3D ResNet18 backbone with transformer
components across three hospitals \cite{yi2023deep}. Baur et al.\ explore a
different pairing, segmenting discs with a U-Net, reconstructing each as a 3D
point cloud and grading it with a graph neural network; performance was uneven
across levels---F1 0.71 at L2/3 against 0.46 overall---which the authors present
as a foundation requiring refinement rather than a competitive result
\cite{baur2024automated}. Hybrids complicate the
clean comparison this thesis undertakes, but they indicate where the field
expects the answer to lie: not in the wholesale replacement of convolution, but
in supplying it with a mechanism for long-range interaction.

\subsection{Published CNN-versus-Transformer Comparisons and Their Limits}
\label{sec:lr-cnn-vs-transformer}

Several studies bear directly on the architectural question, and their results
do not favour transformers as strongly as the general computer vision literature
might suggest.

Kowlagi et al.\ address it explicitly. Observing that recent gains had been
attributed to transformer architectures, they show that a well-tuned three-stage
convolutional pipeline outperforms the state-of-the-art transformer approaches
on Pfirrmann grading, evaluated on the Northern Finland Birth Cohort 1966---a
population cohort rather than a clinical convenience sample---with an ablation
study and subgroup breakdown, and public code \cite{kowlagi2022a}. Their
conclusion is that pipeline design and tuning outweigh backbone choice.

Lee et al.\ provide the closest published analogue to the present study,
developing a CNN-based and a transformer-based model for lumbar stenosis
classification on 564 MRIs with a 100-study external test set, labelled by four
radiologists with expert consensus as reference and read by eight participants
across four experience levels. Both models exceeded 94\% detection recall; on
dichotomised classification the CNN, transformer and human participants achieved
kappas of 0.99/0.99/0.97--0.98 for central canal, 0.98/0.94/0.81--0.94 for
lateral recesses and 0.98/0.95/0.91--0.95 for neural foramina. The CNN was
superior overall, the transformer similar to superior against clinicians
\cite{lee2026deep}.

Udomluck et al.\ compared VGG19, ConvNeXt-Tiny and DINOv2 features with classical
classifiers under external validation, finding conventional VGG19 features
strongest (accuracy 0.9556 with logistic regression) and DINOv2 features
generalising worst, dropping to 0.6222 with LightGBM
\cite{udomluck2026multitask}. Patil et al.\ reached a still more pointed
conclusion through ablation: radiomic features carried essentially the entire
signal for Pfirrmann grading, and frozen ImageNet-pretrained ResNet50 features
added no measurable discriminative value \cite{patil2026mribased}. Silveira et
al.\ likewise found their enhanced U-Net outperformed TransUNet on segmentation
\cite{silveira2025automated}.

Three methodological limitations recur across this evidence. First, the
compared architectures are rarely given identical preprocessing, localisation
and training budgets, so differences may reflect tuning effort rather than
representational capacity---exactly the confound Kowlagi et al.\ identify
\cite{kowlagi2022a}. Second, the input configuration is usually held fixed,
leaving unexamined the interaction between architecture and available context,
even though a transformer's advantage should plausibly be largest where global
context must substitute for missing views. Third, several comparisons rely on
frozen pretrained features rather than fine-tuned backbones
\cite{udomluck2026multitask,patil2026mribased}, which disadvantages transformers
disproportionately given their weaker inductive biases. These three gaps
motivate the experimental design of this thesis.

\subsection{Transfer Learning, Pretraining, and Self-Supervision}
\label{sec:lr-pretraining}

Medical datasets are small by the standards these architectures were designed
for, so initialisation matters. ImageNet pretraining is near-universal, though
Patil et al.'s ablation is a caution that its benefit is not automatic
\cite{patil2026mribased}. Domain-specific pretraining is better supported:
Jamaludin et al.\ obtained equivalent performance from far fewer labels through
longitudinal self-supervision \cite{jamaludin2017selfsupervised}, and Acharya and
Kansakar used contrastive pretraining on disc-centred crops to reduce sensitivity
to irrelevant appearance variation \cite{acharya2026disccentric,chen2020simple}.
Sequential transfer between related targets---initialising a subarticular model
from spinal canal weights---has also been reported to help on smaller datasets, a
strategy of particular relevance to transformers given their lack of
convolutional inductive bias.

% ===========================================================================
\section{Exploiting Spatial Context: Slices, Volumes, and Views}
\label{sec:lr-context}

\subsection{2D, 2.5D, and 3D Input Representations}
\label{sec:lr-dimensionality}

An MRI series is a volume, but severity is assigned per level, so a
representation must be chosen. \emph{2D} models classify individual slices and
are cheap but discard through-plane context; since a disc spans several slices
and a radiologist integrates across them, a single slice may not contain the
evidence. \emph{3D} models consume the volume directly \cite{yi2023deep} at
considerable memory cost---Acharya and Kansakar note that a standard
high-resolution series contains millions of voxels, making full-volume
processing impractical for many architectures \cite{acharya2026disccentric}.

\emph{2.5D} representations occupy the middle ground and have become the
default in this domain. A small stack of adjacent slices is treated as channels
or as a short sequence, capturing local through-plane context at roughly 2D
cost. SpineNet's input was a stack of nine slices per disc
\cite{jamaludin2017spinenet}; M-SCAN describes its approach explicitly as 2.5D,
selecting the three axial slices closest to each level and modelling them with
recurrent layers \cite{batra2025mscan}; the leading challenge solutions combined
2D CNNs with sequence models or attention-based multiple-instance learning over
the slice dimension.

\subsection{Cross-Sequence and Cross-Plane Fusion}
\label{sec:lr-fusion}

Where multiple sequences are available, they must be combined. The simplest
approach concatenates them as channels, which presumes spatial registration that
sagittal and axial acquisitions do not satisfy. More considered strategies
process each sequence separately and fuse at the feature level.

CrossSpine is the most explicit treatment. Nguyen et al.\ identify two failings
of baselines---single-stream designs cannot integrate complementary information
across sequences, and uniform treatment of all discs ignores level-specific
anatomy---and address the first with a cross-sequence attention mechanism that
fuses features at multiple spatial scales, letting each sequence attend to all
others, followed by a weighting mechanism that prioritises the most
diagnostically informative sequence. They report Macro F1 improved by more than
125\%, Mean AUPRC by 99\% and Mean AUROC by 36\% over their baseline
\cite{nguyen2026crossspine}. M-SCAN applies cross-attention to align sagittal and
axial feature maps before classification \cite{batra2025mscan}, and Park et al.\
pursue transformer-based multimodal fusion on the same benchmark
\cite{park2025multiview}. Zeng et al.\ extend fusion beyond imaging entirely,
combining ConvNeXt image features with GPT-4o-generated textual descriptions
encoded by RoBERTa through a Mamba fusion stage \cite{zeng2025gpt4lfs}.

\subsection{Geometric Correspondence Between Planes}
\label{sec:lr-geometry}

Fusing sagittal and axial data requires knowing which axial slice corresponds to
which sagittal level---a geometric problem solved through DICOM spatial
metadata rather than learned. Pan et al.\ derive the correspondence from the
geometric relationship between the two acquisitions, reporting 100\% accuracy in
identifying the disc in each axial slice \cite{pan2021automatically}. M-SCAN
projects 2D sagittal stenosis coordinates into 3D using patient orientation data
from the DICOM headers, then selects the three nearest axial slices for each of
the five levels \cite{batra2025mscan}.

This step is a hidden cost of multi-sequence modelling, and one that bears
directly on the trade-off examined here. It adds a preprocessing dependency, a
failure mode when metadata is absent or inconsistent, and an engineering burden
that a single-sequence pipeline simply does not incur.

\subsection{Multi-Sequence versus Single-Sequence Input}
\label{sec:lr-multiseq}

The clinical standard is unambiguous that all three sequences contribute
\cite{richards2026the,hallinan2021deep}, and the strongest reported results on
the RSNA benchmark use all three \cite{batra2025mscan}. The countervailing
consideration is cost: multi-sequence processing demands greater memory and
compute, cross-plane registration, and longer inference, and the resulting
infrastructure requirements exceed what many hospital environments provide.

The evidence on whether reduced input suffices is genuinely mixed, which is why
the question remains open. On one side, external validation of sagittal-only
SpineNetV2 found agreement dropping specifically for foraminal stenosis and
herniation---the pathologies requiring transverse context
\cite{nigru2024external}---and Wu et al.\ found ordinal Pfirrmann grading
degrading in upper lumbar discs while binary detection held up
\cite{wu2026external}. On the other, Acharya and Kansakar grade severity from
\emph{sagittal T2 alone}, reaching 78.1\% balanced accuracy with a
severe-to-normal misclassification rate of 2.13\%, on the argument that
anatomical focus rather than input breadth is the binding constraint
\cite{acharya2026disccentric}. Kowlagi et al.\ likewise operate on sagittal input
and outperform transformer baselines \cite{kowlagi2022a}.

Two observations sharpen the question. First, the comparison has not been made
under controlled conditions: the sagittal-only and multi-sequence results above
come from different architectures, datasets, localisation strategies and
evaluation protocols, so the contribution of sequence configuration cannot be
isolated. Second, the expected penalty is target-dependent---severe for
foraminal and subarticular grading, mild for central canal---so a single
aggregate accuracy figure would obscure precisely the structure of interest.
Both observations inform the design adopted in Chapter~3.

\subsection{Robustness to Missing or Incomplete Sequences}
\label{sec:lr-missing}

Clinical data is incomplete: sequences are omitted, truncated or corrupted by
artifact. The LumbarDISC inclusion criteria required all three sequences and
excluded studies with substantial artifact \cite{richards2026the}, so a model
trained on it has not been exposed to the degraded conditions of routine
practice. A model that requires three sequences fails outright when one is
absent, whereas a single-sequence model is unaffected---an argument for reduced
input configurations that is about robustness rather than efficiency, and one
that the literature on missing-modality learning addresses only in passing for
this application.

% ===========================================================================
\section{Class Imbalance and Ordinal Severity Modelling}
\label{sec:lr-imbalance}

\subsection{Severity Distributions in Lumbar MRI Datasets}
\label{sec:lr-distributions}

Class imbalance in this domain is not an artefact of sampling but a property of
the population: in any unselected clinical cohort, most disc levels are normal or
mildly degenerate, and severe compression is uncommon. The RSNA LumbarDISC
distributions make the scale explicit \cite{richards2026the}. Of 13{,}474 spinal
canal stenosis grades, 85.4\% are normal or mild, 8.8\% moderate and 5.9\%
severe. Of 26{,}919 neural foraminal grades, roughly 78\% are normal or mild,
18\% moderate and 4.5\% severe. Subarticular grading is the least skewed of the
three, with roughly 69\% normal or mild, 19\% moderate and 11\% severe across
26{,}285 grades.

The same pattern recurs elsewhere. Liawrungrueang et al.\ report a Pfirrmann
distribution of 220 grade~I, 530 grade~II, 170 grade~III, 160 grade~IV and just
20 grade~V across 1{,}000 images \cite{liawrungrueang2023automatic}; Niemeyer et
al.\ describe an approximately Gaussian grade distribution with grades~1 and~5
under-represented \cite{niemeyer2021a}. Two consequences follow. First, a model
trained under unweighted cross-entropy can achieve roughly 85\% accuracy on
spinal canal stenosis by predicting \emph{normal/mild} unconditionally, so
aggregate accuracy is close to uninformative on this task. Second, the ordering
of imbalance across targets is the inverse of the ordering of difficulty:
subarticular stenosis has the most balanced label distribution yet the poorest
human reliability \cite{lurie2008reliability}, confirming that its difficulty is
intrinsic rather than a consequence of scarce positive examples.

\subsection{Data-Level Strategies}
\label{sec:lr-data-level}

The most common data-level intervention is weighted random sampling, in which
the probability of drawing a minority-class example is inflated so that each
minibatch approaches a uniform severity distribution and gradients are not
dominated by the majority class. Aggressive augmentation of minority classes
complements this, since with few severe examples a network may memorise them
rather than learn the morphology of severe compression. Reported techniques
include contrast-limited adaptive histogram equalisation, random scaling, affine
rotation and elastic deformation, alongside mixing methods.

Tsai et al.\ provide the clearest single demonstration of augmentation's value
in the small-data regime, treating it as the subject of the study rather than an
implementation detail: mean average precision reached 92.4\% at 550 images with
augmentation, and augmentation was decisive in preventing both under- and
over-fitting \cite{tsai2021lumbar}. Ahmed et al.\ attack a related form of
imbalance---between anatomical structures rather than severity classes---through
preprocessing that rectifies class inconsistencies before training
\cite{ahmed2024pioneering}.

\subsection{Loss-Level Strategies}
\label{sec:lr-loss-level}

Class-weighted cross-entropy multiplies each class's contribution to the loss by
a scalar, forcing the optimiser to penalise minority-class errors more heavily.
The RSNA challenge encoded this directly into its evaluation metric, assigning
weights of 1, 2 and 4 to normal/mild, moderate and severe respectively, so that
aligning the training objective with the evaluation criterion is a matter of
adopting the same coefficients.

Focal loss goes further, introducing a modulating factor that down-weights
well-classified examples dynamically and concentrates gradient on hard,
borderline cases \cite{lin2017focal}. This is well matched to the present
problem, where the abundant normal cases are mostly easy and the difficulty
concentrates at the mild/moderate and moderate/severe boundaries---precisely
where human readers also disagree. Acharya and Kansakar combine weighted focal
loss with contrastive pretraining \cite{acharya2026disccentric}, and Islam~Sk et
al.\ propose an adaptive PID-Tversky loss that borrows feedback control to adjust
training penalties dynamically toward under-segmented minority instances
\cite{islamsk2026an}.

\subsection{Ordinal-Aware Objectives and Whether They Help}
\label{sec:lr-ordinal-losses}

Given the ordinal label structure of Section~\ref{sec:lr-ordinality}, it is
natural to expect ordinal-aware objectives to outperform categorical ones. The
evidence does not support this expectation as strongly as the intuition
suggests.

Niemeyer et al.\ conducted the most systematic test available, evaluating soft
kappa loss, ordinal cross-entropy and regression losses against standard
cross-entropy for Pfirrmann grading on 1{,}599 patients and 7{,}948 discs, with
extensive hyperparameter optimisation, and additionally trying class pooling and
class-weighted losses to counter imbalance. None of these improved classification
performance overall. The default cross-entropy classifier reached Cohen
$\kappa = 0.92$, average sensitivity 90.2\% and precision 92.5\%, with
predictions within one grade of ground truth in 99.2\% of cases and a mean
absolute error of 0.08 grades \cite{niemeyer2021a}. Patil et al.\ list ordinal
regression among the supplementary analyses they ran, and still report that
intermediate-grade discrimination remains the outstanding difficulty
\cite{patil2026mribased}.

Two readings are possible. Either the ordinal structure is already implicitly
captured---a network trained with cross-entropy on visually ordered classes will
naturally confuse adjacent grades more often than distant ones---or the specific
formulations tested were inadequate. The literature does not currently
distinguish these, which is why this thesis treats the choice of objective as an
empirical question rather than a settled one.

\subsection{Contrastive and Representation-Learning Approaches}
\label{sec:lr-contrastive}

A third strategy attacks imbalance in the representation rather than in the
sampling or the loss. Contrastive learning pulls examples of similar severity
together in the latent space and pushes dissimilar ones apart before any
classifier is fitted \cite{chen2020simple}, so the representation is shaped by
the structure of the problem rather than by class frequency.

Acharya and Kansakar apply this directly, combining contrastive pretraining with
disc-level fine-tuning over a single anatomically localised region of interest
per disc, on the reasoning that this forces attention onto intrinsic disc
features and reduces sensitivity to irrelevant appearance variation. An auxiliary
regression task handles disc localisation and a weighted focal loss addresses the
residual imbalance. The result---78.1\% balanced accuracy with the
severe-to-normal misclassification rate reduced to 2.13\% relative to supervised
training from scratch---is notable less for the headline accuracy than for what
it optimises \cite{acharya2026disccentric}. The lineage runs back to Jamaludin et
al.'s longitudinal self-supervision \cite{jamaludin2017selfsupervised}.

\subsection{The Clinical Asymmetry of Error}
\label{sec:lr-error-asymmetry}

The most important point in this section is that not all errors are equal, and
that standard metrics conceal the difference. Grading a severe stenosis as
normal risks a missed diagnosis and delayed decompression; grading a normal disc
as moderate risks an unnecessary review. Overall accuracy, macro-F1 and even
Cohen's $\kappa$ treat these identically.

Three studies take the asymmetry seriously. Acharya and Kansakar report the
severe-to-normal misclassification rate as a headline result alongside balanced
accuracy, arguing explicitly that conventional metrics do not capture the
differential clinical risk \cite{acharya2026disccentric}. Wu et al.\ characterise
SpineNetv2's error profile as specificity-oriented, with false negatives
exceeding false positives, and conclude that the system is suitable as a
confirmatory second reader but that reliance on its negative findings is not
recommended---a conclusion about deployment that follows from the shape of the
errors rather than their number \cite{wu2026external}. The RSNA weighted metric
encodes the same judgement into the benchmark itself.

Ishimoto et al.\ supply the complementary observation: collapsing four grades to
a severe-versus-rest dichotomy raised agreement from 65.7\% to 94.1\% with
$\kappa = 0.75$ \cite{ishimoto2020could}. If the clinically decisive question is
whether a level is severe, then performance on that question is considerably
better than the multi-class figures suggest, and evaluation should report both.
This chapter therefore recommends, and Chapter~3 adopts, reporting per-class
recall for the severe class alongside aggregate measures.

% ===========================================================================
\section{Datasets and Benchmarks}
\label{sec:lr-datasets}

\subsection{Private and Single-Institution Cohorts}
\label{sec:lr-private-data}

Most work before 2024 used private single-institution data, with cohort sizes
spanning three orders of magnitude---from 75 exams \cite{gao2022automatic} and
146 consecutive patients \cite{lehnen2021detection} through 200 studies
\cite{bharadwaj2023deep} and 542 images \cite{won2020spinal} to 4{,}075 patients
\cite{lu2018deepspine} and 15{,}254 axial images \cite{su2022automatic}.
Aggregate performance correlates only loosely with cohort size, because larger
private datasets are frequently drawn from a single scanner and protocol.

Two consequences follow, and both are structural rather than incidental. Results
are not comparable across studies, since each is measured against a different
label distribution and reference standard; and models cannot be independently
re-evaluated, since the data cannot be shared. Both systematic reviews of the
period identify this as the field's central methodological weakness
\cite{compte2023are,mendes2026deep}.

\subsection{Population and Epidemiological Cohorts}
\label{sec:lr-population-data}

A smaller body of work uses population cohorts, which differ from clinical
convenience samples in a way that matters for generalisation: participants are
not selected by referral, so the severity distribution reflects the population
rather than the imaged-because-symptomatic subset. The Northern Finland Birth
Cohort 1966, comprising lumbar MRI from participants imaged at age 46--47, is the
most used \cite{mcsweeney2023external,kowlagi2022a}; the Wakayama Spine Study
\cite{ishimoto2020could} and the Hong Kong Disc Degeneration Population-Based
Cohort \cite{cheung2022learningbased} serve similar roles. Kowlagi et al.\ are
explicit that evaluating on a population cohort rather than a clinical sample is
a deliberate methodological choice, reflecting the transferability of results to
downstream research \cite{kowlagi2022a}.

\subsection{Public Lumbar MRI Datasets}
\label{sec:lr-public-data}

Public data arrived late to this field relative to CT, where large annotated
vertebral collections have existed for years. SPIDER is the principal
segmentation resource: 447 sagittal T1 and T2 series from 218 patients across
four hospitals, with reference segmentations of vertebrae, discs and spinal
canal, per-level radiological gradings, baseline and nnU-Net reference
performance, and a continuous challenge on a sequestered split
\cite{vandergraaf2024lumbar}. Its influence is visible in the number of
subsequent methods benchmarked on it
\cite{ahmed2024pioneering,silveira2025automated,patil2026mribased,ghobrial2025deep}.
Smaller public collections support classification work
\cite{natalia2024lumbar,shahzadi2023nerve,nikpasand2024automated}.

\subsection{The RSNA 2024 Challenge and LumbarDISC}
\label{sec:lr-lumbardisc}

The RSNA Lumbar Degenerative Imaging Spine Classification dataset is the largest
and most diverse expertly annotated lumbar spondylosis collection assembled, and
it is the dataset used in this thesis \cite{richards2026the}. It comprises MRI
studies from 2{,}697 patients and 8{,}593 image series drawn from eight
institutions across six countries and five continents---a geographic spread
chosen deliberately so that models trained on it face genuine acquisition
heterogeneity rather than a single site's conventions.

Inclusion required a sagittal T2-like series (conventional spin echo T2, STIR or
Dixon), a sagittal T1 series and an axial T2 series covering at least three
lumbar disc levels, from outpatients aged 18 or over imaged for degenerative
disease. Studies with lumbosacral hardware, active non-degenerative pathology,
severe scoliosis or substantial artifact were excluded \cite{richards2026the}.
Annotation was performed by expert volunteer neuroradiologists and
musculoskeletal radiologists recruited through the ASNR
\cite{rsna2024rsna}. The associated challenge \cite{rsna2024rsna} made the
dataset the field's common benchmark and shifted attention decisively from
whole-image classification toward anatomically localised, per-level severity
grading.

\subsection{Label Structure, Distribution, and the Weighted Metric}
\label{sec:lr-label-structure}

The label structure is the defining feature of the benchmark and warrants precise
statement. Each study receives twenty-five graded outputs: five conditions
---spinal canal stenosis, left and right neural foraminal narrowing, left and
right subarticular stenosis---at each of five disc levels from L1/L2 to L5/S1.
Each is assigned one of three ordinal classes: Normal/Mild, Moderate or Severe.
Expert coordinate annotations accompany the grades, so anatomical localisation
can be learned in a supervised manner rather than inferred
\cite{rsna2024rsna,richards2026the}.

Three properties follow. The task is \emph{not} generic three-class image
classification, and results from studies that treat it as such are not
comparable. The severity distribution is severely imbalanced and differently so
for each condition (Section~\ref{sec:lr-distributions}). And the challenge metric
is a sample-weighted log loss with weights of 1, 2 and 4 across the three
classes, which encodes the clinical asymmetry of
Section~\ref{sec:lr-error-asymmetry} directly into the evaluation. A model tuned
for unweighted accuracy is therefore optimising a different objective from the
one the benchmark measures.

\subsection{Challenge Solutions and the Practices They Standardised}
\label{sec:lr-challenge-solutions}

The solutions converged on a common template, and that convergence is itself a
finding. Nearly all adopted multi-stage pipelines in which localisation precedes
classification: a detection model---commonly a U-Net, YOLO variant or Faster
R-CNN with a feature pyramid---identifies each disc level on the sagittal
series, and a second-stage classifier grades tightly cropped regions, typically
using 2.5D representations that combine a 2D backbone with a sequence model or
attention over adjacent slices.

M-SCAN exemplifies the approach and reports AUROC 0.971 on 1{,}975 studies,
using a U-Net with ResNet50 encoder to locate stenosis points on sagittal images,
projecting these into 3D via DICOM orientation metadata to select corresponding
axial slices, and applying cross-attention to align sagittal and axial features
\cite{batra2025mscan}. Liu et al.\ pair MobileNetV2 with U-Net for efficiency,
reporting 94.93\% accuracy under five-fold cross-validation
\cite{liu2025automated}; Park et al.\ apply transformer-based multimodal fusion
\cite{park2025multiview}; and further challenge-derived work continues to appear
\cite{patel2025automated,garciadecelis2025deep}.

Two cautions attach to these figures. Most are cross-validation results on the
public training split rather than external validation, so they measure fit to
LumbarDISC rather than transportability beyond it. And the top solutions rely on
large ensembles, which inflate reported accuracy relative to what a single
deployable model achieves---a gap directly relevant to the efficiency question
this thesis examines.

% ===========================================================================
\section{Evaluation, Validation, and Generalization}
\label{sec:lr-validation}

\subsection{Metrics and Their Appropriateness}
\label{sec:lr-metrics}

Reported metrics in this field are heterogeneous and frequently ill-matched to
the task. Overall accuracy remains the most common headline figure despite being
close to meaningless under the severity distributions of
Section~\ref{sec:lr-distributions}. Area under the ROC curve is widely reported
\cite{batra2025mscan,wang2024machine} but is defined for binary problems and must
be extended by one-versus-rest averaging for three-class grading, which conceals
which class is failing.

More appropriate choices appear in the better studies. Cohen's and Gwet's
$\kappa$ measure agreement above chance and are directly comparable with the
human reliability figures of Section~\ref{sec:lr-reliability}, which is the
correct reference point \cite{lim2022improved,su2022automatic,mcsweeney2023external}.
Balanced accuracy corrects for prevalence \cite{mcsweeney2023external,acharya2026disccentric}.
Mean absolute error in grades exploits the ordinal structure and distinguishes a
one-grade slip from a two-grade error \cite{niemeyer2021a,wu2026external}.
Ishimoto et al.\ report both exact-grade agreement and the collapsed
severe-versus-rest figure, making the difference visible rather than choosing
between them \cite{ishimoto2020could}.

The recommendation that follows is to report a small panel rather than a single
number: a prevalence-corrected aggregate, an agreement statistic comparable to
inter-reader values, per-class recall for the severe class, and the distribution
of error magnitudes across grade boundaries.

\subsection{Internal Validation and Optimistic Reporting}
\label{sec:lr-internal-validation}

A recurring pattern is that internally validated performance is high and
externally validated performance is materially lower. Several reported accuracies
in the reviewed literature exceed 95\%
\cite{shahzadi2023nerve,liawrungrueang2023automatic,liu2025automated,tabarestani2026a},
which sits uneasily beside inter-reader agreement of $\kappa = 0.49$ for
subarticular stenosis \cite{lurie2008reliability}. Where a model appears to
exceed the reproducibility of the process that generated its labels, the
explanation usually lies in the evaluation protocol: single-institution data,
cross-validation rather than held-out external testing, labels from a single
reader, or a task collapsed to a binary decision.

The point is not that these studies are wrong on their own terms, but that
their numbers do not transfer. Compte et al.\ state the general finding
directly: algorithms did not perform as well in replication or external
validation as in their development studies \cite{compte2023are}.

\subsection{External Validation and Domain Shift}
\label{sec:lr-generalization}

External validation on geographically and temporally separate cohorts is the
strongest available evidence of clinical viability, and a small number of
studies supply it.

McSweeney et al.\ evaluated SpineNet on 1{,}331 participants from the Northern
Finland Birth Cohort 1966, a dataset separate from its UK training data in both
geography and time. Balanced accuracy was 78\% for disc degeneration and 86\%
for Modic changes, with Pfirrmann reliability against expert radiologists of Lin
concordance 0.86 and Cohen $\kappa = 0.68$, and Modic detection at
$\kappa = 0.74$; both were essentially unchanged in a low-back-pain subset. The
informative detail is the error structure: 20.83\% of discs were graded
differently from human raters, but only 0.85\% differed by more than one grade
\cite{mcsweeney2023external}. Nigru et al.\ evaluated SpineNetV2 across eleven
features on 1{,}747 discs from 353 patients, finding agreement above 0.7 for most
targets but distinctly lower for foraminal stenosis and herniation, attributed
to the limits of sagittal input \cite{nigru2024external}. Wu et al.\ tested
SpineNetv2 on 491 patients and 2{,}455 discs, reporting overall accuracy of
83.5--97.5\% and superiority over a junior orthopaedic surgeon on canal stenosis,
spondylolisthesis and bilateral foraminal stenosis \cite{wu2026external}. Grob et
al.\ contribute a further independent validation on 882 scans
\cite{grob2022external}.

Studies that build external testing into their development are still uncommon
but instructive. Tumko et al.\ evaluated on 150 studies graded by a panel of
seven radiologists with majority voting and external adjudication
\cite{tumko2024a}; Su et al.\ used a 1{,}273-image external set from a different
hospital and observed accuracy falling by roughly 7--10 points relative to
internal testing \cite{su2022automatic}; Zhang et al.\ found classification
accuracy dropping from 87.70\% to 74.23\% externally \cite{zhang2023deep}; Zeng
et al.\ report an unusually small external drop, from 93.7\% to 92.2\%
\cite{zeng2025gpt4lfs}.

\subsection{Detection Versus Fine-Grained Grading}
\label{sec:lr-detection-vs-grading}

A consistent and under-appreciated finding is that these two capabilities
generalise differently. Zhang et al.\ provide the clearest illustration:
moving to an external institution, detection degraded only modestly---mean IoU
0.82 to 0.70, precision 98.4\% to 96.3\%, sensitivity 99.4\% to 97.8\%---while
classification accuracy fell from 87.70\% to 74.23\% \cite{zhang2023deep}. Wu et
al.\ report the same asymmetry in different terms: binary detection across five
pathologies held up well while ordinal Pfirrmann grading remained the limiting
factor, degrading further in older patients and upper lumbar discs
\cite{wu2026external}. Yilihamu et al.\ show the effect within a single study as
task granularity increases: segmentation IoU remained above 97\% externally,
while classification precision fell from 81.21\% to 74.50\% for an 18-category
problem yet held at 92.51\% to 90.07\% for four-category severity
\cite{yilihamu2025quantification}.

The implication is that where a system fails is predictable. Finding the anatomy
is robust; judging how bad it looks is not. Reporting a single aggregate figure
for a pipeline that performs both obscures this, and evaluation should separate
the two stages.

\subsection{Benchmarking Against Clinicians}
\label{sec:lr-clinician-benchmark}

Comparison against human readers is the most clinically meaningful benchmark, and
the better studies construct it carefully. Won et al.\ provide the cleanest
design: two experts each labelled the data, a model was trained on each expert's
labels, and the key result was that the difference between the two experts was
statistically significant while the difference between each expert and their
corresponding model was not \cite{won2020spinal}. Tumko et al.\ compared against
a seven-radiologist panel average and found the model matched or exceeded it on
all three stenosis types \cite{tumko2024a}. Lee et al.\ benchmarked against eight
participants spanning four experience levels \cite{lee2026deep}, and Wu et al.\
against orthopaedic surgeons of differing seniority \cite{wu2026external}.

Two cautions apply. Comparisons against junior readers or against surgeons rather
than radiologists set a lower bar than subspecialty radiological practice
\cite{ke2025integrating,wu2026external}. And readers in these studies typically
work from cropped images without clinical history, which is not how they read in
practice---so such comparisons measure a narrower task than clinical
interpretation.

\subsection{Longitudinal Stability and Reproducibility}
\label{sec:lr-longitudinal}

One-shot accuracy does not establish that a model is stable over repeat imaging,
which matters for monitoring progression. Murto et al.\ supply the only
longitudinal evidence in this collection, comparing SpineNet's Pfirrmann grading
against radiologists in 19 volunteers imaged at age 37 and again at 51.
Radiologist-to-SpineNet concordance was 0.54 at both timepoints, clearly below
the 0.83 and 0.78 observed between radiologists, and SpineNet systematically
assigned grade~1 to several discs and grade~5 to more discs than the readers did
\cite{murto2025comparison}. The cohort is small, but a systematic offset that is
stable across time is exactly the failure mode that one-shot evaluation cannot
detect. Cheung et al.\ approach the temporal dimension differently, predicting
five-year progression from baseline imaging across 1{,}343 participants
\cite{cheung2022learningbased}.

\subsection{Reporting Quality and Recurring Weaknesses}
\label{sec:lr-reporting}

The systematic reviews converge on a consistent diagnosis. Wang et al.\ pooled
twelve studies covering 15{,}044 patients, obtaining pooled sensitivity 0.84 and
specificity 0.87 with SROC AUC 0.92, but with heterogeneity of $I^2$ near 99\%
and four of twelve studies at high risk of bias; they concluded that no current
model is reliable enough for clinical practice, citing the need for accepted
reference standards, larger datasets, external validation and fuller reporting
\cite{wang2024machine}. Compte et al.\ found no significant difference between
algorithm families and identified the development-to-validation performance drop
as the field's central problem \cite{compte2023are}. Mendes et al.\ found that
most of 56 segmentation studies relied on private datasets and lacked external
validation \cite{mendes2026deep}. Liawrungrueang et al.\ report metrics spanning
71.5\% to 99\% across twenty studies---a range so wide as to indicate that the
underlying tasks are not comparable \cite{liawrungrueang2025artificial}.

Four weaknesses recur: single-institution data without external testing; ground
truth from a single reader with no adjudication; task definitions that vary
between studies while sharing a label vocabulary; and metric selection that
flatters imbalanced problems. The present work adopts the public LumbarDISC
benchmark, its panel-derived labels and its weighted metric specifically to
avoid the first, second and fourth of these.

% ===========================================================================
\section{Interpretability and Clinical Integration}
\label{sec:lr-interpretability}

\subsection{Saliency, Evidence Hotspots, and Grad-CAM}
\label{sec:lr-saliency}

The earliest interpretability contribution in this domain is SpineNet's evidence
hotspots, which map output-layer gradients back to the input to produce a heatmap
per radiological score, obtained without any localisation supervision
\cite{jamaludin2017spinenet}. Grad-CAM \cite{selvaraju2017gradcam} has since
become the default. Tabarestani et al.\ use it to confirm that their cascade
attends to clinically relevant regions \cite{tabarestani2026a}, and Silveira et
al.\ report that attention concentrated on vertebral bodies and nucleus pulposus
in spatial correspondence with manual segmentations \cite{silveira2025automated}.

The value of these methods should not be overstated. A saliency map shows where
a network responded, not why, and confirming that attention falls on plausible
anatomy is a weak check---it excludes gross reliance on background artefact but
does not establish that the model has learned the morphological criteria the
grading system specifies.

\subsection{Inherently Interpretable Models}
\label{sec:lr-interpretable}

A stronger position is to build models whose reasoning is legible by
construction. Bharadwaj et al.\ provide the most striking result in this
literature: alongside a Big Transfer classifier they trained a decision-tree
classifier operating on segmentation-derived features for central canal stenosis,
and the interpretable model \emph{outperformed} the black box, reaching
$\kappa = 0.80$ against 0.54 and approaching inter-reader agreement of 0.74--0.86
\cite{bharadwaj2023deep}. Quantitative morphometry belongs to the same family:
a measured dural sac area \cite{ghobrial2025deep} or a quantitative degeneration
criterion \cite{zheng2022deep} is inspectable in a way a logit is not. Patil et
al.\ make the related point that radiomic features carried the entire
discriminative signal for Pfirrmann grading while deep features added nothing
measurable, and that radiomics additionally yields interpretable shape and
intensity descriptors usable as imaging biomarkers \cite{patil2026mribased}.

Together these findings undercut the assumption that interpretability must be
purchased at the cost of accuracy---at least for targets whose clinical
definition is already explicitly morphological.

\subsection{Vision--Language Models and Report Generation}
\label{sec:lr-vlm}

The most recent direction outputs prose rather than labels. Islam~Sk et al.\
present an explainable vision--language framework that classifies severity,
segments the anatomy and generates radiologist-style reports, built on
biomedical foundation models, with a spatial patch cross-attention module that
preserves the anatomical hierarchies global pooling discards and an adaptive
PID-Tversky loss for class imbalance; they report 90.69\% classification
accuracy, macro-averaged Dice 0.9512 and CIDEr 92.80\%
\cite{islamsk2026an}. Zeng et al.\ take a narrower approach, fusing ConvNeXt
image features with GPT-4o-generated descriptions for foraminal stenosis
classification across three centres and seven scanners \cite{zeng2025gpt4lfs}.
These systems are early and their language metrics measure fluency rather than
clinical correctness, but they indicate where the field is heading.

\subsection{Reader-Assistance Studies and Workflow Impact}
\label{sec:lr-reader-assistance}

The most consequential evidence for clinical value comes not from model metrics
but from measuring what changes when readers use the model. Lim et al.\ conducted
the definitive study of this kind: eight radiologists with 2--13 years of
experience read the same studies with and without deep-learning assistance,
separated by a one-month washout, against a 32-year-veteran musculoskeletal
radiologist as reference. Interpretation time per study fell from 124--274
seconds to 47--71 seconds, and interobserver agreement was superior or equivalent
for every stenosis grading. The largest gain fell to general and in-training
radiologists on four-class foraminal stenosis, where $\kappa$ rose from 0.39 to
0.71 and 0.70 \cite{lim2022improved}. Gao et al.\ observed a comparable effect
for Modic assessment, where AI assistance raised agreement between a junior
radiologist and a senior neuroradiologist from $\kappa = 0.52$ to 0.58
\cite{gao2022automatic}.

The pattern in both is that assistance helps least-experienced readers most and
standardises rather than replaces judgement. This reframes the objective: the
target is not a model that outperforms the best radiologist, but one that raises
the floor.

\subsection{Uncertainty and Calibration}
\label{sec:lr-uncertainty}

Uncertainty estimation is conspicuously underdeveloped in this literature. Where
errors concentrate at grade boundaries and where the reference standard is itself
unreliable, a calibrated confidence would let borderline cases be referred rather
than silently graded---a natural fit for triage. Almost none of the reviewed work
reports calibration. Patil et al.\ are the exception, applying isotonic
calibration on a validation split and reporting Brier scores alongside
discrimination metrics \cite{patil2026mribased}. Since modern networks are
systematically overconfident \cite{guo2017calibration}, an uncalibrated severity
probability should not be interpreted as a clinical confidence, and this remains
an open opportunity rather than a solved problem.

% ===========================================================================
\section{Synthesis and Research Gap}
\label{sec:lr-gap}

\subsection{What the Literature Has Established}
\label{sec:lr-established}

Five findings are supported well enough to be treated as settled.

\textbf{Anatomical localisation precedes classification.} From DeepSPINE
\cite{lu2018deepspine} through the interpretable pipelines of Bharadwaj et al.\
\cite{bharadwaj2023deep} to essentially every leading RSNA challenge solution
\cite{batra2025mscan}, systems that localise the disc level before grading it
outperform those that classify whole images. Detection is close to solved, with
several studies reporting near-perfect level assignment
\cite{pan2021automatically,ghosh2012a}.

\textbf{Automated grading is competitive with human reliability, because human
reliability is low.} Inter-reader agreement is substantial for central canal
stenosis but only moderate for foraminal and poor for subarticular grading
\cite{lurie2008reliability}. Against that benchmark, models sit inside the band
of human disagreement \cite{won2020spinal,mcsweeney2023external,tumko2024a}.
This is a real result, but it is a statement about the modesty of the reference
standard as much as the capability of the models.

\textbf{Errors concentrate at adjacent grade boundaries.} Disagreement above one
grade is rare---0.85\% of discs in the largest external validation
\cite{mcsweeney2023external}---and collapsing to a severe-versus-rest decision
raises agreement substantially \cite{ishimoto2020could}. The difficulty lies in
boundary placement, not gross misclassification.

\textbf{Detection generalises across institutions; fine-grained grading does
not.} The asymmetry is consistent and large
\cite{zhang2023deep,wu2026external,yilihamu2025quantification}.

\textbf{Class imbalance must be addressed explicitly.} With 85.4\% of spinal
canal grades normal or mild \cite{richards2026the}, unweighted objectives yield
degenerate solutions, and weighted sampling, weighted or focal losses and
contrastive pretraining are all established mitigations
\cite{lin2017focal,acharya2026disccentric}.

\subsection{What Remains Contested}
\label{sec:lr-contested}

Three questions remain genuinely open, and the disagreement is substantive
rather than terminological.

\emph{Does architectural sophistication help?} The general computer vision
literature would predict transformers to outperform convolutional networks given
sufficient data, but the lumbar evidence does not bear this out. Kowlagi et al.\
report a tuned CNN pipeline beating transformer approaches \cite{kowlagi2022a};
Lee et al.\ find a CNN superior and a transformer merely comparable
\cite{lee2026deep}; Silveira et al.\ beat TransUNet with an enhanced U-Net
\cite{silveira2025automated}; Udomluck et al.\ find conventional VGG19 features
generalising better than DINOv2 \cite{udomluck2026multitask}; and Patil et al.\
find deep features adding nothing over radiomics \cite{patil2026mribased}. Set
against this, hybrid architectures report gains
\cite{chen2024symtc,zheng2022deep}, and cross-sequence attention reports large
ones \cite{nguyen2026crossspine}.

\emph{Does the ordinal structure of the labels admit exploitation?} Errors
demonstrably cluster at adjacent grades, yet the most systematic test of
ordinal-aware objectives found none improved on plain cross-entropy
\cite{niemeyer2021a}.

\emph{How much imaging is actually required?} Clinical practice and the strongest
benchmark results use all three sequences
\cite{richards2026the,batra2025mscan}, while sagittal-only systems report
competitive performance \cite{acharya2026disccentric,kowlagi2022a} and
sagittal-only systems also fail in a specific, predicted way on the lateral
compartments \cite{nigru2024external}.

\subsection{Gap 1: Architecture Comparison Under Controlled Conditions}
\label{sec:lr-gap-architecture}

The comparative studies above are informative but none isolates the variable of
interest. Three confounds recur.

First, \emph{unequal pipelines}. Compared architectures rarely receive identical
preprocessing, localisation and training budgets, so a reported difference may
reflect engineering effort rather than representational capacity---precisely the
argument Kowlagi et al.\ advance when attributing their CNN's advantage to
pipeline design rather than to convolution as such \cite{kowlagi2022a}. Given the
evidence of Section~\ref{sec:lr-localize-first} that localisation dominates
performance, a comparison in which one arm is better localised than the other
measures preprocessing.

Second, \emph{frozen features}. Several comparisons extract features from frozen
pretrained backbones rather than fine-tuning end to end
\cite{udomluck2026multitask,patil2026mribased}. This systematically
disadvantages transformers, whose weaker inductive biases make them more
dependent on in-domain adaptation.

Third, \emph{fixed input configuration}. Every comparison holds the sequence
configuration constant, so the interaction between architecture and available
context is untested. This matters because the theoretical case for global
attention is strongest exactly where context must be inferred rather than
observed---that is, where a view is missing.

A controlled comparison therefore requires a shared localisation stage, a shared
training and augmentation protocol, comparable capacity and tuning budgets, and
end-to-end fine-tuning of both backbones, evaluated on the same benchmark with
the same metric. No study in this review satisfies all four conditions.

\subsection{Gap 2: Sequence Configuration and the Accuracy--Efficiency Trade-off}
\label{sec:lr-gap-sequence}

The second gap concerns input rather than model. Multi-sequence pipelines carry
costs that are rarely quantified alongside their accuracy: greater memory and
compute, cross-plane geometric registration dependent on DICOM metadata
\cite{batra2025mscan,pan2021automatically}, longer inference, and brittleness
when a sequence is missing or degraded (Section~\ref{sec:lr-missing}). Against
these, the incremental diagnostic yield of each additional sequence has not been
measured under controlled conditions.

The literature supplies both sides of the argument but no direct test. Evidence
that the third sequence matters is indirect, inferred from the failure of
sagittal-only systems on foraminal and subarticular targets
\cite{nigru2024external,wu2026external}; evidence that it may not is likewise
indirect, resting on sagittal-only systems that perform respectably against
different baselines on different data
\cite{acharya2026disccentric,kowlagi2022a}. Because these come from different
architectures, cohorts, localisation strategies and metrics, the contribution of
sequence configuration cannot be separated from everything else that differs.

Two design requirements follow. The comparison must hold architecture and
pipeline constant while varying only the input configuration; and it must report
results per condition rather than in aggregate, since the expected penalty is
concentrated in the foraminal and subarticular targets and an aggregate figure
would average away the structure of interest.

\subsection{Research Questions and Objectives}
\label{sec:lr-questions}

The two gaps combine into a single factorial investigation. Crossing
architecture with input configuration answers both questions and, additionally,
tests whether they interact---whether a transformer's global attention can
partially compensate for absent multiplanar data, which neither gap addresses
alone. This yields the following research questions:

\begin{description}
  \item[RQ1] Under matched localisation, training and evaluation protocols, how
  do convolutional backbones (ResNet50, EfficientNet-B4) compare with a
  hierarchical vision transformer (Swin) for per-level, per-condition severity
  grading of lumbar degenerative disease?

  \item[RQ2] What is the cost in diagnostic performance of reducing the input
  from the full multi-sequence configuration (Sagittal T1, Sagittal T2/STIR,
  Axial T2) to a single-sequence configuration (Sagittal T2), and how is that
  cost distributed across the five graded conditions?

  \item[RQ3] Do architecture and sequence configuration interact---specifically,
  is the penalty for single-sequence input smaller for the transformer than for
  the convolutional backbones?

  \item[RQ4] Which class-imbalance strategies most effectively reduce
  clinically consequential errors, in particular severe cases graded as
  normal or mild?
\end{description}

The corresponding objectives are to implement a shared localisation and
preprocessing stage so that all arms receive identical input regions; to train
each architecture under both configurations with matched budgets; to evaluate
using the benchmark's weighted metric together with per-condition and per-class
reporting; and to quantify computational cost---parameters, memory and inference
time---so that any accuracy difference can be weighed against it.

\subsection{Scope and Delimitations}
\label{sec:lr-scope-delimit}

Four boundaries are stated explicitly. The work addresses \emph{severity
grading}, taking localisation as a shared and largely solved precondition
(Section~\ref{sec:lr-detection}) rather than a contribution. It evaluates on
\emph{LumbarDISC} \cite{richards2026the}; because that collection excludes
studies with hardware, non-degenerative pathology, severe scoliosis or
substantial artifact, conclusions do not extend to those populations, and no
independent external cohort is available to test transportability beyond it. It
compares a \emph{bounded set of architectures}---two convolutional and one
transformer---chosen as representative rather than exhaustive, and does not
evaluate hybrid or vision--language designs
\cite{chen2024symtc,islamsk2026an}. Finally, it reports model performance and
computational cost, not clinical impact: the reader-assistance evidence of
Section~\ref{sec:lr-reader-assistance} \cite{lim2022improved} shows that
workflow benefit requires a prospective reader study, which lies outside the
present scope.

Two constraints inherited from the literature bound the interpretation of any
result. Performance is capped by the reliability of the reference standard
(Section~\ref{sec:lr-label-ceiling}), so figures approaching perfect agreement
should be read as evidence of a protocol artefact rather than of clinical
readiness. And single-benchmark evaluation cannot establish generalisation; the
consistent development-to-validation drop documented in
Section~\ref{sec:lr-generalization} applies to this work as much as to any other.

% ===========================================================================
\section{Chapter Summary}
\label{sec:lr-summary}

This chapter reviewed the literature on automated assessment of lumbar
degenerative disease from MRI, from the clinical grading systems that define the
task through to the architectures that currently attempt it.

The clinical sections established that severity labels originate in
morphological grading systems---Pfirrmann for disc degeneration
\cite{pfirrmann2001magnetic}, Schizas and Lee for central canal
\cite{schizas2010qualitative,lee2011a}, Lee for foraminal
\cite{lee2010a}---whose inter-reader reliability ranges from substantial to poor
depending on the compartment \cite{lurie2008reliability}. This imposes a ceiling
on achievable performance and makes agreement statistics, rather than accuracy,
the appropriate measure. The labels are ordinal, errors concentrate at adjacent
boundaries, and the clinical cost of an error depends on its direction.

The methodological sections traced the field from handcrafted features through
SpineNet's multi-task convolutional grading \cite{jamaludin2017spinenet} to
contemporary multi-stage, multi-view pipelines \cite{batra2025mscan}. The most
robust design principle to emerge is that anatomical localisation should precede
classification \cite{lu2018deepspine,kowlagi2022a}. The most robust empirical
regularity is that detection transfers across institutions while fine-grained
grading does not \cite{zhang2023deep,wu2026external}.

Two questions remain unresolved. Whether attention-based backbones offer any
advantage over well-tuned convolutional ones is contested, with the balance of
lumbar-specific evidence currently favouring CNNs
\cite{kowlagi2022a,lee2026deep,silveira2025automated}, yet no study has compared
them under matched localisation, training and evaluation conditions. And whether
the full three-sequence protocol is necessary, or whether a single ubiquitous
sequence suffices, has been argued from both sides without a controlled test.
These two questions, and their interaction, constitute the gap this thesis
addresses. Chapter~3 sets out the experimental design through which they are
answered.

% ===========================================================================
\printbibliography[heading=bibintoc,title={References}]

\end{document}
 from your main file.
% ===========================================================================

\documentclass[12pt,a4paper]{report}

\usepackage{lmodern}          % scalable fonts; required for microtype expansion
\usepackage[T1]{fontenc}
\usepackage[utf8]{inputenc}
\usepackage[english]{babel}
\usepackage{microtype}
\usepackage[margin=2.5cm]{geometry}
\usepackage{setspace}
\onehalfspacing

\usepackage{amsmath,amssymb}
\usepackage{booktabs}
\usepackage{longtable}
\usepackage{array}
\usepackage{graphicx}
\usepackage[hidelinks]{hyperref}

\usepackage[
  backend=biber,
  style=numeric-comp,
  sorting=none,
  maxbibnames=6,
  minbibnames=3,
  giveninits=true,
  doi=true,
  url=false,
  isbn=false
]{biblatex}

\addbibresource{../lumbar_spine_mri_ai_literature_inventory.bib}
\addbibresource{references-extra.bib}

% The inventory .bib carries internal bookkeeping fields (provenance note, the
% per-paper summary, the local PDF path). These are working metadata, not
% bibliographic data, so they are suppressed from the printed reference list.
% Leaving them in also breaks compilation, since they contain raw underscores.
\AtEveryBibitem{%
  \clearfield{note}%
  \clearfield{comment}%
  \clearfield{annotation}%
  \clearfield{abstract}%
  \clearfield{file}%
}

\begin{document}

\chapter{Literature Review}
\label{ch:litreview}

% ===========================================================================
\section{Introduction}
\label{sec:lr-intro}

\subsection{Scope and Objectives of the Review}
\label{sec:lr-scope}

Lumbar degenerative disease is among the most common indications for magnetic
resonance imaging (MRI) worldwide, and the interpretation of those images is a
repetitive, subjective and time-consuming task that has attracted sustained
attention from the medical image analysis community \cite{lee2024artificial,zhang2024a}.
This chapter reviews that body of work. Its purpose is not to catalogue every
published model, but to establish what is genuinely settled, what remains
contested, and where a controlled experiment can still add knowledge.

The review is organised around a single guiding question: \emph{what determines
whether an automated system can grade lumbar degenerative severity reliably
enough to be clinically useful?} Four candidate answers recur in the
literature---the quality and definition of the reference labels, the degree to
which the anatomy is localised before it is classified, the choice of
representational backbone, and the configuration of imaging sequences supplied
to the model. Each is examined in turn, and the chapter closes by identifying
the two of these that remain inadequately resolved and that this thesis
therefore addresses.

Three specific objectives follow from this. First, to characterise the clinical
grading systems that supply the ground truth for every supervised model in this
field, and to quantify the reliability ceiling those systems impose. Second, to
trace the methodological evolution from handcrafted features to contemporary
convolutional and attention-based architectures, with particular attention to
comparative claims between the two. Third, to assess the evidence on how much
diagnostic information is carried by each MRI sequence and plane, and thus
whether the computational expense of multi-sequence processing is warranted.

\subsection{Search Strategy, Sources, and Inclusion Criteria}
\label{sec:lr-search}

The literature underpinning this chapter was assembled into a structured
inventory of 92 records spanning 2001 to 2026. Sources were identified through
PubMed, Scopus, IEEE~Xplore and arXiv, using combinations of the terms
\emph{lumbar spine}, \emph{MRI}, \emph{spinal stenosis}, \emph{disc
degeneration}, \emph{deep learning}, \emph{convolutional neural network},
\emph{transformer}, \emph{segmentation} and \emph{grading}. Reference lists of
the retrieved systematic reviews \cite{compte2023are,wang2024machine,liawrungrueang2025artificial,mendes2026deep}
were searched by hand for records the database queries had missed, and the
citation graphs of the most influential frameworks were traced forward.

Records were retained if they satisfied at least one of four criteria: they
define or validate a clinical grading standard for lumbar degenerative disease;
they present an automated method for detecting, segmenting or grading lumbar
pathology on MRI; they externally validate such a method; or they describe a
dataset or benchmark on which such methods are trained. Studies confined to
computed tomography or plain radiography were excluded unless they establish a
methodological principle transferable to MRI, as were studies of the cervical
or thoracic spine except where explicitly used for contrast.

Two caveats on the evidence base should be recorded at the outset. First, the
field's publication rate accelerated sharply after 2024, and a substantial
fraction of the most directly relevant work
\cite{batra2025mscan,acharya2026disccentric,nguyen2026crossspine,islamsk2026an}
exists as preprints that have not completed peer review; these are cited as
indicative of methodological direction rather than as settled results. Second,
performance figures reported across studies are frequently not comparable,
because they derive from different datasets, different label definitions and
different evaluation protocols. Wherever numerical results are quoted in this
chapter, the cohort and validation regime that produced them are stated
alongside, and cross-study comparisons are avoided unless the underlying
protocols genuinely align.

\subsection{Structure of the Chapter}
\label{sec:lr-structure}

Sections~\ref{sec:lr-clinical} to~\ref{sec:lr-grading} establish the clinical
substrate: the disease, the imaging protocol used to observe it, and the
grading systems that convert observation into ordinal labels. This material is
not background decoration. The severity classes a model is trained to predict
are defined entirely by these systems, and their documented unreliability
propagates directly into every reported accuracy figure in the chapters that
follow.

Section~\ref{sec:lr-evolution} traces the methodological history of the field,
and Section~\ref{sec:lr-localization} isolates anatomical localisation as the
design principle that most consistently separates effective systems from
ineffective ones. Sections~\ref{sec:lr-backbones} and~\ref{sec:lr-context} then
address the two axes on which this thesis experiments: the representational
backbone, and the spatial and sequential context supplied to it.
Section~\ref{sec:lr-imbalance} treats class imbalance and the ordinal structure
of the label space, and Section~\ref{sec:lr-datasets} the datasets and
benchmarks---principally the RSNA LumbarDISC collection---on which the field now
converges. Sections~\ref{sec:lr-validation} and~\ref{sec:lr-interpretability}
consider validation practice and interpretability respectively.
Section~\ref{sec:lr-gap} synthesises the whole and states the research gap.

% ===========================================================================
\section{Clinical Background: Lumbar Degenerative Disease}
\label{sec:lr-clinical}

\subsection{Epidemiology and Socioeconomic Burden}
\label{sec:lr-epidemiology}

Low back pain is the single largest contributor to years lived with disability
worldwide, and lumbar degenerative change is among its most frequently
implicated structural correlates \cite{compte2023are,zheng2022deep}. Lifetime
prevalence is reported as high as 84\%, and in the United States alone the
combined direct medical and indirect productivity costs are estimated between
\$50 and \$100 billion annually \cite{richards2026the}. Lumbar spinal stenosis
in particular is the most common indication for spinal surgery in patients over
65 \cite{schizas2010qualitative,bharadwaj2023deep}, and demographic ageing is
expected to increase both its incidence and the imaging workload it generates
\cite{lee2024artificial}.

The operational consequence for this thesis is one of volume. MRI referrals for
low back pain have grown faster than the radiology workforce able to report
them \cite{ghosh2014supervised,compte2023are}, and each lumbar study requires
the reader to assign severity at multiple anatomical levels across several
sequences---an average of 14 to 19 minutes of expert time per study
\cite{compte2023are}. It is this combination of high volume, repetitive
structure and documented inter-reader variability that makes the task a
plausible target for automation, rather than any claim that machines perceive
the underlying pathology more acutely than radiologists do.

\subsection{The Degenerative Cascade}
\label{sec:lr-cascade}

Degeneration of the lumbar spine is a progressive, multifactorial process
driven by ageing, mechanical loading, genetic predisposition and lifestyle
factors \cite{esposito2025mribased,zhao2025advances}. It begins in the
intervertebral disc, whose nucleus pulposus loses proteoglycan content and
therefore water from approximately the third decade of life onward. The
resulting dehydration reduces the disc's height, its capacity to distribute
axial load, and its signal intensity on T2-weighted imaging---the last of these
being precisely what the Pfirrmann scale measures \cite{pfirrmann2001magnetic}.

Disc height loss initiates a compensatory cascade through the rest of the motion
segment. The annulus fibrosus weakens and may permit bulging or frank herniation
of nuclear material; the ligamentum flavum buckles and hypertrophies; the
posterior facet joints, now bearing altered load, develop osteoarthritic change
and osteophytosis \cite{bharadwaj2023deep,nikpasand2024automated}. Changes in
the adjacent vertebral endplates and marrow---the Modic
phenotypes~\cite{modic1988degenerative}---frequently accompany this process and
have themselves been the target of automated detection
\cite{gao2022automatic,athertya2021classification}. The clinical significance of
the cascade is that these individually modest changes act cumulatively to
encroach on the spaces occupied by neural tissue.

\subsection{The Three Anatomical Forms of Stenosis}
\label{sec:lr-stenosis-types}

Encroachment is conventionally partitioned by anatomical compartment, and this
partition is directly mirrored in the label structure of the dataset used in
this thesis \cite{richards2026the}:

\begin{enumerate}
  \item \textbf{Spinal canal stenosis} is narrowing of the central canal, which
  contains the thecal sac and the cauda equina. It is best appreciated on axial
  T2-weighted images, where the relationship between cerebrospinal fluid and
  nerve rootlets is visible \cite{schizas2010qualitative,lee2011a}. Severe
  central stenosis characteristically produces neurogenic claudication.

  \item \textbf{Neural foraminal narrowing} is constriction of the lateral
  foramina through which the exiting nerve roots pass. It is assessed on
  parasagittal images, where the epidural fat surrounding the root provides the
  contrast that makes compression visible---T1 weighting being particularly
  suited to this \cite{lee2010a}. It produces dermatomal radicular symptoms.

  \item \textbf{Subarticular stenosis}, also termed lateral recess stenosis,
  affects the zone between the medial edge of the facet and the foramen, where
  it compresses the traversing root before it exits at the level below
  \cite{lurie2008reliability,andreisek2014consensus}.
\end{enumerate}

The distinction matters computationally as well as clinically, because the three
compartments are not equally visible in any single plane. Foraminal narrowing is
poorly assessed from axial images and subarticular stenosis poorly from
sagittal ones, a constraint that Section~\ref{sec:lr-context} shows to have
concrete consequences for model design. It is also the compartment that
correlates with which sequence a model must be given, and thus sits at the
centre of the multi-sequence question this thesis examines.

\subsection{Clinical Presentation and the Treatment Decision}
\label{sec:lr-treatment}

Severity grading is not an academic exercise; it partitions patients into
distinct management pathways. Conservative management---physiotherapy, analgesia,
epidural steroid injection---is appropriate for mild and most moderate disease,
whereas severe compression with corresponding symptoms may warrant decompressive
surgery or fusion \cite{schizas2010qualitative,mysliwiec2010msu}. The evidence
that morphological grade carries genuine prognostic weight is strongest for the
Schizas classification, where grade~C and~D patients were substantially more
likely to fail conservative treatment \cite{schizas2010qualitative}, and for the
MSU herniation grades, where size~2 and~3 lesions identify microdiscectomy
candidates \cite{mysliwiec2010msu}.

Two qualifications temper this. First, the correlation between imaging findings
and symptoms is imperfect: anatomical stenosis is demonstrable in a substantial
proportion of asymptomatic individuals \cite{lurie2008reliability,esposito2025mribased},
so imaging severity constrains rather than determines management. Second, and
more consequentially for automated systems, the clinically decisive boundary is
not evenly distributed across the ordinal scale. The distinction between
\emph{normal} and \emph{mild} rarely changes management, whereas the distinction
between \emph{moderate} and \emph{severe} frequently does. A model that
distributes its errors uniformly across grade boundaries is therefore not
clinically equivalent to one of the same overall accuracy that concentrates its
errors at the harmless end---a point developed in
Section~\ref{sec:lr-error-asymmetry}.

% ===========================================================================
\section{MRI Acquisition: Sequences, Planes, and Diagnostic Yield}
\label{sec:lr-mri}

\subsection{Why MRI is the Reference Modality}
\label{sec:lr-why-mri}

MRI is the reference standard for non-invasive assessment of the lumbar spine.
Computed tomography and plain radiography depict osseous anatomy well but
resolve soft tissue poorly, and both expose the patient to ionising radiation.
MRI resolves the intervertebral disc, thecal sac, individual nerve roots and
ligamentous structures with a contrast unavailable to the alternatives
\cite{zhao2025advances,liawrungrueang2025artificial}, and it is the disposition
of exactly these soft-tissue structures that defines every grading system
reviewed in Section~\ref{sec:lr-grading}. MRI has correspondingly displaced CT
myelography for stenosis assessment \cite{tumko2024a}.

\subsection{Complementary Information Across Sequences and Planes}
\label{sec:lr-sequences}

Lumbar protocols are inherently multiplanar and multisequence, because the spine
is a three-dimensional structure whose pathologies are not all visible in one
projection. Three acquisitions dominate clinical practice and constitute the
input available in the dataset used here \cite{richards2026the}.

\emph{Sagittal T2-weighted} imaging renders fluid hyperintense, so cerebrospinal
fluid in the thecal sac appears bright against darker ligament and bone. It is
the sequence on which disc hydration---and therefore Pfirrmann
grade---is judged \cite{pfirrmann2001magnetic}, and on which the
craniocaudal extent of canal compromise is surveyed. Short tau inversion
recovery (STIR) suppresses the signal from fat and thereby exposes marrow
oedema and inflammatory change, making it the sequence of choice for active
Modic type~1 change \cite{gao2022automatic}.

\emph{Sagittal T1-weighted} imaging renders fat hyperintense. Because the exiting
nerve root within the neural foramen is normally surrounded by epidural fat, T1
provides the contrast against which obliteration of that fat---the basis of the
Lee foraminal grading system---becomes visible \cite{lee2010a}.

\emph{Axial T2-weighted} imaging provides the transverse cross-section through
the disc level. It is the only view in which the morphology of the dural sac,
the arrangement of the cauda equina rootlets and the lateral recesses can be
directly assessed, and it therefore underpins both the Schizas and Lee central
canal systems \cite{schizas2010qualitative,lee2011a} and the assessment of
subarticular stenosis \cite{andreisek2014consensus}.

\subsection{Matching Pathology to Sequence}
\label{sec:lr-sequence-matching}

The correspondence between pathology and sequence is not a matter of
convention but of physics, and the empirical literature bears this out with
unusual clarity. The most direct evidence comes from external validation of
sagittal-only systems. Nigru et al.\ evaluated SpineNetV2 across eleven
radiological features on 1{,}747 lumbosacral discs and found agreement above
$\kappa = 0.7$ for most targets but distinctly lower agreement for foraminal
stenosis and disc herniation, attributing the shortfall specifically to the
absence of transverse anatomical context in sagittal input
\cite{nigru2024external}. Hallinan et al.\ approached the same constraint from
the opposite direction, deliberately assigning axial T2 to central canal and
lateral recess grading while reserving sagittal T1 for the neural foramina
\cite{hallinan2021deep}.

This has a direct bearing on the experiments in this thesis. If a
single-sequence model is to be evaluated, the choice of which sequence to
retain is not arbitrary: sagittal T2 is the most defensible candidate because it
is acquired in essentially every back-pain protocol and because it carries the
signal for canal stenosis and disc degeneration, but the literature predicts
that it should be weakest precisely on the foraminal and subarticular targets.
Whether a sufficiently expressive architecture can compensate for that
predicted weakness is an open question, and one this thesis is designed to
answer.

\subsection{Sources of Acquisition Variability}
\label{sec:lr-variability}

Models trained on one institution's images frequently degrade on another's, and
the proximate causes lie in acquisition. Scanner manufacturer and field strength
(typically 1.5 versus 3.0~T), slice thickness and spacing, in-plane resolution,
echo and repetition times, the presence or absence of fat suppression, and
patient positioning all vary between sites \cite{senzgamboa2023automatic,mendes2026deep}.
Axial acquisitions vary further in whether slices form a continuous stack or are
angled to lie parallel to each disc level---a difference the LumbarDISC
inclusion criteria explicitly accommodate \cite{richards2026the}, and one that
any model consuming axial data must tolerate.

S\'aenz~Gamboa et al.\ treat this variability as the research problem in its own
right rather than as a nuisance, evaluating segmentation across multiple centres
and acquisition parameters and finding that ensembles of U-Net variants were
required to maintain performance \cite{senzgamboa2023automatic}. Batra et al.\
make the complementary design argument: rather than assuming uniform histograms
and a fixed slice count per study, their pipeline is built to accept a variable
number of axial and sagittal slices, on the grounds that this reflects the
conditions under which clinical data actually arrives \cite{batra2025mscan}.
Variability is thus not merely a threat to validity but a specification
constraint, and Section~\ref{sec:lr-generalization} returns to how effectively
the field has met it.

% ===========================================================================
\section{Clinical Grading Systems as the Origin of Ground Truth}
\label{sec:lr-grading}

Every supervised model reviewed in this chapter learns to reproduce labels
generated by a human applying one of the systems described below. The
properties of those systems---their definitions, their granularity and above all
their reliability---therefore bound what any model trained on them can achieve.
This section treats them as the measurement instruments they are.

\subsection{Disc Degeneration: Pfirrmann and its Modified Extension}
\label{sec:lr-pfirrmann}

The Pfirrmann classification, introduced in 2001, remains the most widely
adopted scheme for intervertebral disc degeneration and is applied to sagittal
T2-weighted images \cite{pfirrmann2001magnetic}. It assigns one of five ordinal
grades on the basis of four morphological criteria: homogeneity of the nucleus
pulposus, distinction between nucleus and annulus, T2 signal intensity as a
proxy for hydration, and disc height. Grade~I denotes a homogeneous, bright,
normally proportioned disc; grade~V a collapsed disc of uniformly low signal.

Its virtues are simplicity and universality; its limitations are well
documented. It is a subjective visual judgement, it is insensitive to change
over short intervals, and it captures only intrinsic disc state, saying nothing
about Modic change, high-intensity zones or the downstream compression of neural
structures \cite{esposito2025mribased,zhao2025advances}. It also suffers a
ceiling effect in elderly cohorts, where most discs present as grade~IV or~V and
clinically meaningful variation is compressed into the top of the scale.
Griffith et al.\ addressed this with an eight-level modified scale supported by a
24-image reference panel, validated on 260 discs in 52 subjects of mean age 73;
intra-observer agreement was excellent (weighted $\kappa$ 0.79--0.91) and
inter-observer agreement substantial (0.65--0.67), with most disagreement
confined to a single grade \cite{griffith2007modified}.

Esposito et al.\ place these two systems in a wider context that is worth
recording: a scoping review of 569 studies identified \emph{93} distinct grading
systems for lumbar disc degeneration---63 subjective, 25 quantitative---of which
the Pfirrmann method accounted for just over half of all reported use
\cite{esposito2025mribased}. The apparent standardisation of the field is
therefore partly illusory, and comparisons of automated grading accuracy across
studies must confirm that the same instrument was being reproduced.

\subsection{Central Canal Stenosis: Schizas and Lee}
\label{sec:lr-central-grading}

Two morphological systems have largely displaced simple dural sac
cross-sectional area (DSCA) thresholds for central canal assessment, on the
grounds that area measurements correlate poorly with symptoms and overlap
substantially between symptomatic and asymptomatic individuals
\cite{schizas2010qualitative}.

The \textbf{Schizas classification} grades the morphology of the dural sac on
axial T2 images by the ratio of cerebrospinal fluid to nerve rootlets. Grades~A
and~B retain visible CSF---in~A the rootlets occupy part of the sac, in~B they
fill it but remain individually discernible, giving a characteristic grainy
appearance. Grades~C and~D show no discernible rootlets and no CSF signal,
distinguished from one another by whether posterior epidural fat is preserved
(C) or obliterated (D) \cite{schizas2010qualitative}. Applied to 95 subjects
across surgical, conservatively managed and low-back-pain groups, intra- and
inter-observer agreement were substantial and moderate respectively
($\kappa = 0.65$ and $0.44$). Critically, area measurement over-diagnosed
stenosis in 35 patients and under-diagnosed it in 12 relative to the
morphological grade, and grades~C and~D predicted failure of conservative
management \cite{schizas2010qualitative}.

The \textbf{Lee central canal system} offers a simpler four-grade alternative
based on the degree of separation of the cauda equina rootlets on axial T2:
grade~0 no stenosis; grade~1 obliteration of the anterior CSF space with all
rootlets still separated; grade~2 some rootlets aggregated; grade~3 none
separated \cite{lee2011a}. Assessed by four musculoskeletal radiologists across
81 patients, it achieved inter-reader ICC of 0.730--0.953 and intra-reader
$\kappa$ of 0.863--0.900, and grades tracked cross-sectional area and
anteroposterior diameter at nearly all levels---notably better agreement than
the Schizas system, at the cost of finer morphological detail.

\subsection{Neural Foraminal Stenosis: Lee and Sartoretti}
\label{sec:lr-foraminal-grading}

The \textbf{Lee foraminal system} grades parasagittal images by the pattern of
perineural fat obliteration and the morphology of the nerve root: grade~0
normal; grade~1 fat obliteration in two opposing directions; grade~2
obliteration in all four directions without morphological change to the root;
grade~3 nerve root collapse or distortion \cite{lee2010a}. Evaluated on 576
foramina in 96 patients from L3--L4 to L5--S1 by two radiologists, inter- and
intra-observer agreement were near-perfect ($\kappa$ approximately 0.80--1.0).
This is the direct clinical precursor of the left and right foraminal narrowing
targets in the benchmark used in this thesis. The Sartoretti classification
offers a six-point sagittal alternative (grades~A--F) indexed to progressive
contact between the exiting root and each foraminal boundary.

The near-perfect agreement Lee et al.\ report deserves scepticism when
generalised. It was obtained by two experienced readers applying explicit
criteria to a curated series; the multi-reader, multi-institution agreement
figures in Section~\ref{sec:lr-reliability} are considerably lower, and it is
the latter that describe the conditions under which large training corpora are
actually annotated.

\subsection{Subarticular and Lateral Recess Stenosis}
\label{sec:lr-subarticular-grading}

Subarticular stenosis is graded on axial T2 by the proportion of the lateral
recess compromised: mild at up to one third, moderate between one and two
thirds, severe beyond two thirds, with nerve root involvement described as
touching, displacing or compressing \cite{lurie2008reliability,andreisek2014consensus}.
Compared with the disc and central canal systems this is a notably less
crisply specified instrument, resting on a visually estimated ratio rather than
a discrete morphological pattern. That imprecision shows in the reliability
data: subarticular stenosis is consistently the least reproducible of the
targets \cite{lurie2008reliability}, and it is also the target on which the
severity distribution in LumbarDISC is least skewed
(Section~\ref{sec:lr-lumbardisc}). Its difficulty is thus intrinsic, not merely
a consequence of scarce severe examples.

\subsection{Disc Herniation: Nomenclature and the MSU Classification}
\label{sec:lr-herniation-grading}

Herniation is described qualitatively by the standardised nomenclature of
bulge, protrusion, extrusion and sequestration \cite{fardon2014lumbar}, and
quantitatively by the Michigan State University (MSU) classification, which
grades size against a transverse line joining the medial edges of the facet
joints on the axial slice of maximal herniation \cite{mysliwiec2010msu}.
Grade~1 extends less than half the distance to that line, grade~2 more than
half without crossing it, grade~3 completely beyond it. The MSU system supplies
the reference standard for several automated herniation studies
\cite{zhang2023deep}, and its size~1 versus size~2/3 boundary maps onto the
conservative versus surgical decision.

\subsection{Inter- and Intra-Observer Reliability}
\label{sec:lr-reliability}

The most consequential study for present purposes is that of Lurie et al., who
assessed standardised MRI interpretation of lumbar stenosis using 58 randomly
selected studies from the SPORT trial read by four independent readers, with 20
re-read after at least a month \cite{lurie2008reliability}. Inter-reader
agreement was substantial for central stenosis ($\kappa = 0.73$) but only
moderate for foraminal stenosis ($0.58$) and nerve root impingement ($0.51$),
and poorest for subarticular stenosis ($0.49$), with marked variation between
reader pairs. Quantitative measurement fared no better in absolute terms: mean
absolute inter-reader difference in thecal sac area was 128~mm$^2$, or 13\%.
Intra-reader agreement exceeded inter-reader agreement for every feature.

Two structural findings recur across the wider reliability literature. First,
agreement is systematically worse for the lateral compartments than for the
central canal---the same ordering that later appears in automated model
performance. Second, disagreement is overwhelmingly confined to adjacent
grades: Griffith et al.\ report 84\% of intra-observer and 91\% of
inter-observer disagreement spanning a single grade
\cite{griffith2007modified}, and the same pattern recurs in automated
evaluation \cite{mcsweeney2023external,udomluck2026multitask}. Earlier consensus
efforts confirm that the underlying definitions were themselves contested: the
Delphi survey of Mamisch et al.\ \cite{mamisch2012radiologic} and the consensus
conference reported by Andreisek et al.\ \cite{andreisek2014consensus} were both
convened precisely because radiologic criteria for stenosis varied between
institutions.

\subsection{Label Noise as a Performance Ceiling}
\label{sec:lr-label-ceiling}

The reliability figures above are not a peripheral caveat; they are a hard
constraint on the interpretation of every result in this field. If four expert
readers agree on subarticular severity at $\kappa = 0.49$, then a model trained
on one reader's labels and evaluated against another's cannot exceed that
agreement, and a reported accuracy approaching unity on such a target should
prompt scrutiny of the evaluation protocol rather than admiration. The
LumbarDISC dataset description states the point explicitly, noting that
published inter-rater agreement on these grades ranges from poor to substantial
depending on anatomical location \cite{richards2026the}.

The field has responded to this in two ways. The first is to make the reference
standard more robust by construction---majority voting across a panel with
adjudication of disagreements, as in Tumko et al.\ (seven radiologists)
\cite{tumko2024a} and Su et al.\ (two clinicians with expert adjudication)
\cite{su2022automatic}. The second is to reframe the evaluation question: rather
than asking whether a model matches a single reader, asking whether its
disagreement with readers falls inside the band of disagreement between readers.
Won et al.\ furnish the cleanest example, reporting that the difference between
their two experts was statistically significant while the difference between
each expert and the model trained on that expert's labels was not
\cite{won2020spinal}. Comparatively few studies adopt either safeguard, and
Section~\ref{sec:lr-reporting} returns to this as a systemic weakness.

\subsection{The Ordinal Structure of the Label Space}
\label{sec:lr-ordinality}

Severity labels are ordinal, not categorical: \emph{normal}, \emph{mild},
\emph{moderate} and \emph{severe} stand in a fixed order, and the distance
between them is neither uniform nor clinically symmetric. Treating them as
unordered classes under categorical cross-entropy---as the majority of the
reviewed work does---discards that structure, penalising a
severe-graded-as-normal error exactly as heavily as a moderate-graded-as-severe
one.

Two empirical observations follow. First, errors concentrate overwhelmingly at
adjacent grade boundaries. McSweeney et al.\ report that SpineNet disagreed with
human raters on 20.83\% of discs but that only 0.85\% of discs differed by more
than one grade \cite{mcsweeney2023external}; Udomluck et al.\ observe the same
concentration \cite{udomluck2026multitask}; Ishimoto et al.\ find exact-grade
agreement of 65.7\% rising to 94.1\% when collapsed to a severe-versus-rest
dichotomy \cite{ishimoto2020could}. Much of the apparent error in ordinal
grading is therefore boundary error rather than gross misclassification.

Second, and less expectedly, explicitly ordinal objectives have not reliably
improved on this. Niemeyer et al.\ tested soft-kappa loss, ordinal
cross-entropy and regression losses against plain cross-entropy for Pfirrmann
grading on 1{,}599 patients and 7{,}948 discs, together with class pooling and
class-weighted losses, and found that none improved overall classification
performance; the default cross-entropy classifier reached $\kappa = 0.92$ with
predictions within one grade of ground truth in 99.2\% of cases
\cite{niemeyer2021a}. Patil et al.\ likewise found performance strongest at the
extremes of the Pfirrmann scale and on the binary low-versus-high distinction,
with intermediate grades remaining the difficulty \cite{patil2026mribased}. The
ordinal structure of the problem is thus real and measurable, but the question
of how to exploit it is genuinely open---a point that informs the loss functions
adopted in Chapter~3.

% ===========================================================================
\section{Evolution of Automated Lumbar MRI Analysis}
\label{sec:lr-evolution}

\subsection{Handcrafted Features and Classical Machine Learning}
\label{sec:lr-classical}

Before deep learning, automated analysis of lumbar MRI depended on features
designed by hand. Investigators specified mathematical descriptors of local
appearance and passed them to a conventional classifier. Ghosh et al.\ combined
Histogram of Oriented Gradients (HOG) descriptors with a support vector machine
and a set of geometric heuristics, reporting 99\% disc localisation accuracy
across 53 clinical cases and 318 discs, and---unusually for the period---emitting
a tight bounding box for each disc rather than a single interior point, so that
a downstream system could operate on the region directly
\cite{ghosh2012a}. Athertya and Kumar took a comparable approach to vertebral
degeneration, extracting local binary patterns, Hu moments and gray-level
co-occurrence matrix (GLCM) statistics to classify Modic changes, endplate
defects and focal changes, reaching 85.91\% accuracy for the extent of Modic
change on 100 patients \cite{athertya2021classification}. Related efforts applied
hybrid classical models to disc abnormality detection
\cite{unal2015detection,castromateos2016intervertebral} and semi-automatic
computer-aided classification of degenerative disease \cite{ruizespaa2015semiautomatic},
while Huber et al.\ pursued texture analysis with decision trees as a
quantitative alternative to visual stenosis grading \cite{huber2019qualitative}.
He et al.\ applied a synchronised superpixel representation to define consistent
regions of interest across levels before classifying neural foraminal stenosis,
an early attack on one of the exact conditions graded in this thesis
\cite{he2016automated}.

These systems established feasibility but generalised poorly. Handcrafted
descriptors are sensitive to scanner hardware, field strength, slice thickness
and patient positioning, and the anatomical variability of the lumbar spine
defeats fixed geometric assumptions. One finding from this era nevertheless
retains value: Athertya and Kumar's feature-sensitivity analysis showed that
GLCM entropy alone sufficed to detect focal changes, while endplate-defect
classification depended almost entirely on the first two Hu moments
\cite{athertya2021classification}. Such explicit attribution of predictive power
to individual features is something the deep models that displaced these methods
recovered only later, and imperfectly, through the interpretability techniques of
Section~\ref{sec:lr-interpretability}.

\subsection{The First End-to-End Deep Learning Systems}
\label{sec:lr-first-dl}

Convolutional neural networks removed the need for explicit feature design by
learning hierarchical representations directly from pixels \cite{he2016deep}. In
lumbar imaging the decisive contribution was SpineNet
\cite{jamaludin2017spinenet}. Trained on 2{,}009 patients and 12{,}018 disc
volumes from sagittal T2 images, it predicted six radiological scores
simultaneously---Pfirrmann grade, disc narrowing, upper and lower endplate
defects, and upper and lower marrow changes---from a shared convolutional trunk
that branched into task-specific heads. Two design decisions proved durable.
First, multi-task learning let one architecture serve several targets, reducing
the free parameters relative to training separate networks and exploiting the
correlation between co-occurring pathologies. Second, the system required only
disc-level class labels: no voxel-wise segmentation and no slice-level
localisation, with the evidence hotspots emerging implicitly from classification
training. Reported performance was state-of-the-art and near-human across all
six scores.

The same group demonstrated shortly afterwards that longitudinal clinical scans
could serve as free supervision. A Siamese network was pre-trained with a
contrastive loss on whether two scans belonged to the same patient---information
available from DICOM metadata without knowing the patient's identity---together
with an auxiliary vertebral-level prediction task, then fine-tuned for
degeneration grading. On 1{,}016 subjects, 423 with follow-up imaging, the
pre-trained network outperformed training from scratch and reached equivalent
performance from far fewer annotated examples \cite{jamaludin2017selfsupervised}.
This anticipated by several years the self-supervised pretraining strategies now
routine in medical imaging, and it is the direct ancestor of the contrastive
approach discussed in Section~\ref{sec:lr-contrastive}.

SpineNet's principal limitation was structural rather than architectural: it
consumed preprocessed sagittal disc volumes, and could therefore not assess
pathologies that require the transverse plane. SpineNetV2 extended coverage to
eleven radiological features and moved to a multi-stage design, using a U-Net to
detect and label vertebral bodies and discs before a ResNet classifier graded
each \cite{windsor2022spinenetv2}. The reliance on sagittal input nonetheless
persisted, and Section~\ref{sec:lr-generalization} shows this to be the
consistent locus of its external-validation failures.

\subsection{From Binary Detection to Multi-Level, Multi-Pathology Grading}
\label{sec:lr-task-evolution}

The task definition itself has evolved, and the trajectory matters more than any
individual model. Early deep learning work asked binary questions: is stenosis
present or absent \cite{altun2023lssvgg16}; is a herniation visible in this image
\cite{tsai2021lumbar}. Such formulations are tractable but clinically thin,
because management depends on severity and on which level is affected.

Three shifts followed. The first was to \emph{multi-class severity}: Won et al.\
graded stenosis status on 542 L4--5 axial images against two independent experts
\cite{won2020spinal}; Niemeyer et al.\ predicted the full five-point Pfirrmann
scale across 7{,}948 discs \cite{niemeyer2021a}. The second was to
\emph{multi-pathology} output, in which one model reports several findings at
once: Lehnen et al.\ evaluated disc herniation, bulging, canal stenosis,
nerve-root compression and spondylolisthesis across 888 lumbar segments from a
single CNN \cite{lehnen2021detection}; Su et al.\ graded herniation, central
canal stenosis and nerve-root compression from a shared ResNet-50 trunk with
three classification heads \cite{su2022automatic}. The third was to
\emph{per-level, per-side} prediction, in which the unit of analysis is not the
study or the image but the individual disc level and, for the lateral
compartments, the individual side. Hallinan et al.\ exemplify the mature form,
grading central canal, lateral recess and neural foraminal stenosis at each
level with the sequence appropriate to each \cite{hallinan2021deep}. Single-stage
detectors have since been adapted to deliver the same output more cheaply: Wang
et al.\ modified a YOLOv5 network to detect regions of interest and grade
central, lateral recess and foraminal stenosis simultaneously
\cite{wang2025a}, and a related design grades Pfirrmann degeneration, herniation
and high-intensity zones from sagittal and axial input in one pass
\cite{wang2025automated}. Minin et al.\ pursued the same objective for Pfirrmann
grading with an explicitly open model, motivated by the inaccessibility of
proprietary graders \cite{minin2025spinescan}.

This third formulation is precisely that of the benchmark used in this thesis,
which requires twenty-five graded outputs per study: five conditions across five
disc levels \cite{rsna2024rsna,richards2026the}. The progression is therefore not
merely historical. It explains why performance figures from the binary era
cannot be compared with contemporary results, and why a model that classifies a
whole image is solving a different and substantially easier problem than one
that must attribute severity to a specific anatomical location.

\subsection{Chronological Synthesis}
\label{sec:lr-chronology}

Four phases are discernible. Until roughly 2016 the field was dominated by
handcrafted features and classical classifiers, with localisation as the central
problem \cite{ghosh2012a,ghosh2014supervised}. From 2017 to 2020, CNNs displaced
these methods and multi-task grading became feasible
\cite{jamaludin2017spinenet,lu2018deepspine,han2018spinegan}. From 2021 to 2023
the emphasis moved to clinical credibility: external validation
\cite{mcsweeney2023external}, interpretability \cite{bharadwaj2023deep,gao2022automatic},
reader-assistance evaluation \cite{lim2022improved} and multi-centre testing
\cite{yi2023deep,su2022automatic}. From 2024 onward, the release of a large
public benchmark \cite{richards2026the} redirected effort toward anatomically
localised, multi-view severity grading, alongside the arrival of transformer and
vision-language architectures
\cite{batra2025mscan,nguyen2026crossspine,islamsk2026an,zeng2025gpt4lfs}.

The through-line is a steady increase in the specificity of what is being
predicted, accompanied---until recently---by a decrease in the rigour with which
it is validated. Both systematic reviews of the period make the same complaint:
Compte et al.\ found that algorithms performed markedly worse in replication or
external validation than in their development studies \cite{compte2023are}, and
Wang et al.\ concluded that despite pooled sensitivity of 0.84 and specificity of
0.87 across twelve studies and 15{,}044 patients, no model was reliable enough
for clinical practice \cite{wang2024machine}. Section~\ref{sec:lr-validation}
takes up this tension directly.

% ===========================================================================
\section{Anatomical Localization as a Design Principle}
\label{sec:lr-localization}

\subsection{Disc and Vertebra Detection, Labelling, and Level Assignment}
\label{sec:lr-detection}

Before severity can be attributed to L4--L5, the system must know which structure
is L4--L5. Level assignment is thus a precondition of per-level grading, and it
has been treated as a problem in its own right since the earliest work
\cite{ghosh2012a,lootus2015automated}. Pan et al.\ illustrate the mature
solution: a Faster R-CNN locates vertebral bodies L1--S1 on sagittal images, the
geometric relationship between sagittal and axial acquisitions identifies the
corresponding disc in each axial slice, and the disc region of interest is then
extracted---each of the three stages reported at 100\% accuracy, with the residual
difficulty falling entirely on the subsequent classification
\cite{pan2021automatically}. Detection has, in short, largely been solved;
grading has not.

\subsection{Semantic and Instance Segmentation of Lumbar Structures}
\label{sec:lr-segmentation}

Dense segmentation provides a stronger anatomical substrate than bounding boxes.
The U-Net architecture \cite{ronneberger2015unet} and its descendants dominate,
as a dedicated systematic review of deep-learning-assisted disc segmentation
confirms \cite{wang2024deep}.
Ghosh and Chaudhary segmented vertebrae, disc, dural sac and background jointly
using cascaded random forests over image and sparsely sampled context features,
achieving Dice 0.87 and 0.84 for dural sac and disc across 212 clinical cases
\cite{ghosh2014supervised}. Han et al.\ introduced adversarial training for
multi-structure segmentation \cite{han2018spinegan}, S\'aenz-Gamboa et al.\
established CNN semantic segmentation of the structural elements surrounding the
spinal cord \cite{senzgamboa2021automatic}, and Al-Kafri et al.\ applied
SegNet to boundary delineation for stenosis assessment across 515 symptomatic
patients, contributing two purpose-built metrics---confidence and
consistency---for assessing the quality of the ground truth itself rather than
only the model \cite{alkafri2019boundary}.

Subsequent work has largely consisted of architectural refinement. Wang et al.\
added multilevel feature-map attention and residual blocks to an Attention U-Net
\cite{wang2022automatic}; Ahmed et al.\ targeted class imbalance between
anatomical structures with a custom combined loss and modified upsampling
\cite{ahmed2024pioneering}; Silveira et al.\ added an Inception module for
multi-scale extraction and a dual-output mechanism, reporting mIoU 0.8974 and
F1 0.9444 on SPIDER and---notably---outperforming TransUNet
\cite{silveira2025automated}. M\"oller et al.\ extended coverage to the whole
spine with SPINEPS, segmenting fourteen structures including the posterior
elements, and established a useful general result: segmenting semantically first
and instance-wise second outperformed training directly on instance segmentation,
beating an nnU-Net baseline \cite{isensee2021nnunet} on every structure and metric
\cite{mller2025spinepsautomatic}. Chen et al.\ merged convolutional and attention
pathways in parallel with a novel position embedding, outperforming fifteen
competing models \cite{chen2024symtc}.

\subsection{Public Segmentation Resources and Benchmarks}
\label{sec:lr-segmentation-data}

Progress in this subfield has been unusually well served by shared data. The
SPIDER dataset provides 447 sagittal T1 and T2 series from 218 patients across
four hospitals, with reference segmentations of vertebrae, discs and spinal
canal plus per-level radiological gradings, together with reference performance
for a baseline algorithm and nnU-Net and a continuous public challenge on a
sequestered test split \cite{vandergraaf2024lumbar}. Several of the segmentation
models above are trained and benchmarked on it
\cite{ahmed2024pioneering,silveira2025automated,patil2026mribased}, which makes
their results mutually comparable in a way that is rare elsewhere in this
literature.

\subsection{Quantitative Morphometry}
\label{sec:lr-morphometry}

Segmentation enables measurement, and measurement offers an escape from ordinal
subjectivity. The dural sac cross-sectional area is the canonical example.
Ghobrial and Roth compared U-Net, Attention U-Net and MultiResUNet for automated
DSCA quantification on 683 scans with 50 held back for external validation;
MultiResUNet performed best, with Pearson correlation 0.9917 and mean absolute
error 23.70~mm$^2$, holding up externally at 99.95\% accuracy and F1 0.9393
\cite{ghobrial2025deep}. Zheng et al.\ pursued the same objective for disc
degeneration, segmenting degeneration-related regions with a Swin-transformer
network and computing signal-intensity and geometric features, then proposing a
quantitative criterion to replace subjective Pfirrmann grading on the basis of
strong correlation with grade across a large population \cite{zheng2022deep}.
Related work measures the lordosis angle from segmentation masks
\cite{ji2026mdvmunet}.

A caution is warranted. Ghobrial and Roth's model relies on T1 weighting alone
and a modest dataset, as the authors acknowledge \cite{ghobrial2025deep}, and
Schizas' original finding---that area measurement over- and under-diagnosed
stenosis relative to morphological grade in a substantial minority of patients
\cite{schizas2010qualitative}---applies equally whether the area is traced by
hand or by network. Automating a measurement does not repair a weak correlation
between that measurement and the clinical state.

\subsection{Evidence that Localisation Precedes Classification}
\label{sec:lr-localize-first}

The strongest recurring methodological finding in this literature is that
explicit anatomical localisation before classification improves grading. The
argument was made structurally by DeepSPINE, which chained U-Net vertebral
segmentation and disc-level designation to a multi-task grading network across
4{,}075 patients and more than 22{,}000 disc-level annotations
\cite{lu2018deepspine}, and it has been corroborated repeatedly since.

Four lines of evidence converge. Bharadwaj et al.\ segmented dural sac and disc
with V-Net, localised facet and foramen by geometric rule, and only then
classified, reaching Dice 0.93 and 0.94 with foramen and facet localisation IoU
0.72 and 0.83 \cite{bharadwaj2023deep}. \v{S}u\v{s}ter\v{s}i\v{c} et al.\ showed
that segmenting, cropping and enhancing the region of interest before
classification produced 0.87 and 0.91 accuracy on axial and sagittal views for a
multi-class herniation problem \cite{sustersic2022a}. Kowlagi et al.\ made the
comparative case most directly, arguing that a well-tuned three-stage pipeline of
segmentation, localisation and classification allows convolutional networks to
outperform transformer approaches on Pfirrmann grading in a population cohort
\cite{kowlagi2022a}. Acharya and Kansakar tested the principle in its strongest
form, hypothesising that anatomical focus rather than architecture is the binding
constraint, and grading severity from a single localised region of interest per
disc \cite{acharya2026disccentric}.

The implication for the present work is direct. If localisation dominates
architecture in determining performance, then a comparison between backbones is
only meaningful when both receive identically localised input---otherwise the
comparison measures preprocessing rather than representation. This is a
methodological requirement that Section~\ref{sec:lr-cnn-vs-transformer} shows
much of the existing comparative literature fails to satisfy.

\subsection{Reduced-Annotation and Unsupervised Alternatives}
\label{sec:lr-weak-supervision}

Dense annotation is expensive, and several strands of work reduce the
requirement. Kuang et al.\ eliminated manual labels entirely for vertebral
segmentation: sub-optimal masks from a rule-based method were combined through a
voting mechanism to supervise CNN training, with a regional strategy restricting
attention to a specific vertebral region, validated across 243 volunteers
\cite{kuang2020mrisegflow}. SpineNet's evidence hotspots represent the
complementary weakly supervised route, obtaining localisation from
classification labels alone \cite{jamaludin2017spinenet}, while the SPIDER
annotation protocol adopted a human-in-the-loop compromise, training on a small
subset to produce initial masks that were then corrected and fed back
\cite{vandergraaf2024lumbar}. Self-supervised pretraining on unlabelled or
metadata-labelled data \cite{jamaludin2017selfsupervised} completes the picture.
Together these establish that the annotation burden, long cited as the principal
obstacle to clinical deployment, is substantially negotiable.

% ===========================================================================
\section{Backbone Architectures for Severity Classification}
\label{sec:lr-backbones}

\subsection{Convolutional Networks}
\label{sec:lr-cnns}

Convolutional networks exploit two properties well matched to imaging: local
connectivity, since a filter responds to a neighbourhood rather than the whole
field; and translation equivariance, since the same filter is applied at every
position. Stacking convolutions builds a hierarchy in which early layers respond
to edges and texture and later layers to composite anatomical structure.

\textbf{ResNet} addressed the degradation encountered when networks are made
deeper, introducing residual blocks whose skip connections allow a layer's input
to bypass subsequent convolutions and be added to their output, providing an
unimpeded path for gradients and permitting stable training at fifty layers and
beyond \cite{he2016deep}. In lumbar imaging ResNet is ubiquitous: as the
classification head in multi-task grading \cite{su2022automatic}, as the
detection backbone in staged pipelines \cite{zhang2023deep}, as the encoder in
segmentation networks \cite{batra2025mscan}, and as the classifier in
SpineNetV2 \cite{windsor2022spinenetv2}.

\textbf{EfficientNet} attacked a different problem---how to allocate a fixed
computational budget---by scaling depth, width and input resolution together
under fixed coefficients derived from a neural architecture search, rather than
scaling one dimension in isolation \cite{tan2019efficientnet}. The resulting
models reach comparable accuracy at substantially lower parameter and FLOP cost,
which matters for the deployment scenarios that motivate this thesis.
EfficientNet appears as the classifier in Kowlagi et al.'s Pfirrmann pipeline
\cite{kowlagi2022a} and in the final stage of M-SCAN \cite{batra2025mscan}.
\textbf{ConvNeXt} represents the most recent convolutional line, modernising the
CNN with design choices borrowed from transformers while retaining convolution
\cite{liu2022convnet}; Udomluck et al.\ include ConvNeXt-Tiny in their feature
extractor comparison and find it stable across classifiers
\cite{udomluck2026multitask}.

The architectural limitation of convolution is a bounded receptive field. Deep
stacks aggregate local responses into progressively broader ones, but the
network never explicitly models a relationship between two distant regions---for
instance, between the L1--L2 and L5--S1 levels, whose degenerative states are
biomechanically coupled.

\subsection{Vision Transformers and the Attention Bottleneck}
\label{sec:lr-vit}

Transformers dispense with convolution and rely on self-attention, in which
every element of the input is weighted against every other regardless of
distance \cite{vaswani2017attention}. The Vision Transformer applies this by
partitioning an image into fixed patches, flattening them and treating the
result as a sequence \cite{dosovitskiy2021image}. Global context is available
from the first layer, and there is no locality prior to overcome---but also no
locality prior to exploit, which is why vision transformers are markedly more
data-hungry than convolutional networks of comparable capacity.

The practical obstacle is cost. Global self-attention scales quadratically with
the number of patches, hence quadratically with image area. Clinical MRI demands
high spatial resolution to resolve structures a few millimetres across, so a
naive vision transformer over full-resolution lumbar images is computationally
prohibitive.

\subsection{The Swin Transformer}
\label{sec:lr-swin}

The Swin Transformer resolves this bottleneck through two mechanisms
\cite{liu2021swin}. \emph{Windowed self-attention} restricts attention to
non-overlapping local windows, making cost linear rather than quadratic in image
size; to restore cross-window communication, the window partition is
\emph{shifted} by half a window between consecutive layers, so information
propagates across boundaries over successive blocks. \emph{Hierarchical feature
maps} are produced by patch-merging layers that progressively reduce spatial
resolution while increasing channel depth, yielding a pyramid analogous to a
CNN's and making the backbone usable for dense prediction as well as
classification.

For lumbar grading the appeal is specific. Severity depends simultaneously on
fine local detail---the texture of a nucleus pulposus, the obliteration of
perineural fat---and on regional context, since a rootlet configuration is
interpretable only relative to the surrounding dural sac and the levels above
and below. A hierarchical, windowed attention mechanism is a plausible fit to
that structure, which is the substantive reason for selecting it as the
transformer representative in this thesis. Swin-derived architectures have been
applied to spinal segmentation, and Zheng et al.\ incorporate Swin components in
their quantitative degeneration network \cite{zheng2022deep}.

\subsection{Hybrid CNN--Transformer Architectures}
\label{sec:lr-hybrids}

A third position holds that the two families are complementary. SymTC merges
convolutional and transformer layers in a \emph{parallel dual-path} arrangement
rather than stacking them sequentially, and integrates a position embedding into
the self-attention module to improve spatial awareness; benchmarked against
fifteen existing models on both a private dataset and a released synthetic one,
it performed best for both vertebral bones and discs \cite{chen2024symtc}. Zheng
et al.\ similarly combine convolutional and Swin components
\cite{zheng2022deep}, and Yi et al.\ pair a 3D ResNet18 backbone with transformer
components across three hospitals \cite{yi2023deep}. Baur et al.\ explore a
different pairing, segmenting discs with a U-Net, reconstructing each as a 3D
point cloud and grading it with a graph neural network; performance was uneven
across levels---F1 0.71 at L2/3 against 0.46 overall---which the authors present
as a foundation requiring refinement rather than a competitive result
\cite{baur2024automated}. Hybrids complicate the
clean comparison this thesis undertakes, but they indicate where the field
expects the answer to lie: not in the wholesale replacement of convolution, but
in supplying it with a mechanism for long-range interaction.

\subsection{Published CNN-versus-Transformer Comparisons and Their Limits}
\label{sec:lr-cnn-vs-transformer}

Several studies bear directly on the architectural question, and their results
do not favour transformers as strongly as the general computer vision literature
might suggest.

Kowlagi et al.\ address it explicitly. Observing that recent gains had been
attributed to transformer architectures, they show that a well-tuned three-stage
convolutional pipeline outperforms the state-of-the-art transformer approaches
on Pfirrmann grading, evaluated on the Northern Finland Birth Cohort 1966---a
population cohort rather than a clinical convenience sample---with an ablation
study and subgroup breakdown, and public code \cite{kowlagi2022a}. Their
conclusion is that pipeline design and tuning outweigh backbone choice.

Lee et al.\ provide the closest published analogue to the present study,
developing a CNN-based and a transformer-based model for lumbar stenosis
classification on 564 MRIs with a 100-study external test set, labelled by four
radiologists with expert consensus as reference and read by eight participants
across four experience levels. Both models exceeded 94\% detection recall; on
dichotomised classification the CNN, transformer and human participants achieved
kappas of 0.99/0.99/0.97--0.98 for central canal, 0.98/0.94/0.81--0.94 for
lateral recesses and 0.98/0.95/0.91--0.95 for neural foramina. The CNN was
superior overall, the transformer similar to superior against clinicians
\cite{lee2026deep}.

Udomluck et al.\ compared VGG19, ConvNeXt-Tiny and DINOv2 features with classical
classifiers under external validation, finding conventional VGG19 features
strongest (accuracy 0.9556 with logistic regression) and DINOv2 features
generalising worst, dropping to 0.6222 with LightGBM
\cite{udomluck2026multitask}. Patil et al.\ reached a still more pointed
conclusion through ablation: radiomic features carried essentially the entire
signal for Pfirrmann grading, and frozen ImageNet-pretrained ResNet50 features
added no measurable discriminative value \cite{patil2026mribased}. Silveira et
al.\ likewise found their enhanced U-Net outperformed TransUNet on segmentation
\cite{silveira2025automated}.

Three methodological limitations recur across this evidence. First, the
compared architectures are rarely given identical preprocessing, localisation
and training budgets, so differences may reflect tuning effort rather than
representational capacity---exactly the confound Kowlagi et al.\ identify
\cite{kowlagi2022a}. Second, the input configuration is usually held fixed,
leaving unexamined the interaction between architecture and available context,
even though a transformer's advantage should plausibly be largest where global
context must substitute for missing views. Third, several comparisons rely on
frozen pretrained features rather than fine-tuned backbones
\cite{udomluck2026multitask,patil2026mribased}, which disadvantages transformers
disproportionately given their weaker inductive biases. These three gaps
motivate the experimental design of this thesis.

\subsection{Transfer Learning, Pretraining, and Self-Supervision}
\label{sec:lr-pretraining}

Medical datasets are small by the standards these architectures were designed
for, so initialisation matters. ImageNet pretraining is near-universal, though
Patil et al.'s ablation is a caution that its benefit is not automatic
\cite{patil2026mribased}. Domain-specific pretraining is better supported:
Jamaludin et al.\ obtained equivalent performance from far fewer labels through
longitudinal self-supervision \cite{jamaludin2017selfsupervised}, and Acharya and
Kansakar used contrastive pretraining on disc-centred crops to reduce sensitivity
to irrelevant appearance variation \cite{acharya2026disccentric,chen2020simple}.
Sequential transfer between related targets---initialising a subarticular model
from spinal canal weights---has also been reported to help on smaller datasets, a
strategy of particular relevance to transformers given their lack of
convolutional inductive bias.

% ===========================================================================
\section{Exploiting Spatial Context: Slices, Volumes, and Views}
\label{sec:lr-context}

\subsection{2D, 2.5D, and 3D Input Representations}
\label{sec:lr-dimensionality}

An MRI series is a volume, but severity is assigned per level, so a
representation must be chosen. \emph{2D} models classify individual slices and
are cheap but discard through-plane context; since a disc spans several slices
and a radiologist integrates across them, a single slice may not contain the
evidence. \emph{3D} models consume the volume directly \cite{yi2023deep} at
considerable memory cost---Acharya and Kansakar note that a standard
high-resolution series contains millions of voxels, making full-volume
processing impractical for many architectures \cite{acharya2026disccentric}.

\emph{2.5D} representations occupy the middle ground and have become the
default in this domain. A small stack of adjacent slices is treated as channels
or as a short sequence, capturing local through-plane context at roughly 2D
cost. SpineNet's input was a stack of nine slices per disc
\cite{jamaludin2017spinenet}; M-SCAN describes its approach explicitly as 2.5D,
selecting the three axial slices closest to each level and modelling them with
recurrent layers \cite{batra2025mscan}; the leading challenge solutions combined
2D CNNs with sequence models or attention-based multiple-instance learning over
the slice dimension.

\subsection{Cross-Sequence and Cross-Plane Fusion}
\label{sec:lr-fusion}

Where multiple sequences are available, they must be combined. The simplest
approach concatenates them as channels, which presumes spatial registration that
sagittal and axial acquisitions do not satisfy. More considered strategies
process each sequence separately and fuse at the feature level.

CrossSpine is the most explicit treatment. Nguyen et al.\ identify two failings
of baselines---single-stream designs cannot integrate complementary information
across sequences, and uniform treatment of all discs ignores level-specific
anatomy---and address the first with a cross-sequence attention mechanism that
fuses features at multiple spatial scales, letting each sequence attend to all
others, followed by a weighting mechanism that prioritises the most
diagnostically informative sequence. They report Macro F1 improved by more than
125\%, Mean AUPRC by 99\% and Mean AUROC by 36\% over their baseline
\cite{nguyen2026crossspine}. M-SCAN applies cross-attention to align sagittal and
axial feature maps before classification \cite{batra2025mscan}, and Park et al.\
pursue transformer-based multimodal fusion on the same benchmark
\cite{park2025multiview}. Zeng et al.\ extend fusion beyond imaging entirely,
combining ConvNeXt image features with GPT-4o-generated textual descriptions
encoded by RoBERTa through a Mamba fusion stage \cite{zeng2025gpt4lfs}.

\subsection{Geometric Correspondence Between Planes}
\label{sec:lr-geometry}

Fusing sagittal and axial data requires knowing which axial slice corresponds to
which sagittal level---a geometric problem solved through DICOM spatial
metadata rather than learned. Pan et al.\ derive the correspondence from the
geometric relationship between the two acquisitions, reporting 100\% accuracy in
identifying the disc in each axial slice \cite{pan2021automatically}. M-SCAN
projects 2D sagittal stenosis coordinates into 3D using patient orientation data
from the DICOM headers, then selects the three nearest axial slices for each of
the five levels \cite{batra2025mscan}.

This step is a hidden cost of multi-sequence modelling, and one that bears
directly on the trade-off examined here. It adds a preprocessing dependency, a
failure mode when metadata is absent or inconsistent, and an engineering burden
that a single-sequence pipeline simply does not incur.

\subsection{Multi-Sequence versus Single-Sequence Input}
\label{sec:lr-multiseq}

The clinical standard is unambiguous that all three sequences contribute
\cite{richards2026the,hallinan2021deep}, and the strongest reported results on
the RSNA benchmark use all three \cite{batra2025mscan}. The countervailing
consideration is cost: multi-sequence processing demands greater memory and
compute, cross-plane registration, and longer inference, and the resulting
infrastructure requirements exceed what many hospital environments provide.

The evidence on whether reduced input suffices is genuinely mixed, which is why
the question remains open. On one side, external validation of sagittal-only
SpineNetV2 found agreement dropping specifically for foraminal stenosis and
herniation---the pathologies requiring transverse context
\cite{nigru2024external}---and Wu et al.\ found ordinal Pfirrmann grading
degrading in upper lumbar discs while binary detection held up
\cite{wu2026external}. On the other, Acharya and Kansakar grade severity from
\emph{sagittal T2 alone}, reaching 78.1\% balanced accuracy with a
severe-to-normal misclassification rate of 2.13\%, on the argument that
anatomical focus rather than input breadth is the binding constraint
\cite{acharya2026disccentric}. Kowlagi et al.\ likewise operate on sagittal input
and outperform transformer baselines \cite{kowlagi2022a}.

Two observations sharpen the question. First, the comparison has not been made
under controlled conditions: the sagittal-only and multi-sequence results above
come from different architectures, datasets, localisation strategies and
evaluation protocols, so the contribution of sequence configuration cannot be
isolated. Second, the expected penalty is target-dependent---severe for
foraminal and subarticular grading, mild for central canal---so a single
aggregate accuracy figure would obscure precisely the structure of interest.
Both observations inform the design adopted in Chapter~3.

\subsection{Robustness to Missing or Incomplete Sequences}
\label{sec:lr-missing}

Clinical data is incomplete: sequences are omitted, truncated or corrupted by
artifact. The LumbarDISC inclusion criteria required all three sequences and
excluded studies with substantial artifact \cite{richards2026the}, so a model
trained on it has not been exposed to the degraded conditions of routine
practice. A model that requires three sequences fails outright when one is
absent, whereas a single-sequence model is unaffected---an argument for reduced
input configurations that is about robustness rather than efficiency, and one
that the literature on missing-modality learning addresses only in passing for
this application.

% ===========================================================================
\section{Class Imbalance and Ordinal Severity Modelling}
\label{sec:lr-imbalance}

\subsection{Severity Distributions in Lumbar MRI Datasets}
\label{sec:lr-distributions}

Class imbalance in this domain is not an artefact of sampling but a property of
the population: in any unselected clinical cohort, most disc levels are normal or
mildly degenerate, and severe compression is uncommon. The RSNA LumbarDISC
distributions make the scale explicit \cite{richards2026the}. Of 13{,}474 spinal
canal stenosis grades, 85.4\% are normal or mild, 8.8\% moderate and 5.9\%
severe. Of 26{,}919 neural foraminal grades, roughly 78\% are normal or mild,
18\% moderate and 4.5\% severe. Subarticular grading is the least skewed of the
three, with roughly 69\% normal or mild, 19\% moderate and 11\% severe across
26{,}285 grades.

The same pattern recurs elsewhere. Liawrungrueang et al.\ report a Pfirrmann
distribution of 220 grade~I, 530 grade~II, 170 grade~III, 160 grade~IV and just
20 grade~V across 1{,}000 images \cite{liawrungrueang2023automatic}; Niemeyer et
al.\ describe an approximately Gaussian grade distribution with grades~1 and~5
under-represented \cite{niemeyer2021a}. Two consequences follow. First, a model
trained under unweighted cross-entropy can achieve roughly 85\% accuracy on
spinal canal stenosis by predicting \emph{normal/mild} unconditionally, so
aggregate accuracy is close to uninformative on this task. Second, the ordering
of imbalance across targets is the inverse of the ordering of difficulty:
subarticular stenosis has the most balanced label distribution yet the poorest
human reliability \cite{lurie2008reliability}, confirming that its difficulty is
intrinsic rather than a consequence of scarce positive examples.

\subsection{Data-Level Strategies}
\label{sec:lr-data-level}

The most common data-level intervention is weighted random sampling, in which
the probability of drawing a minority-class example is inflated so that each
minibatch approaches a uniform severity distribution and gradients are not
dominated by the majority class. Aggressive augmentation of minority classes
complements this, since with few severe examples a network may memorise them
rather than learn the morphology of severe compression. Reported techniques
include contrast-limited adaptive histogram equalisation, random scaling, affine
rotation and elastic deformation, alongside mixing methods.

Tsai et al.\ provide the clearest single demonstration of augmentation's value
in the small-data regime, treating it as the subject of the study rather than an
implementation detail: mean average precision reached 92.4\% at 550 images with
augmentation, and augmentation was decisive in preventing both under- and
over-fitting \cite{tsai2021lumbar}. Ahmed et al.\ attack a related form of
imbalance---between anatomical structures rather than severity classes---through
preprocessing that rectifies class inconsistencies before training
\cite{ahmed2024pioneering}.

\subsection{Loss-Level Strategies}
\label{sec:lr-loss-level}

Class-weighted cross-entropy multiplies each class's contribution to the loss by
a scalar, forcing the optimiser to penalise minority-class errors more heavily.
The RSNA challenge encoded this directly into its evaluation metric, assigning
weights of 1, 2 and 4 to normal/mild, moderate and severe respectively, so that
aligning the training objective with the evaluation criterion is a matter of
adopting the same coefficients.

Focal loss goes further, introducing a modulating factor that down-weights
well-classified examples dynamically and concentrates gradient on hard,
borderline cases \cite{lin2017focal}. This is well matched to the present
problem, where the abundant normal cases are mostly easy and the difficulty
concentrates at the mild/moderate and moderate/severe boundaries---precisely
where human readers also disagree. Acharya and Kansakar combine weighted focal
loss with contrastive pretraining \cite{acharya2026disccentric}, and Islam~Sk et
al.\ propose an adaptive PID-Tversky loss that borrows feedback control to adjust
training penalties dynamically toward under-segmented minority instances
\cite{islamsk2026an}.

\subsection{Ordinal-Aware Objectives and Whether They Help}
\label{sec:lr-ordinal-losses}

Given the ordinal label structure of Section~\ref{sec:lr-ordinality}, it is
natural to expect ordinal-aware objectives to outperform categorical ones. The
evidence does not support this expectation as strongly as the intuition
suggests.

Niemeyer et al.\ conducted the most systematic test available, evaluating soft
kappa loss, ordinal cross-entropy and regression losses against standard
cross-entropy for Pfirrmann grading on 1{,}599 patients and 7{,}948 discs, with
extensive hyperparameter optimisation, and additionally trying class pooling and
class-weighted losses to counter imbalance. None of these improved classification
performance overall. The default cross-entropy classifier reached Cohen
$\kappa = 0.92$, average sensitivity 90.2\% and precision 92.5\%, with
predictions within one grade of ground truth in 99.2\% of cases and a mean
absolute error of 0.08 grades \cite{niemeyer2021a}. Patil et al.\ list ordinal
regression among the supplementary analyses they ran, and still report that
intermediate-grade discrimination remains the outstanding difficulty
\cite{patil2026mribased}.

Two readings are possible. Either the ordinal structure is already implicitly
captured---a network trained with cross-entropy on visually ordered classes will
naturally confuse adjacent grades more often than distant ones---or the specific
formulations tested were inadequate. The literature does not currently
distinguish these, which is why this thesis treats the choice of objective as an
empirical question rather than a settled one.

\subsection{Contrastive and Representation-Learning Approaches}
\label{sec:lr-contrastive}

A third strategy attacks imbalance in the representation rather than in the
sampling or the loss. Contrastive learning pulls examples of similar severity
together in the latent space and pushes dissimilar ones apart before any
classifier is fitted \cite{chen2020simple}, so the representation is shaped by
the structure of the problem rather than by class frequency.

Acharya and Kansakar apply this directly, combining contrastive pretraining with
disc-level fine-tuning over a single anatomically localised region of interest
per disc, on the reasoning that this forces attention onto intrinsic disc
features and reduces sensitivity to irrelevant appearance variation. An auxiliary
regression task handles disc localisation and a weighted focal loss addresses the
residual imbalance. The result---78.1\% balanced accuracy with the
severe-to-normal misclassification rate reduced to 2.13\% relative to supervised
training from scratch---is notable less for the headline accuracy than for what
it optimises \cite{acharya2026disccentric}. The lineage runs back to Jamaludin et
al.'s longitudinal self-supervision \cite{jamaludin2017selfsupervised}.

\subsection{The Clinical Asymmetry of Error}
\label{sec:lr-error-asymmetry}

The most important point in this section is that not all errors are equal, and
that standard metrics conceal the difference. Grading a severe stenosis as
normal risks a missed diagnosis and delayed decompression; grading a normal disc
as moderate risks an unnecessary review. Overall accuracy, macro-F1 and even
Cohen's $\kappa$ treat these identically.

Three studies take the asymmetry seriously. Acharya and Kansakar report the
severe-to-normal misclassification rate as a headline result alongside balanced
accuracy, arguing explicitly that conventional metrics do not capture the
differential clinical risk \cite{acharya2026disccentric}. Wu et al.\ characterise
SpineNetv2's error profile as specificity-oriented, with false negatives
exceeding false positives, and conclude that the system is suitable as a
confirmatory second reader but that reliance on its negative findings is not
recommended---a conclusion about deployment that follows from the shape of the
errors rather than their number \cite{wu2026external}. The RSNA weighted metric
encodes the same judgement into the benchmark itself.

Ishimoto et al.\ supply the complementary observation: collapsing four grades to
a severe-versus-rest dichotomy raised agreement from 65.7\% to 94.1\% with
$\kappa = 0.75$ \cite{ishimoto2020could}. If the clinically decisive question is
whether a level is severe, then performance on that question is considerably
better than the multi-class figures suggest, and evaluation should report both.
This chapter therefore recommends, and Chapter~3 adopts, reporting per-class
recall for the severe class alongside aggregate measures.

% ===========================================================================
\section{Datasets and Benchmarks}
\label{sec:lr-datasets}

\subsection{Private and Single-Institution Cohorts}
\label{sec:lr-private-data}

Most work before 2024 used private single-institution data, with cohort sizes
spanning three orders of magnitude---from 75 exams \cite{gao2022automatic} and
146 consecutive patients \cite{lehnen2021detection} through 200 studies
\cite{bharadwaj2023deep} and 542 images \cite{won2020spinal} to 4{,}075 patients
\cite{lu2018deepspine} and 15{,}254 axial images \cite{su2022automatic}.
Aggregate performance correlates only loosely with cohort size, because larger
private datasets are frequently drawn from a single scanner and protocol.

Two consequences follow, and both are structural rather than incidental. Results
are not comparable across studies, since each is measured against a different
label distribution and reference standard; and models cannot be independently
re-evaluated, since the data cannot be shared. Both systematic reviews of the
period identify this as the field's central methodological weakness
\cite{compte2023are,mendes2026deep}.

\subsection{Population and Epidemiological Cohorts}
\label{sec:lr-population-data}

A smaller body of work uses population cohorts, which differ from clinical
convenience samples in a way that matters for generalisation: participants are
not selected by referral, so the severity distribution reflects the population
rather than the imaged-because-symptomatic subset. The Northern Finland Birth
Cohort 1966, comprising lumbar MRI from participants imaged at age 46--47, is the
most used \cite{mcsweeney2023external,kowlagi2022a}; the Wakayama Spine Study
\cite{ishimoto2020could} and the Hong Kong Disc Degeneration Population-Based
Cohort \cite{cheung2022learningbased} serve similar roles. Kowlagi et al.\ are
explicit that evaluating on a population cohort rather than a clinical sample is
a deliberate methodological choice, reflecting the transferability of results to
downstream research \cite{kowlagi2022a}.

\subsection{Public Lumbar MRI Datasets}
\label{sec:lr-public-data}

Public data arrived late to this field relative to CT, where large annotated
vertebral collections have existed for years. SPIDER is the principal
segmentation resource: 447 sagittal T1 and T2 series from 218 patients across
four hospitals, with reference segmentations of vertebrae, discs and spinal
canal, per-level radiological gradings, baseline and nnU-Net reference
performance, and a continuous challenge on a sequestered split
\cite{vandergraaf2024lumbar}. Its influence is visible in the number of
subsequent methods benchmarked on it
\cite{ahmed2024pioneering,silveira2025automated,patil2026mribased,ghobrial2025deep}.
Smaller public collections support classification work
\cite{natalia2024lumbar,shahzadi2023nerve,nikpasand2024automated}.

\subsection{The RSNA 2024 Challenge and LumbarDISC}
\label{sec:lr-lumbardisc}

The RSNA Lumbar Degenerative Imaging Spine Classification dataset is the largest
and most diverse expertly annotated lumbar spondylosis collection assembled, and
it is the dataset used in this thesis \cite{richards2026the}. It comprises MRI
studies from 2{,}697 patients and 8{,}593 image series drawn from eight
institutions across six countries and five continents---a geographic spread
chosen deliberately so that models trained on it face genuine acquisition
heterogeneity rather than a single site's conventions.

Inclusion required a sagittal T2-like series (conventional spin echo T2, STIR or
Dixon), a sagittal T1 series and an axial T2 series covering at least three
lumbar disc levels, from outpatients aged 18 or over imaged for degenerative
disease. Studies with lumbosacral hardware, active non-degenerative pathology,
severe scoliosis or substantial artifact were excluded \cite{richards2026the}.
Annotation was performed by expert volunteer neuroradiologists and
musculoskeletal radiologists recruited through the ASNR
\cite{rsna2024rsna}. The associated challenge \cite{rsna2024rsna} made the
dataset the field's common benchmark and shifted attention decisively from
whole-image classification toward anatomically localised, per-level severity
grading.

\subsection{Label Structure, Distribution, and the Weighted Metric}
\label{sec:lr-label-structure}

The label structure is the defining feature of the benchmark and warrants precise
statement. Each study receives twenty-five graded outputs: five conditions
---spinal canal stenosis, left and right neural foraminal narrowing, left and
right subarticular stenosis---at each of five disc levels from L1/L2 to L5/S1.
Each is assigned one of three ordinal classes: Normal/Mild, Moderate or Severe.
Expert coordinate annotations accompany the grades, so anatomical localisation
can be learned in a supervised manner rather than inferred
\cite{rsna2024rsna,richards2026the}.

Three properties follow. The task is \emph{not} generic three-class image
classification, and results from studies that treat it as such are not
comparable. The severity distribution is severely imbalanced and differently so
for each condition (Section~\ref{sec:lr-distributions}). And the challenge metric
is a sample-weighted log loss with weights of 1, 2 and 4 across the three
classes, which encodes the clinical asymmetry of
Section~\ref{sec:lr-error-asymmetry} directly into the evaluation. A model tuned
for unweighted accuracy is therefore optimising a different objective from the
one the benchmark measures.

\subsection{Challenge Solutions and the Practices They Standardised}
\label{sec:lr-challenge-solutions}

The solutions converged on a common template, and that convergence is itself a
finding. Nearly all adopted multi-stage pipelines in which localisation precedes
classification: a detection model---commonly a U-Net, YOLO variant or Faster
R-CNN with a feature pyramid---identifies each disc level on the sagittal
series, and a second-stage classifier grades tightly cropped regions, typically
using 2.5D representations that combine a 2D backbone with a sequence model or
attention over adjacent slices.

M-SCAN exemplifies the approach and reports AUROC 0.971 on 1{,}975 studies,
using a U-Net with ResNet50 encoder to locate stenosis points on sagittal images,
projecting these into 3D via DICOM orientation metadata to select corresponding
axial slices, and applying cross-attention to align sagittal and axial features
\cite{batra2025mscan}. Liu et al.\ pair MobileNetV2 with U-Net for efficiency,
reporting 94.93\% accuracy under five-fold cross-validation
\cite{liu2025automated}; Park et al.\ apply transformer-based multimodal fusion
\cite{park2025multiview}; and further challenge-derived work continues to appear
\cite{patel2025automated,garciadecelis2025deep}.

Two cautions attach to these figures. Most are cross-validation results on the
public training split rather than external validation, so they measure fit to
LumbarDISC rather than transportability beyond it. And the top solutions rely on
large ensembles, which inflate reported accuracy relative to what a single
deployable model achieves---a gap directly relevant to the efficiency question
this thesis examines.

% ===========================================================================
\section{Evaluation, Validation, and Generalization}
\label{sec:lr-validation}

\subsection{Metrics and Their Appropriateness}
\label{sec:lr-metrics}

Reported metrics in this field are heterogeneous and frequently ill-matched to
the task. Overall accuracy remains the most common headline figure despite being
close to meaningless under the severity distributions of
Section~\ref{sec:lr-distributions}. Area under the ROC curve is widely reported
\cite{batra2025mscan,wang2024machine} but is defined for binary problems and must
be extended by one-versus-rest averaging for three-class grading, which conceals
which class is failing.

More appropriate choices appear in the better studies. Cohen's and Gwet's
$\kappa$ measure agreement above chance and are directly comparable with the
human reliability figures of Section~\ref{sec:lr-reliability}, which is the
correct reference point \cite{lim2022improved,su2022automatic,mcsweeney2023external}.
Balanced accuracy corrects for prevalence \cite{mcsweeney2023external,acharya2026disccentric}.
Mean absolute error in grades exploits the ordinal structure and distinguishes a
one-grade slip from a two-grade error \cite{niemeyer2021a,wu2026external}.
Ishimoto et al.\ report both exact-grade agreement and the collapsed
severe-versus-rest figure, making the difference visible rather than choosing
between them \cite{ishimoto2020could}.

The recommendation that follows is to report a small panel rather than a single
number: a prevalence-corrected aggregate, an agreement statistic comparable to
inter-reader values, per-class recall for the severe class, and the distribution
of error magnitudes across grade boundaries.

\subsection{Internal Validation and Optimistic Reporting}
\label{sec:lr-internal-validation}

A recurring pattern is that internally validated performance is high and
externally validated performance is materially lower. Several reported accuracies
in the reviewed literature exceed 95\%
\cite{shahzadi2023nerve,liawrungrueang2023automatic,liu2025automated,tabarestani2026a},
which sits uneasily beside inter-reader agreement of $\kappa = 0.49$ for
subarticular stenosis \cite{lurie2008reliability}. Where a model appears to
exceed the reproducibility of the process that generated its labels, the
explanation usually lies in the evaluation protocol: single-institution data,
cross-validation rather than held-out external testing, labels from a single
reader, or a task collapsed to a binary decision.

The point is not that these studies are wrong on their own terms, but that
their numbers do not transfer. Compte et al.\ state the general finding
directly: algorithms did not perform as well in replication or external
validation as in their development studies \cite{compte2023are}.

\subsection{External Validation and Domain Shift}
\label{sec:lr-generalization}

External validation on geographically and temporally separate cohorts is the
strongest available evidence of clinical viability, and a small number of
studies supply it.

McSweeney et al.\ evaluated SpineNet on 1{,}331 participants from the Northern
Finland Birth Cohort 1966, a dataset separate from its UK training data in both
geography and time. Balanced accuracy was 78\% for disc degeneration and 86\%
for Modic changes, with Pfirrmann reliability against expert radiologists of Lin
concordance 0.86 and Cohen $\kappa = 0.68$, and Modic detection at
$\kappa = 0.74$; both were essentially unchanged in a low-back-pain subset. The
informative detail is the error structure: 20.83\% of discs were graded
differently from human raters, but only 0.85\% differed by more than one grade
\cite{mcsweeney2023external}. Nigru et al.\ evaluated SpineNetV2 across eleven
features on 1{,}747 discs from 353 patients, finding agreement above 0.7 for most
targets but distinctly lower for foraminal stenosis and herniation, attributed
to the limits of sagittal input \cite{nigru2024external}. Wu et al.\ tested
SpineNetv2 on 491 patients and 2{,}455 discs, reporting overall accuracy of
83.5--97.5\% and superiority over a junior orthopaedic surgeon on canal stenosis,
spondylolisthesis and bilateral foraminal stenosis \cite{wu2026external}. Grob et
al.\ contribute a further independent validation on 882 scans
\cite{grob2022external}.

Studies that build external testing into their development are still uncommon
but instructive. Tumko et al.\ evaluated on 150 studies graded by a panel of
seven radiologists with majority voting and external adjudication
\cite{tumko2024a}; Su et al.\ used a 1{,}273-image external set from a different
hospital and observed accuracy falling by roughly 7--10 points relative to
internal testing \cite{su2022automatic}; Zhang et al.\ found classification
accuracy dropping from 87.70\% to 74.23\% externally \cite{zhang2023deep}; Zeng
et al.\ report an unusually small external drop, from 93.7\% to 92.2\%
\cite{zeng2025gpt4lfs}.

\subsection{Detection Versus Fine-Grained Grading}
\label{sec:lr-detection-vs-grading}

A consistent and under-appreciated finding is that these two capabilities
generalise differently. Zhang et al.\ provide the clearest illustration:
moving to an external institution, detection degraded only modestly---mean IoU
0.82 to 0.70, precision 98.4\% to 96.3\%, sensitivity 99.4\% to 97.8\%---while
classification accuracy fell from 87.70\% to 74.23\% \cite{zhang2023deep}. Wu et
al.\ report the same asymmetry in different terms: binary detection across five
pathologies held up well while ordinal Pfirrmann grading remained the limiting
factor, degrading further in older patients and upper lumbar discs
\cite{wu2026external}. Yilihamu et al.\ show the effect within a single study as
task granularity increases: segmentation IoU remained above 97\% externally,
while classification precision fell from 81.21\% to 74.50\% for an 18-category
problem yet held at 92.51\% to 90.07\% for four-category severity
\cite{yilihamu2025quantification}.

The implication is that where a system fails is predictable. Finding the anatomy
is robust; judging how bad it looks is not. Reporting a single aggregate figure
for a pipeline that performs both obscures this, and evaluation should separate
the two stages.

\subsection{Benchmarking Against Clinicians}
\label{sec:lr-clinician-benchmark}

Comparison against human readers is the most clinically meaningful benchmark, and
the better studies construct it carefully. Won et al.\ provide the cleanest
design: two experts each labelled the data, a model was trained on each expert's
labels, and the key result was that the difference between the two experts was
statistically significant while the difference between each expert and their
corresponding model was not \cite{won2020spinal}. Tumko et al.\ compared against
a seven-radiologist panel average and found the model matched or exceeded it on
all three stenosis types \cite{tumko2024a}. Lee et al.\ benchmarked against eight
participants spanning four experience levels \cite{lee2026deep}, and Wu et al.\
against orthopaedic surgeons of differing seniority \cite{wu2026external}.

Two cautions apply. Comparisons against junior readers or against surgeons rather
than radiologists set a lower bar than subspecialty radiological practice
\cite{ke2025integrating,wu2026external}. And readers in these studies typically
work from cropped images without clinical history, which is not how they read in
practice---so such comparisons measure a narrower task than clinical
interpretation.

\subsection{Longitudinal Stability and Reproducibility}
\label{sec:lr-longitudinal}

One-shot accuracy does not establish that a model is stable over repeat imaging,
which matters for monitoring progression. Murto et al.\ supply the only
longitudinal evidence in this collection, comparing SpineNet's Pfirrmann grading
against radiologists in 19 volunteers imaged at age 37 and again at 51.
Radiologist-to-SpineNet concordance was 0.54 at both timepoints, clearly below
the 0.83 and 0.78 observed between radiologists, and SpineNet systematically
assigned grade~1 to several discs and grade~5 to more discs than the readers did
\cite{murto2025comparison}. The cohort is small, but a systematic offset that is
stable across time is exactly the failure mode that one-shot evaluation cannot
detect. Cheung et al.\ approach the temporal dimension differently, predicting
five-year progression from baseline imaging across 1{,}343 participants
\cite{cheung2022learningbased}.

\subsection{Reporting Quality and Recurring Weaknesses}
\label{sec:lr-reporting}

The systematic reviews converge on a consistent diagnosis. Wang et al.\ pooled
twelve studies covering 15{,}044 patients, obtaining pooled sensitivity 0.84 and
specificity 0.87 with SROC AUC 0.92, but with heterogeneity of $I^2$ near 99\%
and four of twelve studies at high risk of bias; they concluded that no current
model is reliable enough for clinical practice, citing the need for accepted
reference standards, larger datasets, external validation and fuller reporting
\cite{wang2024machine}. Compte et al.\ found no significant difference between
algorithm families and identified the development-to-validation performance drop
as the field's central problem \cite{compte2023are}. Mendes et al.\ found that
most of 56 segmentation studies relied on private datasets and lacked external
validation \cite{mendes2026deep}. Liawrungrueang et al.\ report metrics spanning
71.5\% to 99\% across twenty studies---a range so wide as to indicate that the
underlying tasks are not comparable \cite{liawrungrueang2025artificial}.

Four weaknesses recur: single-institution data without external testing; ground
truth from a single reader with no adjudication; task definitions that vary
between studies while sharing a label vocabulary; and metric selection that
flatters imbalanced problems. The present work adopts the public LumbarDISC
benchmark, its panel-derived labels and its weighted metric specifically to
avoid the first, second and fourth of these.

% ===========================================================================
\section{Interpretability and Clinical Integration}
\label{sec:lr-interpretability}

\subsection{Saliency, Evidence Hotspots, and Grad-CAM}
\label{sec:lr-saliency}

The earliest interpretability contribution in this domain is SpineNet's evidence
hotspots, which map output-layer gradients back to the input to produce a heatmap
per radiological score, obtained without any localisation supervision
\cite{jamaludin2017spinenet}. Grad-CAM \cite{selvaraju2017gradcam} has since
become the default. Tabarestani et al.\ use it to confirm that their cascade
attends to clinically relevant regions \cite{tabarestani2026a}, and Silveira et
al.\ report that attention concentrated on vertebral bodies and nucleus pulposus
in spatial correspondence with manual segmentations \cite{silveira2025automated}.

The value of these methods should not be overstated. A saliency map shows where
a network responded, not why, and confirming that attention falls on plausible
anatomy is a weak check---it excludes gross reliance on background artefact but
does not establish that the model has learned the morphological criteria the
grading system specifies.

\subsection{Inherently Interpretable Models}
\label{sec:lr-interpretable}

A stronger position is to build models whose reasoning is legible by
construction. Bharadwaj et al.\ provide the most striking result in this
literature: alongside a Big Transfer classifier they trained a decision-tree
classifier operating on segmentation-derived features for central canal stenosis,
and the interpretable model \emph{outperformed} the black box, reaching
$\kappa = 0.80$ against 0.54 and approaching inter-reader agreement of 0.74--0.86
\cite{bharadwaj2023deep}. Quantitative morphometry belongs to the same family:
a measured dural sac area \cite{ghobrial2025deep} or a quantitative degeneration
criterion \cite{zheng2022deep} is inspectable in a way a logit is not. Patil et
al.\ make the related point that radiomic features carried the entire
discriminative signal for Pfirrmann grading while deep features added nothing
measurable, and that radiomics additionally yields interpretable shape and
intensity descriptors usable as imaging biomarkers \cite{patil2026mribased}.

Together these findings undercut the assumption that interpretability must be
purchased at the cost of accuracy---at least for targets whose clinical
definition is already explicitly morphological.

\subsection{Vision--Language Models and Report Generation}
\label{sec:lr-vlm}

The most recent direction outputs prose rather than labels. Islam~Sk et al.\
present an explainable vision--language framework that classifies severity,
segments the anatomy and generates radiologist-style reports, built on
biomedical foundation models, with a spatial patch cross-attention module that
preserves the anatomical hierarchies global pooling discards and an adaptive
PID-Tversky loss for class imbalance; they report 90.69\% classification
accuracy, macro-averaged Dice 0.9512 and CIDEr 92.80\%
\cite{islamsk2026an}. Zeng et al.\ take a narrower approach, fusing ConvNeXt
image features with GPT-4o-generated descriptions for foraminal stenosis
classification across three centres and seven scanners \cite{zeng2025gpt4lfs}.
These systems are early and their language metrics measure fluency rather than
clinical correctness, but they indicate where the field is heading.

\subsection{Reader-Assistance Studies and Workflow Impact}
\label{sec:lr-reader-assistance}

The most consequential evidence for clinical value comes not from model metrics
but from measuring what changes when readers use the model. Lim et al.\ conducted
the definitive study of this kind: eight radiologists with 2--13 years of
experience read the same studies with and without deep-learning assistance,
separated by a one-month washout, against a 32-year-veteran musculoskeletal
radiologist as reference. Interpretation time per study fell from 124--274
seconds to 47--71 seconds, and interobserver agreement was superior or equivalent
for every stenosis grading. The largest gain fell to general and in-training
radiologists on four-class foraminal stenosis, where $\kappa$ rose from 0.39 to
0.71 and 0.70 \cite{lim2022improved}. Gao et al.\ observed a comparable effect
for Modic assessment, where AI assistance raised agreement between a junior
radiologist and a senior neuroradiologist from $\kappa = 0.52$ to 0.58
\cite{gao2022automatic}.

The pattern in both is that assistance helps least-experienced readers most and
standardises rather than replaces judgement. This reframes the objective: the
target is not a model that outperforms the best radiologist, but one that raises
the floor.

\subsection{Uncertainty and Calibration}
\label{sec:lr-uncertainty}

Uncertainty estimation is conspicuously underdeveloped in this literature. Where
errors concentrate at grade boundaries and where the reference standard is itself
unreliable, a calibrated confidence would let borderline cases be referred rather
than silently graded---a natural fit for triage. Almost none of the reviewed work
reports calibration. Patil et al.\ are the exception, applying isotonic
calibration on a validation split and reporting Brier scores alongside
discrimination metrics \cite{patil2026mribased}. Since modern networks are
systematically overconfident \cite{guo2017calibration}, an uncalibrated severity
probability should not be interpreted as a clinical confidence, and this remains
an open opportunity rather than a solved problem.

% ===========================================================================
\section{Synthesis and Research Gap}
\label{sec:lr-gap}

\subsection{What the Literature Has Established}
\label{sec:lr-established}

Five findings are supported well enough to be treated as settled.

\textbf{Anatomical localisation precedes classification.} From DeepSPINE
\cite{lu2018deepspine} through the interpretable pipelines of Bharadwaj et al.\
\cite{bharadwaj2023deep} to essentially every leading RSNA challenge solution
\cite{batra2025mscan}, systems that localise the disc level before grading it
outperform those that classify whole images. Detection is close to solved, with
several studies reporting near-perfect level assignment
\cite{pan2021automatically,ghosh2012a}.

\textbf{Automated grading is competitive with human reliability, because human
reliability is low.} Inter-reader agreement is substantial for central canal
stenosis but only moderate for foraminal and poor for subarticular grading
\cite{lurie2008reliability}. Against that benchmark, models sit inside the band
of human disagreement \cite{won2020spinal,mcsweeney2023external,tumko2024a}.
This is a real result, but it is a statement about the modesty of the reference
standard as much as the capability of the models.

\textbf{Errors concentrate at adjacent grade boundaries.} Disagreement above one
grade is rare---0.85\% of discs in the largest external validation
\cite{mcsweeney2023external}---and collapsing to a severe-versus-rest decision
raises agreement substantially \cite{ishimoto2020could}. The difficulty lies in
boundary placement, not gross misclassification.

\textbf{Detection generalises across institutions; fine-grained grading does
not.} The asymmetry is consistent and large
\cite{zhang2023deep,wu2026external,yilihamu2025quantification}.

\textbf{Class imbalance must be addressed explicitly.} With 85.4\% of spinal
canal grades normal or mild \cite{richards2026the}, unweighted objectives yield
degenerate solutions, and weighted sampling, weighted or focal losses and
contrastive pretraining are all established mitigations
\cite{lin2017focal,acharya2026disccentric}.

\subsection{What Remains Contested}
\label{sec:lr-contested}

Three questions remain genuinely open, and the disagreement is substantive
rather than terminological.

\emph{Does architectural sophistication help?} The general computer vision
literature would predict transformers to outperform convolutional networks given
sufficient data, but the lumbar evidence does not bear this out. Kowlagi et al.\
report a tuned CNN pipeline beating transformer approaches \cite{kowlagi2022a};
Lee et al.\ find a CNN superior and a transformer merely comparable
\cite{lee2026deep}; Silveira et al.\ beat TransUNet with an enhanced U-Net
\cite{silveira2025automated}; Udomluck et al.\ find conventional VGG19 features
generalising better than DINOv2 \cite{udomluck2026multitask}; and Patil et al.\
find deep features adding nothing over radiomics \cite{patil2026mribased}. Set
against this, hybrid architectures report gains
\cite{chen2024symtc,zheng2022deep}, and cross-sequence attention reports large
ones \cite{nguyen2026crossspine}.

\emph{Does the ordinal structure of the labels admit exploitation?} Errors
demonstrably cluster at adjacent grades, yet the most systematic test of
ordinal-aware objectives found none improved on plain cross-entropy
\cite{niemeyer2021a}.

\emph{How much imaging is actually required?} Clinical practice and the strongest
benchmark results use all three sequences
\cite{richards2026the,batra2025mscan}, while sagittal-only systems report
competitive performance \cite{acharya2026disccentric,kowlagi2022a} and
sagittal-only systems also fail in a specific, predicted way on the lateral
compartments \cite{nigru2024external}.

\subsection{Gap 1: Architecture Comparison Under Controlled Conditions}
\label{sec:lr-gap-architecture}

The comparative studies above are informative but none isolates the variable of
interest. Three confounds recur.

First, \emph{unequal pipelines}. Compared architectures rarely receive identical
preprocessing, localisation and training budgets, so a reported difference may
reflect engineering effort rather than representational capacity---precisely the
argument Kowlagi et al.\ advance when attributing their CNN's advantage to
pipeline design rather than to convolution as such \cite{kowlagi2022a}. Given the
evidence of Section~\ref{sec:lr-localize-first} that localisation dominates
performance, a comparison in which one arm is better localised than the other
measures preprocessing.

Second, \emph{frozen features}. Several comparisons extract features from frozen
pretrained backbones rather than fine-tuning end to end
\cite{udomluck2026multitask,patil2026mribased}. This systematically
disadvantages transformers, whose weaker inductive biases make them more
dependent on in-domain adaptation.

Third, \emph{fixed input configuration}. Every comparison holds the sequence
configuration constant, so the interaction between architecture and available
context is untested. This matters because the theoretical case for global
attention is strongest exactly where context must be inferred rather than
observed---that is, where a view is missing.

A controlled comparison therefore requires a shared localisation stage, a shared
training and augmentation protocol, comparable capacity and tuning budgets, and
end-to-end fine-tuning of both backbones, evaluated on the same benchmark with
the same metric. No study in this review satisfies all four conditions.

\subsection{Gap 2: Sequence Configuration and the Accuracy--Efficiency Trade-off}
\label{sec:lr-gap-sequence}

The second gap concerns input rather than model. Multi-sequence pipelines carry
costs that are rarely quantified alongside their accuracy: greater memory and
compute, cross-plane geometric registration dependent on DICOM metadata
\cite{batra2025mscan,pan2021automatically}, longer inference, and brittleness
when a sequence is missing or degraded (Section~\ref{sec:lr-missing}). Against
these, the incremental diagnostic yield of each additional sequence has not been
measured under controlled conditions.

The literature supplies both sides of the argument but no direct test. Evidence
that the third sequence matters is indirect, inferred from the failure of
sagittal-only systems on foraminal and subarticular targets
\cite{nigru2024external,wu2026external}; evidence that it may not is likewise
indirect, resting on sagittal-only systems that perform respectably against
different baselines on different data
\cite{acharya2026disccentric,kowlagi2022a}. Because these come from different
architectures, cohorts, localisation strategies and metrics, the contribution of
sequence configuration cannot be separated from everything else that differs.

Two design requirements follow. The comparison must hold architecture and
pipeline constant while varying only the input configuration; and it must report
results per condition rather than in aggregate, since the expected penalty is
concentrated in the foraminal and subarticular targets and an aggregate figure
would average away the structure of interest.

\subsection{Research Questions and Objectives}
\label{sec:lr-questions}

The two gaps combine into a single factorial investigation. Crossing
architecture with input configuration answers both questions and, additionally,
tests whether they interact---whether a transformer's global attention can
partially compensate for absent multiplanar data, which neither gap addresses
alone. This yields the following research questions:

\begin{description}
  \item[RQ1] Under matched localisation, training and evaluation protocols, how
  do convolutional backbones (ResNet50, EfficientNet-B4) compare with a
  hierarchical vision transformer (Swin) for per-level, per-condition severity
  grading of lumbar degenerative disease?

  \item[RQ2] What is the cost in diagnostic performance of reducing the input
  from the full multi-sequence configuration (Sagittal T1, Sagittal T2/STIR,
  Axial T2) to a single-sequence configuration (Sagittal T2), and how is that
  cost distributed across the five graded conditions?

  \item[RQ3] Do architecture and sequence configuration interact---specifically,
  is the penalty for single-sequence input smaller for the transformer than for
  the convolutional backbones?

  \item[RQ4] Which class-imbalance strategies most effectively reduce
  clinically consequential errors, in particular severe cases graded as
  normal or mild?
\end{description}

The corresponding objectives are to implement a shared localisation and
preprocessing stage so that all arms receive identical input regions; to train
each architecture under both configurations with matched budgets; to evaluate
using the benchmark's weighted metric together with per-condition and per-class
reporting; and to quantify computational cost---parameters, memory and inference
time---so that any accuracy difference can be weighed against it.

\subsection{Scope and Delimitations}
\label{sec:lr-scope-delimit}

Four boundaries are stated explicitly. The work addresses \emph{severity
grading}, taking localisation as a shared and largely solved precondition
(Section~\ref{sec:lr-detection}) rather than a contribution. It evaluates on
\emph{LumbarDISC} \cite{richards2026the}; because that collection excludes
studies with hardware, non-degenerative pathology, severe scoliosis or
substantial artifact, conclusions do not extend to those populations, and no
independent external cohort is available to test transportability beyond it. It
compares a \emph{bounded set of architectures}---two convolutional and one
transformer---chosen as representative rather than exhaustive, and does not
evaluate hybrid or vision--language designs
\cite{chen2024symtc,islamsk2026an}. Finally, it reports model performance and
computational cost, not clinical impact: the reader-assistance evidence of
Section~\ref{sec:lr-reader-assistance} \cite{lim2022improved} shows that
workflow benefit requires a prospective reader study, which lies outside the
present scope.

Two constraints inherited from the literature bound the interpretation of any
result. Performance is capped by the reliability of the reference standard
(Section~\ref{sec:lr-label-ceiling}), so figures approaching perfect agreement
should be read as evidence of a protocol artefact rather than of clinical
readiness. And single-benchmark evaluation cannot establish generalisation; the
consistent development-to-validation drop documented in
Section~\ref{sec:lr-generalization} applies to this work as much as to any other.

% ===========================================================================
\section{Chapter Summary}
\label{sec:lr-summary}

This chapter reviewed the literature on automated assessment of lumbar
degenerative disease from MRI, from the clinical grading systems that define the
task through to the architectures that currently attempt it.

The clinical sections established that severity labels originate in
morphological grading systems---Pfirrmann for disc degeneration
\cite{pfirrmann2001magnetic}, Schizas and Lee for central canal
\cite{schizas2010qualitative,lee2011a}, Lee for foraminal
\cite{lee2010a}---whose inter-reader reliability ranges from substantial to poor
depending on the compartment \cite{lurie2008reliability}. This imposes a ceiling
on achievable performance and makes agreement statistics, rather than accuracy,
the appropriate measure. The labels are ordinal, errors concentrate at adjacent
boundaries, and the clinical cost of an error depends on its direction.

The methodological sections traced the field from handcrafted features through
SpineNet's multi-task convolutional grading \cite{jamaludin2017spinenet} to
contemporary multi-stage, multi-view pipelines \cite{batra2025mscan}. The most
robust design principle to emerge is that anatomical localisation should precede
classification \cite{lu2018deepspine,kowlagi2022a}. The most robust empirical
regularity is that detection transfers across institutions while fine-grained
grading does not \cite{zhang2023deep,wu2026external}.

Two questions remain unresolved. Whether attention-based backbones offer any
advantage over well-tuned convolutional ones is contested, with the balance of
lumbar-specific evidence currently favouring CNNs
\cite{kowlagi2022a,lee2026deep,silveira2025automated}, yet no study has compared
them under matched localisation, training and evaluation conditions. And whether
the full three-sequence protocol is necessary, or whether a single ubiquitous
sequence suffices, has been argued from both sides without a controlled test.
These two questions, and their interaction, constitute the gap this thesis
addresses. Chapter~3 sets out the experimental design through which they are
answered.

% ===========================================================================
\printbibliography[heading=bibintoc,title={References}]

\end{document}
